% Dénombrement — Mathématiques 1ère D — Cours signé OnBuch+
% Programme officiel MINESEC. Compiler avec : tectonic N10-denombrement.tex
\newcommand{\DOCMATIERE}{Mathématiques}
\newcommand{\DOCNIVEAU}{1ère D}
\newcommand{\DOCMODULE}{Organisation et gestion des données}
\newcommand{\DOCLECON}{N10}
\newcommand{\DOCTITRE}{Dénombrement}
% OnBuch+ — préambule « cours signés » ENRICHI (compilé avec tectonic / XeLaTeX)
% Système visuel « Pop » d'OnBuch+ : fond crème (jamais blanc), bordures encre
% épaisses, ombres « dures » décalées, titres Archivo Black en capitales.
% Version MATHÉMATIQUES : graphes (pgfplots/tikz), boîtes pédagogiques
% (définition, propriété, méthode, attention, le savais-tu, exercice résolu,
% exercice, TP, bilan), page de garde et signature OnBuch+ ancrée en bas.
%
% Macros attendues dans main.tex AVANT \input{preamble} :
%   \DOCMATIERE  \DOCNIVEAU  \DOCMODULE  \DOCLECON (numéro)  \DOCTITRE
\documentclass[11pt]{article}

\usepackage{fontspec}
\usepackage[french]{babel}
\usepackage[a4paper,margin=2cm,top=3.4cm,bottom=2.8cm,headheight=1.4cm,footskip=1.3cm]{geometry}
\usepackage{amsmath,amssymb,amsfonts}
\usepackage{mathtools} % \xmapsto, \coloneqq... souvent utilisés par les modèles
\usepackage{cancel} % simplifications barrées
% \abs{z} (module, valeur absolue) et \norm{u} (norme), utilisés par les modèles
\providecommand{\abs}{}\let\abs\relax\DeclarePairedDelimiter{\abs}{\lvert}{\rvert}
\providecommand{\norm}{}\let\norm\relax\DeclarePairedDelimiter{\norm}{\lVert}{\rVert}
\usepackage[table]{xcolor} % [table] : fournit aussi \rowcolors (alternance de lignes)
\usepackage{pagecolor}
\usepackage{tikz}
\usetikzlibrary{arrows.meta,positioning,calc,shapes.geometric,shapes.misc,shapes.arrows,decorations.pathmorphing,decorations.pathreplacing,decorations.markings,patterns,shadows,mindmap,trees,fit,backgrounds,matrix,intersections,angles,quotes}
\usepackage{pgfplots}
\usepackage{pgf-pie} % diagrammes circulaires \pie (statistiques)
\pgfplotsset{compat=1.18}
\usepgfplotslibrary{fillbetween,groupplots} % aires entre courbes (calcul intégral)
\usepackage{enumitem}
\usepackage{fancyhdr}
% [most] charge toutes les bibliothèques tcolorbox (listings, minted, raster,
% documentation...) et gonfle inutilement la mémoire TeX sur un document long
% (~300 boîtes) : on ne charge que ce qui est réellement utilisé.
\usepackage[skins,breakable]{tcolorbox}
\usepackage{graphicx}
\usepackage{colortbl}
\usepackage{booktabs}
\usepackage{tabularx}
\usepackage{diagbox} % case d'en-tête diagonale des tableaux à double entrée
% Colonnes extensibles centrée / gauche / droite (C, L, R), souvent utilisées par les modèles
\newcolumntype{C}{>{\centering\arraybackslash}X}
\newcolumntype{L}{>{\raggedright\arraybackslash}X}
\newcolumntype{R}{>{\raggedleft\arraybackslash}X}
\usepackage{array}
\usepackage{multicol}
\usepackage{siunitx}
% parse-numbers=false : accepte aussi \SI{2,5 \times 10^{-3}}{...}
\sisetup{locale=FR,output-decimal-marker={,},group-minimum-digits=4,parse-numbers=false}
\DeclareSIUnit\ohm{\ensuremath{\symup{\Omega}}} % U+2126 absent de la police
\usepackage{qrcode}
% \degree : commande fréquemment utilisée par les modèles pour un angle ou une
% température en dehors de \SI{}{} (non fournie par les paquets chargés).
\providecommand{\degree}{^{\circ}}
% \qed (fin de démonstration), utilisé par les modèles sans amsthm
\providecommand{\qed}{\hfill$\square$}
% \barre : double barre des valeurs interdites dans un tableau de variations (tkz-tab)
\providecommand{\barre}{\ensuremath{\|}}
% environnement proof (amsthm non chargé) utilisé par les modèles
\ifdefined\proof\else\newenvironment{proof}[1][Démonstration]{\par\noindent\textit{#1.}\ }{\qed\par}\fi
\usepackage{lastpage}
\usepackage{hyperref}
\hypersetup{hidelinks}

% ============================================================
% Palette « Pop » OnBuch+ — mêmes hex que la classe P (plus_ui.dart)
% ============================================================
\definecolor{poporange}{HTML}{F59321}
\definecolor{popdark}{HTML}{E07A0C}
\definecolor{popcream}{HTML}{FFF6E8}
\definecolor{popink}{HTML}{1C1714}
\definecolor{poppurple}{HTML}{7A5AE0}
\definecolor{popgreen}{HTML}{2E9E6A}
\definecolor{poppink}{HTML}{E0457B}
\definecolor{popblue}{HTML}{2F7DE1}
\definecolor{popgold}{HTML}{C98A12}
\definecolor{popmuted}{HTML}{8A7F76}
\definecolor{popline}{HTML}{EADFD2}
\definecolor{popgray}{HTML}{8A7F76}
\colorlet{popgrey}{popgray}
% Alias tolérants (le modèle nomme parfois les couleurs différemment)
\colorlet{popwhite}{white}
\colorlet{popblack}{popink}
\colorlet{popred}{poppink}
\colorlet{popyellow}{popgold}
% Teintes claires pour fonds de tableaux / zones
\colorlet{popblueL}{popblue!10!white}
\colorlet{poporangeL}{poporange!14!white}
\colorlet{popgreenL}{popgreen!12!white}
\colorlet{poppurpleL}{poppurple!12!white}
\colorlet{poppinkL}{poppink!12!white}
\colorlet{popgoldL}{popgold!14!white}

\pagecolor{popcream}

% ============================================================
% Typographie
% ============================================================
\defaultfontfeatures{Ligatures=TeX}
\setmainfont{PlusJakartaSans-Medium.ttf}[Path=../../../fonts/,BoldFont=PlusJakartaSans-Bold.ttf]
% Maths en sans-serif (Fira Math) avec chiffres et lettres droites de Plus
% Jakarta, pour que les formules se fondent dans le texte.
\usepackage[math-style=ISO,bold-style=ISO]{unicode-math}
\setmathfont{FiraMath-Regular.otf}[Path=../../../fonts/]
\setmathfont{PlusJakartaSans-Medium.ttf}[Path=../../../fonts/,range={up/num,up/{Latin,latin},\mathup}]
\usepackage{icomma} % virgule décimale sans espace en mode math
% Caractères mathématiques Unicode absents des polices de texte (Plus Jakarta,
% Archivo Black) : rendus par leur équivalent mathématique, en texte comme en
% formule — sinon ils s'affichent en carré vide (ex. « Arithmétique dans ℤ »).
\usepackage{newunicodechar}
\newunicodechar{ℝ}{\ensuremath{\mathbb{R}}}
\newunicodechar{ℤ}{\ensuremath{\mathbb{Z}}}
\newunicodechar{ℂ}{\ensuremath{\mathbb{C}}}
\newunicodechar{ℕ}{\ensuremath{\mathbb{N}}}
\newunicodechar{ℚ}{\ensuremath{\mathbb{Q}}}
\newunicodechar{𝔻}{\ensuremath{\mathbb{D}}}
\newunicodechar{α}{\ensuremath{\alpha}}
\newunicodechar{β}{\ensuremath{\beta}}
\newunicodechar{γ}{\ensuremath{\gamma}}
\newunicodechar{δ}{\ensuremath{\delta}}
\newunicodechar{ε}{\ensuremath{\varepsilon}}
\newunicodechar{θ}{\ensuremath{\theta}}
\newunicodechar{λ}{\ensuremath{\lambda}}
\newunicodechar{μ}{\ensuremath{\mu}}
\newunicodechar{σ}{\ensuremath{\sigma}}
\newunicodechar{τ}{\ensuremath{\tau}}
\newunicodechar{φ}{\ensuremath{\varphi}}
\newunicodechar{ψ}{\ensuremath{\psi}}
\newunicodechar{ω}{\ensuremath{\omega}}
\newunicodechar{Σ}{\ensuremath{\Sigma}}
% majuscules produites par \MakeUppercase dans les titres (α-aminés -> Α)
\newunicodechar{Α}{\ensuremath{\alpha}}
\newunicodechar{Β}{\ensuremath{\beta}}
\newunicodechar{Γ}{\ensuremath{\Gamma}}
\newunicodechar{Δ}{\ensuremath{\Delta}}
\newunicodechar{Ε}{\ensuremath{\varepsilon}}
\newunicodechar{Θ}{\ensuremath{\Theta}}
\newunicodechar{Λ}{\ensuremath{\Lambda}}
\newunicodechar{Μ}{\ensuremath{\mu}}
\newunicodechar{Τ}{\ensuremath{\tau}}
\newunicodechar{Φ}{\ensuremath{\Phi}}
\newunicodechar{Ψ}{\ensuremath{\Psi}}
\newunicodechar{Ω}{\ensuremath{\Omega}}
\newunicodechar{↦}{\ensuremath{\mapsto}}
\newunicodechar{∈}{\ensuremath{\in}}
\newunicodechar{∉}{\ensuremath{\notin}}
\newunicodechar{∧}{\ensuremath{\wedge}}
\newunicodechar{⇔}{\ensuremath{\Leftrightarrow}}
\newunicodechar{⟺}{\ensuremath{\Longleftrightarrow}}
\newunicodechar{⇒}{\ensuremath{\Rightarrow}}
\newunicodechar{⟹}{\ensuremath{\Longrightarrow}}
\newunicodechar{∘}{\ensuremath{\circ}}
\newunicodechar{∪}{\ensuremath{\cup}}
\newunicodechar{∩}{\ensuremath{\cap}}
\newunicodechar{≡}{\ensuremath{\equiv}}
\newunicodechar{⊂}{\ensuremath{\subset}}
\newunicodechar{⊥}{\ensuremath{\perp}}
\newunicodechar{∃}{\ensuremath{\exists}}
\newunicodechar{∀}{\ensuremath{\forall}}
\newunicodechar{∛}{\ensuremath{\sqrt[3]{\,}}}
\newunicodechar{∜}{\ensuremath{\sqrt[4]{\,}}}
\newunicodechar{⃗}{\ensuremath{{}^{\to}}}
\newunicodechar{ⁿ}{\ifmmode^{n}\else\textsuperscript{\ensuremath{n}}\fi}
\newunicodechar{ˣ}{\ifmmode^{x}\else\textsuperscript{\ensuremath{x}}\fi}
\newunicodechar{ᵇ}{\ifmmode^{b}\else\textsuperscript{\ensuremath{b}}\fi}
\newunicodechar{ʳ}{\ifmmode^{r}\else\textsuperscript{\ensuremath{r}}\fi}
\newunicodechar{ᵘ}{\ifmmode^{u}\else\textsuperscript{\ensuremath{u}}\fi}
\newunicodechar{ʸ}{\ifmmode^{y}\else\textsuperscript{\ensuremath{y}}\fi}
\newunicodechar{ᵃ}{\ifmmode^{a}\else\textsuperscript{\ensuremath{a}}\fi}
\newunicodechar{ᵉ}{\ifmmode^{e}\else\textsuperscript{\ensuremath{e}}\fi}
\newunicodechar{ᵏ}{\ifmmode^{k}\else\textsuperscript{\ensuremath{k}}\fi}
\newunicodechar{ᵖ}{\ifmmode^{p}\else\textsuperscript{\ensuremath{p}}\fi}
\newunicodechar{ᴺ}{\ifmmode^{N}\else\textsuperscript{\ensuremath{N}}\fi}
\newunicodechar{ᶜ}{\ifmmode^{c}\else\textsuperscript{\ensuremath{c}}\fi}
\newunicodechar{ʰ}{\ifmmode^{h}\else\textsuperscript{\ensuremath{h}}\fi}
\newunicodechar{ᵠ}{\ifmmode^{q}\else\textsuperscript{\ensuremath{q}}\fi}
\newunicodechar{ₙ}{\ifmmode_{n}\else\textsubscript{\ensuremath{n}}\fi}
\newunicodechar{ₐ}{\ifmmode_{a}\else\textsubscript{\ensuremath{a}}\fi}
\newunicodechar{ᵢ}{\ifmmode_{i}\else\textsubscript{\ensuremath{i}}\fi}
\newunicodechar{ᵣ}{\ifmmode_{r}\else\textsubscript{\ensuremath{r}}\fi}
\newunicodechar{ₖ}{\ifmmode_{k}\else\textsubscript{\ensuremath{k}}\fi}
\newunicodechar{ᵦ}{\ifmmode_{\beta}\else\textsubscript{\ensuremath{\beta}}\fi}
\newunicodechar{ₓ}{\ifmmode_{x}\else\textsubscript{\ensuremath{x}}\fi}
\newfontfamily\popdisplay{ArchivoBlack-Regular.ttf}[Path=../../../fonts/]
\newfontfamily\poplogo{SpaceGrotesk-Bold.ttf}[Path=../../../fonts/]
\newfontfamily\popsemi{PlusJakartaSans-SemiBold.ttf}[Path=../../../fonts/]
\newfontfamily\popxbold{PlusJakartaSans-ExtraBold.ttf}[Path=../../../fonts/]
\newfontfamily\popxboldls{PlusJakartaSans-ExtraBold.ttf}[Path=../../../fonts/,LetterSpace=3.0]

\setlist[enumerate]{leftmargin=1.7em,itemsep=5pt,topsep=4pt}
\setlist[itemize]{leftmargin=1.7em,itemsep=4pt,topsep=4pt,label={\color{poporange}\textbullet}}
\setlength{\parindent}{0pt}
\setlength{\parskip}{5pt}
\renewcommand{\arraystretch}{1.25}

% pgfplots : style Pop par défaut
\pgfplotsset{
  every axis/.append style={axis line style={popink,line width=1pt},tick style={popink},
    grid style={popline,line width=0.5pt},label style={font=\small},tick label style={font=\footnotesize}},
  popcurve/.style={poporange,line width=1.8pt,no markers,smooth},
}
% Même style disponible dans un \draw TikZ simple (hors pgfplots)
\tikzset{popcurve/.style={poporange,line width=1.8pt}}

% ============================================================
% Pill / squiggle
% ============================================================
\newcommand{\poppill}[3]{%
  \tikz[baseline=-0.55ex]\node[draw=popink,line width=1.2pt,rounded corners=2.6pt,%
    fill=#2,inner xsep=6.5pt,inner ysep=3pt]{%
    \popxboldls\fontsize{7}{7}\selectfont\color{#3}\MakeUppercase{#1}};%
}
\newcommand{\popsquiggle}[1]{%
  \tikz[baseline=-0.4ex]{
    \draw[#1,line width=2.6pt,line cap=round]
      plot [smooth] coordinates {(0,0.09) (0.34,0.01) (0.68,0.21) (1.02,0.01) (1.36,0.10)};
  }%
}
% Mot-clé surligné (marqueur orange sous le texte)
\newcommand{\cle}[1]{{\popxbold\color{popdark}#1}}

% ============================================================
% En-tête / pied
% ============================================================
\pagestyle{fancy}
\fancyhf{}
\renewcommand{\headrulewidth}{0pt}
\renewcommand{\footrulewidth}{0pt}
\lhead{\raisebox{-0.15ex}{\poplogo\fontsize{13}{13}\selectfont\color{popink}OnBuch\color{poporange}+}}
\rhead{\poppill{\DOCMATIERE\ \textbullet\ \DOCNIVEAU\ \textbullet\ Leçon \DOCLECON}{white}{popink}}
\lfoot{\popxbold\fontsize{7.6}{7.6}\selectfont\color{popmuted}\MakeUppercase{Cours signé OnBuch+}\ \textbullet\ {\popsemi\DOCTITRE}}
\rfoot{\tikz[baseline=-0.6ex]\node[rounded corners=8.5pt,fill=popink,minimum height=17pt,minimum width=17pt,inner xsep=5pt,inner sep=0]{\popxbold\fontsize{8}{8}\selectfont\color{popcream}\thepage\,/\,\pageref*{LastPage}};}

% ============================================================
% Titres
% ============================================================
\newcommand{\chaptitre}[2]{%
  \begin{tcolorbox}[enhanced,colback=poporange,colframe=popink,boxrule=2.6pt,arc=11pt,
      drop shadow=popink,left=15pt,right=15pt,top=12pt,bottom=11pt,before skip=3pt,after skip=18pt]
    \poppill{#2}{popink}{popcream}\par\vspace{9pt}
    {\raggedright\hyphenpenalty=10000\exhyphenpenalty=10000\popdisplay\fontsize{22}{24}\selectfont\color{popcream}\MakeUppercase{#1}\par}
  \end{tcolorbox}
}
% Section numérotée : pastille orange avec le numéro + titre Archivo
\newcounter{popsec}
\newcommand{\coursec}[1]{%
  \stepcounter{popsec}\setcounter{popsub}{0}%
  \par\vspace{16pt}\noindent
  \begin{minipage}{\linewidth}
  \tikz[baseline=-0.7ex]\node[draw=popink,line width=1.6pt,fill=poporange,rounded corners=4pt,minimum size=22pt,inner sep=0]{\popdisplay\fontsize{12}{12}\selectfont\color{popcream}\thepopsec};%
  \hspace{8pt}{\popdisplay\fontsize{15}{17}\selectfont\color{popink}\MakeUppercase{#1}}\par
  \vspace{4pt}\noindent{\color{poporange}\rule{36pt}{3.6pt}}
  \end{minipage}\par\nopagebreak\vspace{8pt}
}
\newcounter{popsub}
\newcommand{\courssub}[1]{%
  \stepcounter{popsub}%
  \par\vspace{9pt}\noindent{\popxbold\fontsize{11.5}{13}\selectfont\color{popink}{\color{poporange}\thepopsec.\thepopsub}\hspace{6pt}#1}\par\nopagebreak\vspace{4pt}
}

% ============================================================
% Boîtes pédagogiques — même recette : fond blanc, bordure épaisse colorée,
% ombre dure encre, badge plein. Titre optionnel [#1].
% ============================================================
\tcbset{popbase/.style={breakable,enhanced,colback=white,boxrule=2.3pt,arc=9pt,
  shadow={3pt}{-3pt}{0pt}{popink},left=14pt,right=14pt,top=11pt,bottom=11pt,before skip=11pt,after skip=15pt,
  fontupper=\popsemi\fontsize{10.6}{14.5}\selectfont\color{popink},
  fontlower=\popsemi\fontsize{10.6}{14.5}\selectfont\color{popink}}}
% Coche dessinée (le glyphe ✓ n'existe pas dans les polices du document)
\newcommand{\popcheck}[1]{\tikz[baseline=-0.5ex]\draw[#1,line width=1.6pt,line cap=round,line join=round](-0.13,0.0)--(-0.04,-0.09)--(0.13,0.1);}
\newcommand{\popboxhead}[3]{\poppill{#1}{#2}{white}\ifx\relax#3\relax\else\hspace{6pt}{\popxbold\fontsize{10.6}{12}\selectfont\color{popink}#3}\fi\par\vspace{7pt}}

\newtcolorbox{aretenir}[1][]{popbase,colframe=popgreen,before upper={\popboxhead{À retenir}{popgreen}{#1}}}
\newtcolorbox{exemplebox}[1][]{popbase,colframe=poporange,before upper={\popboxhead{Exemple}{poporange}{#1}}}
\newtcolorbox{definition}[1][]{popbase,colframe=popblue,before upper={\popboxhead{Définition}{popblue}{#1}}}
\newtcolorbox{propriete}[1][]{popbase,colframe=poppurple,before upper={\popboxhead{Propriété}{poppurple}{#1}}}
\newtcolorbox{methode}[1][]{popbase,colframe=popgold,before upper={\popboxhead{Méthode}{popgold}{#1}}}
\newtcolorbox{attention}[1][]{popbase,colframe=poppink,before upper={\popboxhead{Attention}{poppink}{#1}}}
\newtcolorbox{experience}[1][]{popbase,colframe=popblue,colback=popblueL,before upper={\popboxhead{Expérience}{popblue}{#1}}}
% « Le savais-tu ? » : carte violette pleine (contexte, culture, Cameroun)
\newtcolorbox{savaistu}[1][]{popbase,colback=poppurple,colframe=popink,
  fontupper=\popsemi\fontsize{10.4}{14.2}\selectfont\color{white},
  before upper={\poppill{Le savais-tu\,?}{white}{poppurple}\ifx\relax#1\relax\else\hspace{6pt}{\popxbold\color{white}#1}\fi\par\vspace{7pt}}}
% Exercice résolu : énoncé en haut, \tcblower puis la solution sur fond crème
\newcounter{popexo}\newcounter{popexr}
\newtcolorbox{exoresolu}[1][]{popbase,colframe=poporange,colbacklower=popcream,
  segmentation style={popink,line width=1.2pt,dashed},
  before upper={\stepcounter{popexr}\popboxhead{Exercice résolu \thepopexr}{poporange}{#1}},
  before lower={\poppill{Solution}{popink}{popcream}\par\vspace{6pt}}}
% Exercice (non résolu en place) — la correction est dans la partie Corrigés
\newtcolorbox{exercice}[1][]{popbase,colframe=popink,
  before upper={\stepcounter{popexo}\popboxhead{Exercice \thepopexo}{popink}{#1}}}
% Corrigé
\newtcolorbox{corrige}[1][]{popbase,colframe=popgreen,colback=popgreenL,
  before upper={\popboxhead{Corrigé}{popgreen}{#1}}}
% Objectifs / prérequis / bilan
\newtcolorbox{objectifs}{popbase,colframe=popink,colback=poporangeL,
  before upper={\popboxhead{Objectifs de la leçon}{popink}{}}}
\newtcolorbox{prerequis}{popbase,colframe=popink,colback=popblueL,
  before upper={\popboxhead{Prérequis}{popblue}{}}}
\newtcolorbox{bilan}{popbase,colframe=popink,colback=white,boxrule=2.8pt,
  before upper={\popboxhead{Fiche bilan}{poporange}{L'essentiel à connaître pour l'examen}}}
% Figure encadrée avec légende
\newtcolorbox{popfigure}[1][]{enhanced,breakable=false,colback=white,colframe=popline,boxrule=1.4pt,arc=8pt,
  left=10pt,right=10pt,top=10pt,bottom=8pt,before skip=10pt,after skip=12pt,halign=center,
  lower separated=false,halign lower=center,
  fontlower=\popsemi\fontsize{9}{11.5}\selectfont\color{popmuted},
  #1}
\newcommand{\legende}[1]{\par\vspace{6pt}{\popsemi\fontsize{9}{11.5}\selectfont\color{popmuted}#1}}

% ============================================================
% Signature OnBuch+
% ============================================================
\newcommand{\poplogoinline}[1]{{\poplogo\fontsize{#1}{#1}\selectfont\color{popink}OnBuch\color{poporange}+}}
\newcommand{\signatureonbuch}{%
  \begin{tcolorbox}[enhanced,colback=popink,colframe=popink,boxrule=2.2pt,arc=10pt,
      drop shadow=poporange,left=14pt,right=14pt,top=10pt,bottom=10pt,before skip=0pt,after skip=0pt]
    \tikz[baseline=-0.7ex]\node[circle,fill=popgreen,draw=popcream,line width=1.3pt,minimum size=22pt,inner sep=0]{\popcheck{white}};%
    \hspace{9pt}\begin{minipage}[c]{0.62\linewidth}
      {\popxbold\fontsize{12}{13}\selectfont\color{popcream}Signé\ \textbullet\ {\poplogo OnBuch\color{poporange}+}}\par\vspace{2pt}
      {\raggedright\popsemi\fontsize{8}{10.5}\selectfont\color{popline}\MakeUppercase{Équipe pédagogique OnBuch+}\\\MakeUppercase{Conforme au programme officiel MINESEC}\par}
    \end{minipage}\hfill
    {\poplogo\fontsize{22}{22}\selectfont\color{popcream}OnBuch\color{poporange}+}
  \end{tcolorbox}%
}

% ============================================================
% Page de garde
% #1 titre · #2 matière · #3 niveau · #4 module · #5 n° leçon · #6 durée ·
% #7 description / objectifs (texte ou itemize)
% ============================================================
\newcommand{\pagegarde}[7]{%
  \thispagestyle{empty}
  \newgeometry{a4paper,margin=1.7cm,top=1.6cm,bottom=1.4cm}%
  \begin{tikzpicture}[remember picture,overlay]
    \fill[poporange!18] ([xshift=-3.2cm,yshift=-3.6cm]current page.north east) circle (4.6cm);
    \fill[poppurple!14] ([xshift=2.4cm,yshift=7.6cm]current page.south west) circle (3.8cm);
  \end{tikzpicture}%
  {\poplogo\fontsize{30}{30}\selectfont\color{popink}OnBuch\color{poporange}+}\hfill
  \raisebox{0.6ex}{\poppill{Cours signé \textbullet\ Premium}{popink}{popcream}}\par
  \vspace{1pt}\popsquiggle{poppurple}\par
  \vspace{8pt}
  \begin{tcolorbox}[enhanced,colback=poporange,colframe=popink,boxrule=2.8pt,arc=13pt,
      drop shadow=popink,left=18pt,right=18pt,top=12pt,bottom=13pt,after skip=10pt]
    \poppill{#2 \textbullet\ #3}{popink}{popcream}\hspace{5pt}\poppill{#4}{white}{popink}\par\vspace{8pt}
    {\popdisplay\fontsize{14}{14}\selectfont\color{popink}LEÇON #5}\par\vspace{4pt}
    {\raggedright\hyphenpenalty=10000\exhyphenpenalty=10000\popdisplay\fontsize{25}{27}\selectfont\color{popcream}\MakeUppercase{#1}\par}
    \vspace{8pt}
    {\popxbold\fontsize{9.5}{9.5}\selectfont\color{popink}\MakeUppercase{Durée conseillée : #6}}
  \end{tcolorbox}
  \begin{tcolorbox}[enhanced,colback=white,colframe=popink,boxrule=2.2pt,arc=10pt,
      drop shadow=popink,left=16pt,right=16pt,top=10pt,bottom=10pt,after skip=8pt]
    \poppill{Dans cette leçon}{poppurple}{white}\par\vspace{5pt}
    \popsemi\fontsize{9.3}{12.6}\selectfont\color{popink}#7
  \end{tcolorbox}
  \noindent\begin{minipage}[t]{0.487\linewidth}
    \begin{tcolorbox}[enhanced,colback=popgreen,colframe=popink,boxrule=2.2pt,arc=10pt,
        drop shadow=popink,left=11pt,right=11pt,top=10pt,bottom=10pt,height=3.1cm,valign=center]
      \tcbox[colback=white,colframe=popink,boxrule=1.3pt,arc=5pt,boxsep=0pt,nobeforeafter,
        left=3pt,right=3pt,top=3pt,bottom=3pt,valign=center]{\qrcode[height=1.9cm]{https://play.google.com/store/apps/details?id=com.luvvix.onbuch&hl=fr}}%
      \hspace{8pt}\begin{minipage}[c]{0.5\linewidth}
        {\popdisplay\fontsize{11}{12.5}\selectfont\color{white}\MakeUppercase{Télécharge OnBuch}}\par\vspace{4pt}
        {\raggedright\popsemi\fontsize{8.4}{11}\selectfont\color{white}Gratuit \textbullet\ Google Play\\Tuteur IA, annales, concours, résultats\par}
      \end{minipage}
    \end{tcolorbox}
  \end{minipage}\hfill
  \begin{minipage}[t]{0.487\linewidth}
    \begin{tcolorbox}[enhanced,colback=poppurple,colframe=popink,boxrule=2.2pt,arc=10pt,
        drop shadow=popink,left=11pt,right=11pt,top=10pt,bottom=10pt,height=3.1cm,valign=center]
      \tikz[baseline=-0.7ex]\node[circle,fill=white,draw=popink,line width=1.3pt,minimum size=1.6cm,inner sep=0]{\popdisplay\fontsize{24}{24}\selectfont\color{poppurple}+};%
      \hspace{8pt}\begin{minipage}[c]{0.6\linewidth}
        {\popdisplay\fontsize{11}{12.5}\selectfont\color{white}\MakeUppercase{Passe à OnBuch+}}\par\vspace{4pt}
        {\raggedright\popsemi\fontsize{8.4}{11}\selectfont\color{white}Tous les cours signés, Tuteur IA illimité, fiches PDF — onglet OnBuch+ de l'app\par}
      \end{minipage}
    \end{tcolorbox}
  \end{minipage}
  \vspace{8pt}
  \popcontenu
  \vfill
  \signatureonbuch
  \restoregeometry
  \newpage
}

% Rangée « Ce cours contient » (page de garde)
\newcommand{\poptile}[3]{% #1 couleur #2 gros texte #3 légende
  \begin{tcolorbox}[enhanced,colback=#1,colframe=popink,boxrule=2pt,arc=9pt,shadow={2.5pt}{-2.5pt}{0pt}{popink},
      left=6pt,right=6pt,top=9pt,bottom=9pt,halign=center,valign=center,height=1.75cm,width=\linewidth,nobeforeafter]
    {\popdisplay\fontsize{12.5}{13}\selectfont\color{white}\MakeUppercase{#2}}\par\vspace{3pt}
    {\centering\popsemi\fontsize{7.8}{9.5}\selectfont\color{white}#3\par}
  \end{tcolorbox}}
\newcommand{\popcontenu}{%
  \par\noindent{\popxboldls\fontsize{7.5}{7.5}\selectfont\color{popmuted}CE COURS SIGNÉ CONTIENT}\par\vspace{7pt}
  \noindent\begin{minipage}[t]{0.235\linewidth}\poptile{popblue}{Cours}{complet, illustré, pas à pas}\end{minipage}\hfill
  \begin{minipage}[t]{0.235\linewidth}\poptile{poporange}{Méthodes}{et exercices résolus}\end{minipage}\hfill
  \begin{minipage}[t]{0.235\linewidth}\poptile{poppink}{Exercices}{type examen + corrigés}\end{minipage}\hfill
  \begin{minipage}[t]{0.235\linewidth}\poptile{popgreen}{Fiche bilan}{l'essentiel à retenir}\end{minipage}\par}

% Sommaire
\newtcolorbox{sommaire}{enhanced,colback=white,colframe=popink,boxrule=2.2pt,arc=10pt,
  drop shadow=popink,left=18pt,right=18pt,top=13pt,bottom=13pt,before skip=6pt,after skip=20pt,
  before upper={\poppill{Sommaire}{poppurple}{white}\par\vspace{9pt}\popsemi\fontsize{10.6}{14.5}\selectfont\color{popink}}}

% Signature de fin de cours — poussée tout en bas de la dernière page
\newcommand{\findecours}{%
  \par\vfill\nopagebreak
  \begin{center}\popsquiggle{poporange}\end{center}\vspace{4pt}
  \signatureonbuch
}
\AtBeginDocument{\def\llbracket{\mathopen{[\mkern-3.5mu[}}\def\rrbracket{\mathclose{]\mkern-3.5mu]}}} % intervalles d'entiers (glyphe ⟦ absent de la police)
% Glyphes absents de Fira Math : redéfinis à partir de glyphes disponibles
\AtBeginDocument{%
  \def\setminus{\mathbin{\backslash}}\let\smallsetminus\setminus
  \def\vdots{\mathord{\vbox{\baselineskip3.2pt\lineskiplimit0pt\kern4pt\hbox{.}\hbox{.}\hbox{.}}}}%
  \def\ddots{\mathinner{\mkern1mu\raise6.5pt\hbox{.}\mkern2mu\raise3.6pt\hbox{.}\mkern2mu\raise0.7pt\hbox{.}\mkern1mu}}%
  \def\triangle{\mathord{\scalebox{0.9}{$\symup{\Delta}$}}}%
  \def\nabla{\mathord{\raisebox{\depth}{\scalebox{1}[-1]{$\symup{\Delta}$}}}}%
  \def\checkmark{\popcheck{popdark}}%
  \def\star{\mathbin{\ast}}\def\bigstar{\mathord{\ast}}%
  \def\diamond{\mathbin{\scalebox{0.6}{$\Diamond$}}}%
  \def\bigcup{\mathop{\mathchoice{\raisebox{-0.3em}{\scalebox{1.7}{$\cup$}}}{\scalebox{1.25}{$\cup$}}{\cup}{\cup}}\displaylimits}%
  \def\bigcap{\mathop{\mathchoice{\raisebox{-0.3em}{\scalebox{1.7}{$\cap$}}}{\scalebox{1.25}{$\cap$}}{\cap}{\cap}}\displaylimits}%
  \def\lesseqqgtr{\mathrel{\lessgtr}}\def\bigoplus{\mathop{\oplus}\displaylimits}%
}
\providecommand{\ssi}{si et seulement si} % abréviation fréquente
\AtBeginDocument{\providecommand{\wideparen}{\overparen}} % arc de cercle

%% FIGFIX-BEGIN v5 — correctifs globaux des figures (docs/onbuch-generation/tools/figfix)
\usepackage{adjustbox}\usepackage{etoolbox}\tcbuselibrary{raster}
% toute figure TikZ/pgfplots plus large que la ligne est réduite à la largeur disponible
\BeforeBeginEnvironment{tikzpicture}{\begin{adjustbox}{max width=\linewidth}}
\AfterEndEnvironment{tikzpicture}{\end{adjustbox}}
% légendes pgfplots lisibles par-dessus les courbes
\pgfplotsset{every axis/.append style={legend style={fill=white,fill opacity=0.92,text opacity=1,draw=black!15,inner sep=3pt}}}
% étiquettes d'arêtes (syntaxe edge["..."]) avec fond blanc
\tikzset{every edge quotes/.append style={fill=white,inner sep=1.2pt}}
%% FIGFIX-END
\begin{document}

\pagegarde{Dénombrement}{Mathématiques}{1ère D}{Organisation et gestion des données}{N10}{6 h}{Cette leçon fournit les outils fondamentaux pour compter les éléments d'ensembles finis : diagrammes de Venn, tableaux, arbres, p-uplets, arrangements, combinaisons et formule du binôme de Newton. Elle prépare directement le calcul des probabilités de Terminale et constitue un chapitre incontournable du Probatoire.
\vspace{4pt}
\begin{multicols}{2}\small
\begin{itemize}
\item À la fin de cette leçon, je sais utiliser un diagramme de Venn ou un tableau à double entrée pour dénombrer les éléments de réunions et d'intersections d'ensembles finis
\item À la fin de cette leçon, je sais construire un arbre de choix pour dénombrer des situations successives
\item À la fin de cette leçon, je sais déterminer le nombre de p-uplets avec répétition (n\^{}p) et sans répétition d'éléments pris parmi n
\item À la fin de cette leçon, je sais calculer un arrangement A(n,p) et une factorielle n!
\item À la fin de cette leçon, je sais calculer une combinaison C(n,p) et utiliser la relation A(n,p) = C(n,p) × p!
\item À la fin de cette leçon, je sais déterminer le nombre de parties d'un ensemble à n éléments (2\^{}n)
\end{itemize}\end{multicols}}

\begin{sommaire}
\begin{multicols}{2}
\begin{enumerate}
\item Cardinal d'un ensemble fini et premiers outils
\item Arbres de choix et p-uplets
\item Factorielle, permutations et arrangements
\item Combinaisons et parties d'un ensemble
\item Triangle de Pascal et formule du binôme
\item Synthèse : choisir le bon outil de dénombrement
\item Activité d'intégration
\item Méthodes et exercices résolus
\item Exercices
\item Corrigés des exercices
\item Fiche bilan
\end{enumerate}
\end{multicols}
\end{sommaire}

\begin{objectifs}
\begin{itemize}
\item À la fin de cette leçon, je sais utiliser un diagramme de Venn ou un tableau à double entrée pour dénombrer les éléments de réunions et d'intersections d'ensembles finis
\item À la fin de cette leçon, je sais construire un arbre de choix pour dénombrer des situations successives
\item À la fin de cette leçon, je sais déterminer le nombre de p-uplets avec répétition (n\^{}p) et sans répétition d'éléments pris parmi n
\item À la fin de cette leçon, je sais calculer un arrangement A(n,p) et une factorielle n!
\item À la fin de cette leçon, je sais calculer une combinaison C(n,p) et utiliser la relation A(n,p) = C(n,p) × p!
\item À la fin de cette leçon, je sais déterminer le nombre de parties d'un ensemble à n éléments (2\^{}n)
\item À la fin de cette leçon, je sais construire le triangle de Pascal et développer (a+b)\^{}n
\item À la fin de cette leçon, je sais choisir l'outil de dénombrement adapté (liste, arrangement, combinaison) selon que l'ordre compte ou non
\end{itemize}
\end{objectifs}

\begin{prerequis}
\begin{itemize}
\item Notions sur les ensembles : appartenance, inclusion, réunion, intersection, complémentaire (Seconde)
\item Cardinal d'un ensemble fini et opérations sur les entiers naturels
\item Calcul littéral : puissances, développement de (a+b)² et (a+b)³
\item Lecture et construction de tableaux à double entrée (statistiques de Seconde)
\end{itemize}
\end{prerequis}

\newpage

\coursec{Cardinal d'un ensemble fini et premiers outils}

Au lycée de Ngaoundéré, c'est la grande effervescence. Le directeur organise le tournoi de football inter-classes : $8$ classes participent, et il veut savoir combien de matchs programmer si chaque classe rencontre chaque autre une seule fois. Pendant ce temps, à la cantine scolaire, la cuisinière propose chaque jour un menu composé d'un plat, d'un accompagnement et d'une boisson : avec $4$ plats, $3$ accompagnements et $2$ boissons, combien de menus différents peut-elle composer ? Un peu plus loin, un élève de Première C fixe l'écran de son téléphone : combien existe-t-il de codes PIN à $4$ chiffres ?

Ces trois questions semblent très différentes. Pourtant, elles ont un point commun fondamental : il s'agit de \cle{compter sans énumérer}. Personne ne songerait à écrire les dix mille codes PIN un par un ! Cet art de compter intelligemment s'appelle le \cle{dénombrement}. Dans cette première section, nous posons les fondations : la notion de cardinal, les formules de la réunion et du complémentaire, et deux outils visuels redoutables : les diagrammes de Venn et les tableaux à double entrée. Nous terminerons par le principe multiplicatif, clé de voûte de toute la leçon.

\courssub{Ensemble fini et cardinal}

\begin{definition}[Ensemble fini, cardinal]
Un ensemble $E$ est dit \cle{fini} s'il possède un nombre fini d'éléments. Le \cle{cardinal} de $E$, noté $\operatorname{card}(E)$ ou $|E|$, est le nombre d'éléments de $E$.

Par convention, l'ensemble vide vérifie $\operatorname{card}(\varnothing) = 0$.
\end{definition}

\begin{exemplebox}[Premiers cardinaux]
\begin{itemize}
\item Si $E = \{2 ; 4 ; 6 ; 8\}$, alors $\operatorname{card}(E) = 4$.
\item L'ensemble des jours de la semaine a pour cardinal $7$.
\item L'ensemble des élèves d'une classe de Première D de Garoua a pour cardinal $60$ (par exemple).
\item L'ensemble $\mathbb{N}$ des entiers naturels n'est \emph{pas} fini : on ne peut pas lui attribuer de cardinal au sens de cette leçon.
\end{itemize}
\end{exemplebox}

\begin{attention}[Un ensemble n'a pas de répétition]
Dans un ensemble, chaque élément n'est compté \textbf{qu'une seule fois}. Si $E = \{1 ; 2 ; 2 ; 3\}$, alors $E = \{1 ; 2 ; 3\}$ et $\operatorname{card}(E) = 3$, et non $4$. De même, dans un problème concret, un même élève ne doit jamais être compté deux fois : c'est toute la difficulté du dénombrement !
\end{attention}

\courssub{Cardinal de la réunion : la formule fondamentale}

Imaginons une classe où certains élèves étudient l'allemand (ensemble $A$) et d'autres l'espagnol (ensemble $B$). Si l'on additionne naïvement $\operatorname{card}(A) + \operatorname{card}(B)$, les élèves qui étudient \emph{les deux} langues sont comptés deux fois. Il faut donc retrancher une fois l'intersection.

\begin{propriete}[Cardinal de la réunion de deux ensembles]
Soient $A$ et $B$ deux parties d'un ensemble fini $E$. Alors :
\[ \operatorname{card}(A \cup B) = \operatorname{card}(A) + \operatorname{card}(B) - \operatorname{card}(A \cap B). \]
En particulier, si $A$ et $B$ sont \cle{disjoints} (c'est-à-dire $A \cap B = \varnothing$), alors :
\[ \operatorname{card}(A \cup B) = \operatorname{card}(A) + \operatorname{card}(B). \]
\end{propriete}

\begin{experience}[Justification guidée par le diagramme de Venn]
Découpons la réunion $A \cup B$ en trois régions \emph{sans chevauchement} :
\begin{enumerate}
\item la région $R_1$ : les éléments qui sont dans $A$ mais pas dans $B$ (c'est $A \setminus B$) ;
\item la région $R_2$ : les éléments qui sont dans $A \cap B$ ;
\item la région $R_3$ : les éléments qui sont dans $B$ mais pas dans $A$.
\end{enumerate}
Ces trois régions sont deux à deux disjointes et leur réunion est $A \cup B$, donc :
\[ \operatorname{card}(A \cup B) = \operatorname{card}(R_1) + \operatorname{card}(R_2) + \operatorname{card}(R_3). \]
Or $\operatorname{card}(A) = \operatorname{card}(R_1) + \operatorname{card}(R_2)$ et $\operatorname{card}(B) = \operatorname{card}(R_2) + \operatorname{card}(R_3)$. En additionnant :
\[ \operatorname{card}(A) + \operatorname{card}(B) = \operatorname{card}(R_1) + 2\,\operatorname{card}(R_2) + \operatorname{card}(R_3). \]
La différence avec $\operatorname{card}(A \cup B)$ est exactement $\operatorname{card}(R_2) = \operatorname{card}(A \cap B)$, compté une fois de trop. D'où la formule. $\blacksquare$
\end{experience}

\begin{popfigure}
\begin{center}
\begin{tikzpicture}[scale=1.1]
% Fond E
\draw[rounded corners, fill=popcream] (-3.4,-2.2) rectangle (3.4,2.4);
\node[popdark] at (3.0,2.1) {$E$};
% Ensemble A
\draw[thick, popblue, fill=popblueL, fill opacity=0.6] (-0.9,0) circle (1.6);
% Ensemble B
\draw[thick, poporange, fill=poporangeL, fill opacity=0.6] (0.9,0) circle (1.6);
% Réintersection pour lisibilité
\begin{scope}
\clip (-0.9,0) circle (1.6);
\fill[poppurple, fill opacity=0.45] (0.9,0) circle (1.6);
\end{scope}
\draw[thick, popblue] (-0.9,0) circle (1.6);
\draw[thick, poporange] (0.9,0) circle (1.6);
\node[popblue, font=\bfseries] at (-2.0,1.15) {$A$};
\node[poporange, font=\bfseries] at (2.0,1.15) {$B$};
\node at (-1.35,0) {$23$};
\node at (0,0) {$12$};
\node at (1.35,0) {$16$};
\node[popmuted] at (-2.7,-1.6) {$9$};
\node[popmuted, font=\small] at (-2.7,-1.95) {aucune};
\end{tikzpicture}
\end{center}
\legende{Diagramme de Venn : $A$ (allemand) et $B$ (espagnol). L'intersection, en violet, contient les $12$ élèves qui étudient les deux langues.}
\end{popfigure}

\begin{exoresolu}[Langues vivantes en Première]
\textbf{Énoncé.} Dans une classe de $60$ élèves, $35$ étudient l'allemand, $28$ étudient l'espagnol et $12$ étudient les deux langues. Combien d'élèves n'étudient aucune de ces deux langues ?
\tcblower
\textbf{Solution détaillée.} Notons $E$ l'ensemble des élèves de la classe, $A$ l'ensemble des élèves étudiant l'allemand et $B$ celui des élèves étudiant l'espagnol. On a :
\[ \operatorname{card}(E) = 60, \quad \operatorname{card}(A) = 35, \quad \operatorname{card}(B) = 28, \quad \operatorname{card}(A \cap B) = 12. \]
\textbf{Étape 1 :} nombre d'élèves étudiant \emph{au moins une} des deux langues :
\[ \operatorname{card}(A \cup B) = 35 + 28 - 12 = 51. \]
\textbf{Étape 2 :} les élèves n'étudiant aucune des deux langues forment le complémentaire de $A \cup B$ dans $E$ :
\[ \operatorname{card}\bigl(\overline{A \cup B}\bigr) = 60 - 51 = 9. \]
\textbf{Conclusion :} $9$ élèves n'étudient ni l'allemand ni l'espagnol.

\emph{Vérification par régions :} allemand seul : $35 - 12 = 23$ ; espagnol seul : $28 - 12 = 16$ ; les deux : $12$ ; aucune : $9$. Total : $23 + 16 + 12 + 9 = 60$. La classe est bien complète.
\end{exoresolu}

\begin{attention}[L'erreur classique du double comptage]
L'erreur la plus fréquente au Probatoire est d'écrire $\operatorname{card}(A \cup B) = \operatorname{card}(A) + \operatorname{card}(B)$ \emph{sans vérifier} si les ensembles sont disjoints. Dans l'exemple précédent, cela donnerait $35 + 28 = 63 > 60$, ce qui est absurde puisque la classe ne compte que $60$ élèves ! Dès que l'intersection est non vide, \textbf{il faut retrancher} $\operatorname{card}(A \cap B)$. Réflexe à acquérir : avant d'additionner, toujours se demander « y a-t-il des éléments communs ? ».
\end{attention}

\courssub{Cardinal du complémentaire}

\begin{propriete}[Cardinal du complémentaire]
Soit $A$ une partie d'un ensemble fini $E$. Le complémentaire de $A$ dans $E$, noté $\bar{A}$ ou $E \setminus A$, vérifie :
\[ \operatorname{card}(\bar{A}) = \operatorname{card}(E) - \operatorname{card}(A). \]
\end{propriete}

Cette propriété découle immédiatement du fait que $A$ et $\bar{A}$ sont disjoints et que $A \cup \bar{A} = E$. Elle est précieuse dans une stratégie très puissante : pour compter les éléments vérifiant une propriété, il est parfois plus simple de compter ceux qui \emph{ne la vérifient pas} et de soustraire. C'est le \cle{raisonnement par le complémentaire}, que nous avons déjà utilisé à l'étape 2 de l'exercice résolu. Nous le retrouverons constamment : « au moins un » se traduit souvent par « total moins aucun ».

\courssub{Diagrammes de Venn : deux puis trois ensembles}

Le \cle{diagramme de Venn} représente les ensembles par des disques (ou des ovales) tracés dans un rectangle qui figure l'ensemble de référence $E$. Chaque région du dessin correspond à une situation précise, et l'on y inscrit les \emph{effectifs}.

\begin{methode}[Résoudre un problème avec un diagramme de Venn]
\begin{enumerate}
\item Nommer les ensembles et traduire les données de l'énoncé.
\item Dessiner le rectangle $E$ et les disques (deux ou trois selon le nombre de critères).
\item \textbf{Commencer toujours par la région la plus centrale} : l'intersection de tous les ensembles.
\item Remplir les régions voisines \emph{par soustractions successives} (par exemple : « $A$ seul » $= \operatorname{card}(A) - \operatorname{card}(A \cap B)$).
\item Terminer par la région extérieure (« aucun ») en soustrayant du total.
\item Vérifier que la somme de toutes les régions redonne $\operatorname{card}(E)$.
\end{enumerate}
\end{methode}

Avec trois ensembles, le diagramme possède $2^3 = 8$ régions. Prenons l'exemple des sports pratiqués par les $50$ élèves d'une classe de Première TI : football ($F$), handball ($H$), athlétisme ($A$).

\begin{exemplebox}[Trois sports, huit régions]
Dans une classe de $50$ élèves : $30$ jouent au football, $20$ au handball, $15$ font de l'athlétisme ; $10$ pratiquent football et handball, $6$ football et athlétisme, $5$ handball et athlétisme, et $3$ pratiquent les trois sports. Remplissage progressif :
\begin{itemize}
\item les trois sports : $3$ ;
\item $F \cap H$ seulement : $10 - 3 = 7$ ; \quad $F \cap A$ seulement : $6 - 3 = 3$ ; \quad $H \cap A$ seulement : $5 - 3 = 2$ ;
\item $F$ seul : $30 - 7 - 3 - 3 = 17$ ; \quad $H$ seul : $20 - 7 - 2 - 3 = 8$ ; \quad $A$ seul : $15 - 3 - 2 - 3 = 7$ ;
\item aucun sport : $50 - (17 + 8 + 7 + 7 + 3 + 2 + 3) = 50 - 47 = 3$.
\end{itemize}
Vérification : $17 + 8 + 7 + 7 + 3 + 2 + 3 + 3 = 50$. Le compte est bon.
\end{exemplebox}

\begin{popfigure}
\begin{center}
\begin{tikzpicture}[scale=0.95]
\draw[rounded corners, fill=popcream] (-3.6,-2.9) rectangle (3.6,3.0);
\node[popdark] at (3.2,2.7) {$E$};
\draw[thick, popblue, fill=popblueL, fill opacity=0.45] (-0.8,0.55) circle (1.5);
\draw[thick, poporange, fill=poporangeL, fill opacity=0.45] (0.8,0.55) circle (1.5);
\draw[thick, popgreen, fill=popgreenL, fill opacity=0.45] (0,-0.75) circle (1.5);
\node[popblue, font=\bfseries] at (-1.9,1.75) {$F$};
\node[poporange, font=\bfseries] at (1.9,1.75) {$H$};
\node[popgreen!60!black, font=\bfseries] at (0,-2.55) {$A$};
\node at (0,0.85) {\small $3$};
\node at (-0.85,0.35) {\small $7$};
\node at (0.85,0.35) {\small $2$};
\node at (0,-0.55) {\small $3$};
\node at (-1.35,1.15) {\small $17$};
\node at (1.35,1.15) {\small $8$};
\node at (0,-1.55) {\small $7$};
\node[popmuted] at (-2.9,-2.3) {\small $3$};
\end{tikzpicture}
\end{center}
\legende{Diagramme de Venn à trois ensembles : football, handball, athlétisme. La région centrale ($3$ élèves) est remplie en premier.}
\end{popfigure}

\courssub{Tableaux à double entrée}

Lorsqu'on croise \textbf{deux} critères (et non trois), le \cle{tableau à double entrée} est souvent plus rapide et plus lisible qu'un diagramme de Venn. Chaque case contient l'effectif d'une catégorie croisée, et les marges donnent les totaux.

\begin{exemplebox}[Répartition des élèves par sexe et par série]
Un lycée de Bafoussam compte $200$ élèves en Première, répartis entre les séries C et D. On sait qu'il y a $110$ garçons, $90$ filles, $120$ élèves en série C, et que $70$ garçons sont en série C. Complétons le tableau :
\begin{center}
\begin{tabularx}{\linewidth}{|l|X|X|X|}
\hline
\rowcolor{popblueL} & \textbf{Série C} & \textbf{Série D} & \textbf{Total} \\
\hline
\textbf{Garçons} & $70$ & $110 - 70 = 40$ & $110$ \\
\hline
\textbf{Filles} & $120 - 70 = 50$ & $90 - 50 = 40$ & $90$ \\
\hline
\textbf{Total} & $120$ & $200 - 120 = 80$ & $200$ \\
\hline
\end{tabularx}
\end{center}
Chaque case se déduit des autres par simple soustraction : le tableau « s'équilibre » tout seul, comme un sudoku. On lit par exemple que $40$ filles sont en série D. Vérification croisée : $40 + 40 = 80$ élèves en série D, conforme au total.
\end{exemplebox}

\begin{attention}[Venn ou tableau ?]
Le tableau à double entrée ne fonctionne que pour croiser \textbf{deux} critères. Dès que trois critères entrent en jeu (comme les trois sports ci-dessus), il faut revenir au diagramme de Venn. Choisir le bon outil fait déjà partie du savoir-faire exigé.
\end{attention}

\courssub{Produit cartésien et principe multiplicatif}

Revenons à la cantine scolaire : un menu est un \emph{triplet} (plat ; accompagnement ; boisson). Compter les menus, c'est compter les éléments d'un \cle{produit cartésien}.

\begin{definition}[Produit cartésien]
Le produit cartésien de deux ensembles $E$ et $F$, noté $E \times F$, est l'ensemble de tous les couples $(x ; y)$ avec $x \in E$ et $y \in F$ :
\[ E \times F = \bigl\{ (x ; y) \mid x \in E,\ y \in F \bigr\}. \]
\end{definition}

\begin{propriete}[Cardinal du produit cartésien — admise]
Si $E$ et $F$ sont deux ensembles finis, alors :
\[ \operatorname{card}(E \times F) = \operatorname{card}(E) \times \operatorname{card}(F). \]
Cette propriété est admise après l'illustration ci-dessous ; elle se généralise à un produit de $p$ ensembles : les cardinaux se multiplient.
\end{propriete}

\begin{exemplebox}[Illustration : les menus de la cantine]
Notons $P = \{\text{ndolé} ; \text{poulet DG} ; \text{eru} ; \text{poisson braisé}\}$ les $4$ plats, $A = \{\text{riz} ; \text{plantain} ; \text{bâton de manioc}\}$ les $3$ accompagnements. Un choix (plat ; accompagnement) est un couple de $P \times A$. On peut disposer tous ces couples dans un tableau à $4$ lignes et $3$ colonnes : il y a exactement $4 \times 3 = 12$ cases, donc $12$ couples possibles. En ajoutant les $2$ boissons (eau, jus de foléré), chaque couple se dédouble en $2$ menus :
\[ 4 \times 3 \times 2 = 24 \text{ menus différents.} \]
\end{exemplebox}

\begin{aretenir}[Principe additif et principe multiplicatif]
Deux grands réflexes gouvernent tout le dénombrement :
\begin{itemize}
\item \textbf{Principe additif} : si une situation se décompose en cas \emph{incompatibles} (qui ne peuvent pas se produire en même temps), on \textbf{additionne} les effectifs des cas. Mot-clé dans l'énoncé : « \emph{ou} ».
\item \textbf{Principe multiplicatif} : si une situation se décompose en \emph{étapes successives}, chaque étape offrant un certain nombre de choix, on \textbf{multiplie} les nombres de choix. Mots-clés dans l'énoncé : « \emph{puis} », « \emph{et} ».
\end{itemize}
Exemple : composer un menu, c'est choisir un plat \emph{puis} un accompagnement \emph{puis} une boisson : $4 \times 3 \times 2 = 24$. Choisir un élève de série C \emph{ou} de série D parmi des groupes disjoints : on additionne.
\end{aretenir}

\begin{exoresolu}[Le code PIN du téléphone]
\textbf{Énoncé.} Combien existe-t-il de codes PIN à $4$ chiffres (chaque chiffre de $0$ à $9$, répétitions autorisées, un code pouvant commencer par $0$) ?
\tcblower
\textbf{Solution détaillée.} Composer un code PIN, c'est effectuer $4$ choix successifs : le $1^{\text{er}}$ chiffre, \emph{puis} le $2^{\text{e}}$, \emph{puis} le $3^{\text{e}}$, \emph{puis} le $4^{\text{e}}$. Chaque choix offre $10$ possibilités, indépendamment des précédents. Par le principe multiplicatif :
\[ 10 \times 10 \times 10 \times 10 = 10^4 = \num{10000} \text{ codes.} \]
Il s'agit exactement du cardinal du produit cartésien $\{0 ; 1 ; \dots ; 9\}^4$. \textbf{Conclusion :} il existe $\num{10000}$ codes PIN à $4$ chiffres, de $0000$ à $9999$.
\end{exoresolu}

\begin{savaistu}[Le dénombrement, gardien de ta vie numérique]
La sécurité des codes PIN, des mots de passe et des transactions Mobile Money (MTN MoMo, Orange Money) repose entièrement sur le dénombrement : plus le nombre de combinaisons possibles est grand, plus un attaquant mettra de temps à toutes les essayer. Un code à $4$ chiffres offre $\num{10000}$ possibilités — cassable en quelques minutes par une machine —, tandis qu'un mot de passe de $8$ caractères parmi $62$ symboles offre $62^8 \approx 2{,}2 \times 10^{14}$ possibilités. Les mathématiques de cette leçon protègent littéralement ton argent !
\end{savaistu}

\begin{exercice}[À toi de jouer]
Dans un lycée de Douala, on a interrogé les $80$ élèves d'une classe de Première sur leurs clubs : $45$ sont inscrits au club informatique, $38$ au club théâtre, et $15$ aux deux clubs.
\begin{enumerate}
\item Combien d'élèves sont inscrits uniquement au club informatique ?
\item Combien d'élèves sont inscrits à au moins un des deux clubs ? En déduire le nombre d'élèves inscrits à aucun des deux clubs.
\item Un menu du réfectoire se compose d'une entrée et d'un plat. Avec $3$ entrées et $5$ plats, combien de menus différents peut-on former ?
\end{enumerate}
\end{exercice}

\begin{corrige}[Exercice 1]
\begin{enumerate}
\item Informatique seul : $45 - 15 = 30$ élèves.
\item $\operatorname{card}(I \cup T) = 45 + 38 - 15 = 68$ élèves inscrits à au moins un club. Par le complémentaire : $80 - 68 = 12$ élèves ne sont inscrits à aucun des deux clubs.
\item Principe multiplicatif : $3 \times 5 = 15$ menus différents.
\end{enumerate}
\end{corrige}

\begin{aretenir}[L'essentiel de la section]
\begin{itemize}
\item Le \cle{cardinal} d'un ensemble fini est son nombre d'éléments.
\item $\operatorname{card}(A \cup B) = \operatorname{card}(A) + \operatorname{card}(B) - \operatorname{card}(A \cap B)$ ; on additionne sans retrancher \emph{uniquement} si les ensembles sont disjoints.
\item $\operatorname{card}(\bar{A}) = \operatorname{card}(E) - \operatorname{card}(A)$ : penser au complémentaire quand le décompte direct est compliqué.
\item Diagramme de Venn : on remplit \textbf{toujours} en commençant par l'intersection centrale, puis par soustractions successives.
\item Deux critères : tableau à double entrée ; trois critères : diagramme de Venn.
\item $\operatorname{card}(E \times F) = \operatorname{card}(E) \times \operatorname{card}(F)$ : des étapes successives se multiplient, des cas incompatibles s'additionnent.
\end{itemize}
\end{aretenir}

\coursec{Arbres de choix et p-uplets}

Dans la section précédente, nous avons appris à compter les éléments d'un ensemble fini à l'aide du cardinal, des diagrammes de Venn et des tableaux à double entrée. Ces outils sont parfaits pour décrire des \emph{situations} (combien d'élèves pratiquent tel sport, combien parlent telle langue). Mais ils atteignent vite leurs limites dès qu'il s'agit de compter des \emph{possibilités} issues de choix successifs : combien de menus différents peut-on composer dans un restaurant de Yaoundé ? Combien de codes PIN à 4 chiffres existe-t-il pour ta carte SIM Orange ou MTN ? Combien de classements possibles pour les 8 finalistes d'une course ?

Pour répondre à ces questions, il nous faut un nouvel outil : l'\cle{arbre de choix}, et un nouveau concept fondamental, le \cle{p-uplet} (ou \cle{p-liste}). C'est toute l'ambition de cette section.

\courssub{Les arbres de choix}

\textbf{Une situation pour commencer.} Une cantine scolaire propose à ses élèves de composer leur menu en choisissant : un plat parmi 2 (ndolé ou poulet DG), un accompagnement parmi 2 (riz ou plantain), et une boisson parmi 2 (eau ou jus). Combien de menus différents un élève peut-il composer ?

On pourrait essayer d'énumérer les menus à la main, mais on risque vite d'en oublier ou d'en compter deux fois. La bonne idée consiste à \emph{visualiser} les choix successifs sous forme d'un arbre : chaque choix correspond à un étage, chaque option à une branche, et chaque menu complet à un \emph{chemin} allant de la racine à une feuille.

\begin{popfigure}
\begin{center}
\begin{tikzpicture}[
  grow=right,
  level distance=3.2cm,
  level 1/.style={sibling distance=4.4cm},
  level 2/.style={sibling distance=2.2cm},
  level 3/.style={sibling distance=1.1cm},
  every node/.style={font=\small},
  edge from parent/.style={draw=popmuted, thick}
]
\node[fill=popink, text=white, rounded corners, inner sep=3pt] {Menu}
  child { node[fill=poporangeL, rounded corners, inner sep=2pt] {Poulet DG}
    child { node[fill=popgreenL, rounded corners, inner sep=2pt] {Plantain}
      child { node[fill=popblueL, rounded corners, inner sep=2pt] {Jus} }
      child { node[fill=popblueL, rounded corners, inner sep=2pt] {Eau} }
    }
    child { node[fill=popgreenL, rounded corners, inner sep=2pt] {Riz}
      child { node[fill=popblueL, rounded corners, inner sep=2pt] {Jus} }
      child { node[fill=popblueL, rounded corners, inner sep=2pt] {Eau} }
    }
  }
  child { node[fill=poporangeL, rounded corners, inner sep=2pt] {Ndolé}
    child { node[fill=popgreenL, rounded corners, inner sep=2pt] {Plantain}
      child { node[fill=popblueL, rounded corners, inner sep=2pt] {Jus} }
      child { node[fill=popblueL, rounded corners, inner sep=2pt] {Eau} }
    }
    child { node[fill=popgreenL, rounded corners, inner sep=2pt] {Riz}
      child { node[fill=popblueL, rounded corners, inner sep=2pt] {Jus} }
      child { node[fill=popblueL, rounded corners, inner sep=2pt] {Eau} }
    }
  };
\end{tikzpicture}
\end{center}
\legende{Arbre de choix pour composer un menu : 2 plats $\times$ 2 accompagnements $\times$ 2 boissons $= 8$ menus possibles.}
\end{popfigure}

\begin{definition}[Arbre de choix]
Un \cle{arbre de choix} est un schéma qui représente des choix successifs :
\begin{itemize}
\item chaque \emph{niveau} (ou étage) de l'arbre correspond à un choix ;
\item chaque \emph{branche} issue d'un nœud correspond à une option possible pour ce choix ;
\item chaque \emph{chemin} allant de la racine à une feuille (extrémité) correspond à une issue complète de la succession de choix.
\end{itemize}
Le nombre total d'issues est le nombre de chemins de la racine vers les feuilles.
\end{definition}

Sur notre exemple, on compte 8 feuilles, donc 8 menus. Mais on n'a pas besoin de dessiner l'arbre en entier pour trouver ce nombre : à chacun des 2 plats correspondent 2 accompagnements, soit $2 \times 2 = 4$ couples (plat, accompagnement) ; et à chacun de ces 4 couples correspondent 2 boissons, soit $4 \times 2 = 8$ menus. C'est le \emph{principe multiplicatif}.

\begin{propriete}[Principe multiplicatif]
Si une situation comporte $p$ choix successifs, le premier offrant $n_1$ possibilités, le deuxième $n_2$ possibilités, \dots, le $p$-ième $n_p$ possibilités, alors le nombre total d'issues est :
\[ n_1 \times n_2 \times \cdots \times n_p. \]
\end{propriete}

\begin{attention}[Le principe multiplicatif exige l'indépendance des nombres de choix]
Le principe multiplicatif s'applique lorsque le nombre d'options à chaque étape \emph{ne dépend pas} des choix précédents. Si le nombre de boissons disponibles dépendait du plat choisi, il faudrait dessiner l'arbre en entier et compter les feuilles une à une. Vérifie toujours cette condition avant de multiplier !
\end{attention}

\courssub{Les p-uplets (ou p-listes)}

Formalisons ce que nous venons de faire. Un menu est une \emph{liste ordonnée} de 3 éléments : (plat, accompagnement, boisson). L'ordre des cases est fixé par la nature du problème : (ndolé, riz, eau) et (riz, ndolé, eau) ne désignent pas le même objet, car la première case est réservée au plat. En mathématiques, une telle liste s'appelle un \emph{p-uplet}.

\begin{definition}[p-uplet d'un ensemble]
Soit $E$ un ensemble non vide et $p$ un entier naturel non nul. On appelle \cle{p-uplet} (ou \cle{p-liste}) d'éléments de $E$ toute liste ordonnée de $p$ éléments de $E$ :
\[ (x_1, x_2, \ldots, x_p) \quad \text{où } x_1 \in E,\ x_2 \in E,\ \ldots,\ x_p \in E. \]
L'ensemble de tous les p-uplets d'éléments de $E$ se note $E^p$ (produit cartésien $E \times E \times \cdots \times E$, $p$ fois).
\end{definition}

\begin{aretenir}[Deux caractéristiques essentielles d'un p-uplet]
\begin{itemize}
\item \textbf{L'ordre compte} : $(a, b) \neq (b, a)$ dès que $a \neq b$. Un p-uplet n'est pas un ensemble : $\{a, b\} = \{b, a\}$ mais $(a,b) \neq (b,a)$.
\item \textbf{La répétition est autorisée} (sauf mention contraire) : $(a, a, b)$ est un 3-uplet parfaitement valide, comme le code PIN 1125.
\end{itemize}
\end{aretenir}

\begin{exemplebox}[Premiers exemples de p-uplets]
\begin{itemize}
\item Un code PIN est un 4-uplet d'éléments de $E = \{0, 1, 2, \ldots, 9\}$ : par exemple $(1, 9, 9, 4)$.
\item Un mot de 5 lettres (ayant un sens ou non) est un 5-uplet d'éléments de l'alphabet.
\item Le résultat de 3 lancers successifs d'un dé est un 3-uplet d'éléments de $\{1, 2, 3, 4, 5, 6\}$.
\end{itemize}
\end{exemplebox}

\courssub{p-uplets avec répétition : le nombre $n^p$}

\begin{propriete}[Nombre de p-uplets d'un ensemble à n éléments]
Soit $E$ un ensemble à $n$ éléments ($n \geq 1$) et $p$ un entier naturel non nul. Le nombre de p-uplets d'éléments de $E$ est :
\[ \operatorname{Card}(E^p) = n^p. \]
\emph{Cette démonstration est formatrice et peut être demandée au Probatoire sous forme guidée.}
\end{propriete}

\begin{experience}[Démonstration par le principe multiplicatif]
Construire un p-uplet $(x_1, x_2, \ldots, x_p)$ d'éléments de $E$, c'est effectuer $p$ choix successifs :
\begin{enumerate}
\item choisir le premier élément $x_1$ : il y a $n$ possibilités (tout élément de $E$ convient) ;
\item choisir le deuxième élément $x_2$ : il y a encore $n$ possibilités, car la répétition est autorisée (on peut reprendre le même élément) ;
\item[$\vdots$]
\item[$p$.] choisir le $p$-ième élément $x_p$ : il y a toujours $n$ possibilités.
\end{enumerate}
Le nombre d'options à chaque étape ne dépend pas des choix précédents : le principe multiplicatif s'applique et donne
\[ \underbrace{n \times n \times \cdots \times n}_{p \text{ facteurs}} = n^p. \]
\end{experience}

\begin{exoresolu}[Codes PIN et codes à chiffres distincts]
Une carte SIM est protégée par un code PIN de 4 chiffres (chaque chiffre est choisi dans $\{0, 1, \ldots, 9\}$).
\begin{enumerate}
\item Combien de codes PIN différents existe-t-il ?
\item Combien de codes PIN ont leurs 4 chiffres deux à deux distincts ?
\end{enumerate}
\tcblower
\textbf{1. Codes PIN quelconques.} Un code PIN est un 4-uplet d'éléments d'un ensemble à $n = 10$ éléments, avec répétition possible (le code 2233 est valide). Le nombre de codes est donc :
\[ 10^4 = 10 \times 10 \times 10 \times 10 = 10\,000. \]
Il existe \textbf{10 000 codes PIN} différents, de 0000 à 9999.

\textbf{2. Codes à chiffres distincts.} Cette fois, la répétition est interdite : c'est un tirage \emph{sans remise}. On applique le principe multiplicatif étape par étape :
\begin{itemize}
\item 1er chiffre : 10 possibilités ;
\item 2e chiffre : 9 possibilités (on ne peut pas reprendre le premier) ;
\item 3e chiffre : 8 possibilités ;
\item 4e chiffre : 7 possibilités.
\end{itemize}
Le nombre de codes est donc :
\[ 10 \times 9 \times 8 \times 7 = 5040. \]
Il existe \textbf{5040 codes PIN} à chiffres deux à deux distincts. Remarque au passage : interdire la répétition a presque divisé le nombre de codes par deux.
\end{exoresolu}

\courssub{p-uplets d'éléments distincts (sans répétition)}

Dans de nombreuses situations, un élément déjà choisi ne peut plus l'être : on ne peut pas attribuer deux fois la même médaille, on ne peut pas tirer deux fois la même carte d'un paquet, on ne peut pas élire deux fois la même personne à deux postes distincts. On parle alors de p-uplets \emph{d'éléments distincts}, ou de tirages \emph{sans remise}.

\begin{propriete}[Nombre de p-uplets d'éléments distincts]
Soit $E$ un ensemble à $n$ éléments et $p$ un entier tel que $1 \leq p \leq n$. Le nombre de p-uplets d'éléments de $E$ \emph{deux à deux distincts} est :
\[ n(n-1)(n-2)\cdots(n-p+1). \]
C'est un produit de $p$ facteurs entiers consécutifs, en partant de $n$ et en descendant.
\end{propriete}

\begin{experience}[Démonstration guidée par l'arbre]
Imaginons l'arbre de choix correspondant à la construction d'un p-uplet d'éléments distincts.
\begin{enumerate}
\item \textbf{Niveau 1} (choix de $x_1$) : la racine a $n$ branches, car tout élément de $E$ est disponible.
\item \textbf{Niveau 2} (choix de $x_2$) : de chaque nœud du niveau 1 partent $n-1$ branches, car l'élément déjà choisi est désormais interdit. Il y a donc $n(n-1)$ chemins à deux étages.
\item \textbf{Niveau 3} (choix de $x_3$) : deux éléments sont déjà utilisés, il reste $n-2$ options par nœud, soit $n(n-1)(n-2)$ chemins.
\item[$\vdots$]
\item \textbf{Niveau $p$} (choix de $x_p$) : $p-1$ éléments sont déjà utilisés, il reste donc $n-(p-1) = n-p+1$ options.
\end{enumerate}
Par le principe multiplicatif, le nombre total de chemins est :
\[ n \times (n-1) \times (n-2) \times \cdots \times (n-p+1), \]
un produit de exactement $p$ facteurs. La condition $p \leq n$ est indispensable : sinon, à un moment, il ne reste plus aucun élément disponible et le produit vaut 0, ce qui traduit l'impossibilité de la tâche.
\end{experience}

\begin{methode}[Compter un produit $n(n-1)\cdots(n-p+1)$ sans se tromper]
Pour calculer correctement ce produit :
\begin{enumerate}
\item repère $n$ (le nombre d'éléments de départ) et $p$ (la longueur du p-uplet) ;
\item écris $p$ facteurs, ni plus, ni moins, en partant de $n$ et en diminuant de 1 à chaque fois ;
\item vérifie que le dernier facteur est bien $n - p + 1$ (et non $n - p$).
\end{enumerate}
Par exemple, pour $n = 10$ et $p = 4$ : les facteurs sont $10, 9, 8, 7$ ; le dernier est $10 - 4 + 1 = 7$. Correct.
\end{methode}

\begin{exemplebox}[Podium d'une course]
À la finale du 100 m des jeux scolaires, 8 athlètes sont au départ. Un podium est un 3-uplet (or, argent, bronze) d'athlètes distincts : l'ordre compte et nul ne peut occuper deux places. Le nombre de podiums possibles est :
\[ 8 \times 7 \times 6 = 336. \]
\end{exemplebox}

\courssub{Avec ou sans répétition ? La question décisive}

Toute la difficulté pratique du dénombrement tient en une question : \emph{un élément déjà choisi peut-il être choisi à nouveau ?} La réponse dépend de la nature physique de la situation, pas de tes préférences.

\begin{popfigure}
\begin{center}
\begin{tikzpicture}[
  grow=right,
  level distance=2.6cm,
  level 1/.style={sibling distance=3cm},
  level 2/.style={sibling distance=1.6cm},
  every node/.style={font=\small},
  edge from parent/.style={draw=popmuted, thick}
]
\node[fill=popink, text=white, rounded corners, inner sep=3pt] {Code à 3 chiffres}
  child { node[fill=popblueL, rounded corners, inner sep=2pt] {0--9}
    child { node[fill=popgreenL, rounded corners, inner sep=2pt] {10 branches} }
    child { node[fill=popgreenL, rounded corners, inner sep=2pt] {$\cdots$} }
    child { node[fill=popgreenL, rounded corners, inner sep=2pt] {10 branches} }
  }
  child { node[fill=popblueL, rounded corners, inner sep=2pt] {$\vdots$} edge from parent[draw=none] }
  child { node[fill=popblueL, rounded corners, inner sep=2pt] {0--9}
    child { node[fill=popgreenL, rounded corners, inner sep=2pt] {10 branches} }
    child { node[fill=popgreenL, rounded corners, inner sep=2pt] {$\cdots$} }
    child { node[fill=popgreenL, rounded corners, inner sep=2pt] {10 branches} }
  };
\node[font=\small\itshape, text=popdark] at (8.2,-3.4) {Total : $10 \times 10 \times 10 = 10^3 = 1000$ codes};
\end{tikzpicture}
\end{center}
\legende{Arbre partiel d'un code à 3 chiffres : chaque nœud a 10 branches, car la répétition est autorisée. On ne dessine que quelques branches, mais on sait les compter toutes.}
\end{popfigure}

\begin{methode}[Choisir entre « avec » et « sans » répétition]
Face à un problème de dénombrement de listes, pose-toi dans l'ordre les questions suivantes :
\begin{enumerate}
\item \textbf{L'ordre compte-t-il ?} Si oui, on est dans le monde des p-uplets (cette section). Si non, il faudra les combinaisons (section 4).
\item \textbf{Un élément peut-il être réutilisé ?} Demande-toi si la situation l'autorise physiquement : un chiffre d'un code peut réapparaître (avec remise) ; une même personne ne peut pas être à la fois première et deuxième (sans remise) ; une boule remise dans l'urne peut être retirée (avec remise), sinon non.
\item \textbf{Conclusion} : avec répétition $\rightarrow$ $n^p$ ; sans répétition $\rightarrow$ $n(n-1)\cdots(n-p+1)$.
\end{enumerate}
Le vocabulaire des urnes est éclairant : tirage \emph{avec remise} $=$ répétition possible ; tirage \emph{sans remise} $=$ éléments distincts.
\end{methode}

\begin{center}
\begin{tabularx}{\linewidth}{|l|X|X|}
\hline
\rowcolor{popblueL} \textbf{Situation} & \textbf{Répétition ?} & \textbf{Nombre} \\
\hline
Code PIN à 4 chiffres & Oui (1122 valide) & $10^4 = 10\,000$ \\
\hline
Podium de 8 athlètes & Non (un athlète, une place) & $8 \times 7 \times 6 = 336$ \\
\hline
Mot de 3 lettres & Oui & $26^3 = 17\,576$ \\
\hline
Bureau (président, trésorier) parmi 15 élèves & Non (deux postes, deux personnes) & $15 \times 14 = 210$ \\
\hline
3 lancers successifs d'un dé & Oui & $6^3 = 216$ \\
\hline
\end{tabularx}
\end{center}

\begin{attention}[L'erreur classique : confondre $n^p$ et $p^n$]
C'est l'une des erreurs les plus fréquentes au Probatoire. Retiens bien :
\begin{itemize}
\item $n$ est le nombre d'\textbf{éléments} de l'ensemble $E$ (le nombre d'options par case) ;
\item $p$ est la \textbf{longueur} de la liste (le nombre de cases).
\end{itemize}
Le nombre de p-uplets est $n^p$ : \emph{c'est le nombre d'options qui est élevé à la puissance « nombre de cases »}. Ainsi, le nombre de codes à 4 chiffres est $10^4 = 10\,000$, et surtout pas $4^{10} = 1\,048\,576$ ! Vérifie toujours sur un petit cas : avec $n = 2$ symboles et $p = 3$ cases, on liste $(0,0,0), (0,0,1), \ldots$ : on trouve $8 = 2^3$ listes, et non $3^2 = 9$.
\end{attention}

\begin{savaistu}[Les plaques d'immatriculation, un dénombrement roulant]
Une plaque d'immatriculation camerounaise de la série actuelle comporte des lettres et des chiffres disposés dans un ordre précis. Si un format comporte 2 lettres suivies de 3 chiffres, le nombre de plaques possibles est $26^2 \times 10^3 = 676 \times 1000 = 676\,000$. Lorsque ce stock s'épuise avec l'augmentation du parc automobile, il suffit d'ajouter un caractère pour multiplier la capacité par 10 ou 26 : c'est toute la puissance (exponentielle !) du principe multiplicatif. Le même principe explique pourquoi les opérateurs téléphoniques allongent les numéros quand les abonnés se multiplient.
\end{savaistu}

\begin{exercice}[Menus, mots de passe et classements]
\begin{enumerate}
\item Un restaurant de Douala propose 3 entrées, 4 plats principaux et 2 desserts. Combien de menus complets (entrée, plat, dessert) peut-on composer ?
\item Un mot de passe est formé de 6 caractères choisis parmi les 26 lettres minuscules. Combien de mots de passe sont possibles ? Combien ont leurs 6 lettres deux à deux distinctes ?
\item Dans une classe de 20 élèves, on désigne un délégué, un délégué adjoint et un trésorier (trois élèves différents). Combien de bureaux différents peut-on former ?
\end{enumerate}
\end{exercice}

\begin{corrige}[Exercice : Menus, mots de passe et classements]
\textbf{1.} Un menu est un 3-uplet (entrée, plat, dessert). Par le principe multiplicatif : $3 \times 4 \times 2 = 24$ menus.

\textbf{2.} Avec répétition possible : $26^6 = 308\,915\,776$ mots de passe. Avec lettres distinctes : $26 \times 25 \times 24 \times 23 \times 22 \times 21 = 165\,765\,600$ mots de passe.

\textbf{3.} Un bureau est un 3-uplet d'élèves distincts (l'ordre compte car les postes sont différents, et un élève ne peut cumuler deux postes) : $20 \times 19 \times 18 = 6840$ bureaux possibles.
\end{corrige}

\begin{aretenir}[L'essentiel de la section]
\begin{itemize}
\item Un \textbf{arbre de choix} représente des choix successifs ; chaque chemin de la racine à une feuille est une issue.
\item \textbf{Principe multiplicatif} : $n_1 \times n_2 \times \cdots \times n_p$ issues pour $p$ choix successifs.
\item Un \textbf{p-uplet} est une liste ordonnée de $p$ éléments ; l'ordre compte.
\item Avec répétition : $n^p$ p-uplets. Sans répétition : $n(n-1)\cdots(n-p+1)$ p-uplets d'éléments distincts ($p$ facteurs, $p \leq n$).
\item La question décisive : \emph{un élément peut-il être réutilisé ?}
\end{itemize}
Dans la prochaine section, nous donnerons un nom et une notation au produit $n(n-1)\cdots(n-p+1)$ : ce seront la \emph{factorielle} et les \emph{arrangements}.
\end{aretenir}

\coursec{Factorielle, permutations et arrangements}

Dans la section précédente, tu as appris à compter des $p$-uplets grâce aux arbres de choix : avec répétition, on obtient $n^p$ listes ; sans répétition, on obtient $n(n-1)(n-2)\cdots$ jusqu'à épuisement des $p$ choix. Ces produits d'entiers consécutifs décroissants apparaissent si souvent en dénombrement que les mathématiciens leur ont donné un nom et une notation : la \cle{factorielle}. Cette section lui est consacrée, ainsi qu'à ses deux grandes applications : les \cle{permutations} (ranger $n$ objets en file) et les \cle{arrangements} (ranger $p$ objets choisis parmi $n$, dans un ordre qui compte). Tu verras que les podiums d'athlétisme, les classements au Probatoire et les anagrammes relèvent tous du même outil.

\courssub{La factorielle d'un entier naturel}

\textbf{Une notation pour un produit qui revient sans cesse.}\par

Imagine que tu doives ranger 4 livres différents sur une étagère, de gauche à droite. Tu as 4 choix pour la première place, puis 3 pour la deuxième, puis 2, puis 1 : au total $4\times 3\times 2\times 1 = 24$ rangements. Avec 10 livres, le calcul devient $10\times 9\times 8\times\cdots\times 2\times 1$, un produit long à écrire. La factorielle est exactement l'abréviation de ce produit.

\begin{definition}[Factorielle d'un entier naturel]
Soit $n$ un entier naturel non nul. On appelle \cle{factorielle} de $n$, et on note $n!$ (lire « factorielle $n$ » ou « $n$ factorielle »), le produit de tous les entiers de $1$ à $n$ :
\[ n! = n\times (n-1)\times (n-2)\times \cdots \times 2\times 1. \]
Par convention, on pose $\boxed{0! = 1}$.
\end{definition}

\begin{exemplebox}[Premières valeurs]
$1! = 1$ ; \quad $2! = 2\times 1 = 2$ ; \quad $3! = 3\times 2\times 1 = 6$ ; \quad $4! = 24$ ; \quad $5! = 120$ ; \quad $6! = 720$.
\end{exemplebox}

\begin{aretenir}[Valeurs de $n!$ pour $n$ de 0 à 10]
\begin{center}
\begin{tabularx}{\linewidth}{|l|X|X|X|X|X|X|}\hline
\rowcolor{popblueL} $n$ & 0 & 1 & 2 & 3 & 4 & 5 \\ \hline
$n!$ & 1 & 1 & 2 & 6 & 24 & 120 \\ \hline
\rowcolor{popblueL} $n$ & 6 & 7 & 8 & 9 & 10 & \\ \hline
$n!$ & 720 & \num{5040} & \num{40320} & \num{362880} & \num{3628800} & \\ \hline
\end{tabularx}
\end{center}
La factorielle \emph{explose} très vite : $10!$ dépasse déjà trois millions. C'est pourquoi on préfère souvent garder l'écriture factorielle plutôt que de tout calculer.
\end{aretenir}

\begin{savaistu}[Pourquoi $0! = 1$ ?]
Cette convention n'est pas arbitraire : elle garantit que toutes les formules de dénombrement restent vraies aux cas limites. Par exemple, il existe exactement \emph{une} façon de ranger zéro objet (ne rien faire !), et la formule des arrangements que tu verras plus loin exige $0! = 1$ pour donner le bon résultat quand $p = n$.
\end{savaistu}

\courssub{Calculer et simplifier avec les factorielles}

La propriété fondamentale de la factorielle découle directement de sa définition.

\begin{propriete}[Relation de récurrence]
Pour tout entier $n \geq 1$ :
\[ (n+1)! = (n+1)\times n! \qquad\text{et de manière équivalente}\qquad n! = n\times (n-1)!. \]
\end{propriete}

En effet, $(n+1)! = (n+1)\times \underbrace{n\times (n-1)\times\cdots\times 1}_{=\,n!}$. Cette relation permet de « déplier » une factorielle autant que nécessaire.

\begin{methode}[Simplifier un quotient de factorielles sans tout calculer]
Pour simplifier $\dfrac{n!}{p!}$ avec $n \geq p$ :
\begin{enumerate}
\item Repérer la \emph{plus grande} des deux factorielles (ici $n!$).
\item La déplier jusqu'à faire apparaître la plus petite : $n! = n(n-1)\cdots(p+1)\times p!$.
\item Simplifier par $p!$ au numérateur et au dénominateur.
\item Il reste le produit des entiers de $p+1$ à $n$.
\end{enumerate}
Exemple fondateur : $\dfrac{10!}{8!} = \dfrac{10\times 9\times 8!}{8!} = 10\times 9 = 90$, car le facteur $8!$ apparaît en même temps au numérateur et au dénominateur.
\end{methode}

\begin{exemplebox}[Simplifications types]
\begin{align*}
\frac{12!}{10!} &= 12\times 11 = 132 ; &
\frac{7!}{4!} &= 7\times 6\times 5 = 210 ; &
\frac{(n+1)!}{(n-1)!} &= (n+1)\times n = n^2 + n.
\end{align*}
\end{exemplebox}

\begin{attention}[Erreurs fréquentes avec les factorielles]
\begin{itemize}
\item La factorielle n'est \textbf{pas distributive} : en général $(n+p)! \neq n! + p!$. Par exemple $(2+3)! = 5! = 120$, alors que $2! + 3! = 2 + 6 = 8$.
\item De même, $(2n)! \neq 2\times n!$ : $(2\times 3)! = 6! = 720$ mais $2\times 3! = 12$.
\item On ne peut simplifier que des \emph{produits} : $\dfrac{n!+1}{n!}$ ne se simplifie pas en $1$ ; on écrit $\dfrac{n!+1}{n!} = 1 + \dfrac{1}{n!}$.
\end{itemize}
\end{attention}

\courssub{Permutations de $n$ éléments}

\textbf{Intuition.}\par Trois coureurs — Amadou, Béatrice et Chantal — arrivent à la fin d'un relais. De combien de façons peut-on les classer du premier au troisième rang, sans ex æquo ? Pour la première place : 3 choix ; pour la deuxième : 2 choix restants ; pour la troisième : 1 seul. Total : $3\times 2\times 1 = 3! = 6$ classements. Ranger $n$ objets distincts dans un ordre complet, c'est exactement construire un $n$-uplet sans répétition utilisant tous les objets.

\begin{definition}[Permutation]
Soit $E$ un ensemble fini à $n$ éléments. On appelle \cle{permutation} de $E$ tout $n$-uplet d'éléments \emph{distincts} de $E$, c'est-à-dire toute façon d'ordonner la totalité des $n$ éléments.
\end{definition}

\begin{propriete}[Nombre de permutations]
Le nombre de permutations d'un ensemble à $n$ éléments est
\[ n! = n\times (n-1)\times \cdots \times 2\times 1. \]
\end{propriete}

\begin{experience}[Démonstration guidée]
Construisons une permutation de $n$ éléments, place par place :
\begin{enumerate}
\item Pour la première place : $n$ choix possibles.
\item Pour la deuxième place : les éléments étant distincts, il reste $n-1$ choix.
\item On continue ainsi : à la $k$-ième place, il reste $n-k+1$ choix.
\item Pour la dernière place : $1$ seul choix.
\end{enumerate}
Par le principe multiplicatif (arbre de choix), le nombre total est $n\times(n-1)\times\cdots\times 2\times 1 = n!$. \hfill$\blacksquare$
\end{experience}

\begin{popfigure}
\begin{center}
\begin{tikzpicture}[grow=right, level distance=2.6cm,
  level 1/.style={sibling distance=3.4cm},
  level 2/.style={sibling distance=1.6cm},
  level 3/.style={sibling distance=0.75cm},
  every node/.style={font=\small}]
\node {Départ}
  child { node[popblue] {A}
    child { node[popblue] {B} child { node[popdark] {C} } }
    child { node[popblue] {C} child { node[popdark] {B} } }
  }
  child { node[poporange] {B}
    child { node[poporange] {A} child { node[popdark] {C} } }
    child { node[poporange] {C} child { node[popdark] {A} } }
  }
  child { node[popgreen] {C}
    child { node[popgreen] {A} child { node[popdark] {B} } }
    child { node[popgreen] {B} child { node[popdark] {A} } }
  };
\end{tikzpicture}
\end{center}
\legende{Arbre des $3! = 6$ permutations des coureurs A (Amadou), B (Béatrice), C (Chantal). Chaque chemin de la racine à une feuille donne un classement complet : ABC, ACB, BAC, BCA, CAB, CBA.}
\end{popfigure}

\begin{exemplebox}[Anagrammes]
Combien de mots de 4 lettres (ayant un sens ou non) peut-on former avec les lettres A, N, G, O utilisées chacune exactement une fois ? C'est le nombre de permutations de 4 lettres distinctes : $4! = 24$. Le mot ANGO en est une, GONA une autre.
\end{exemplebox}

\courssub{Arrangements de $p$ éléments parmi $n$}

\textbf{Intuition.}\par Reprenons nos athlètes, mais cette fois 8 athlètes disputent une finale et seules les trois premières places (or, argent, bronze) sont récompensées. On ne classe plus \emph{tous} les athlètes : on choisit 3 personnes parmi 8, \emph{et l'ordre compte} (être médaillé d'or n'est pas être médaillé de bronze). C'est la situation typique d'un arrangement.

\begin{definition}[Arrangement]
Soit $E$ un ensemble à $n$ éléments et $p$ un entier tel que $1 \leq p \leq n$. On appelle \cle{arrangement} de $p$ éléments de $E$ tout $p$-uplet $(x_1, x_2, \ldots, x_p)$ d'éléments de $E$ \emph{deux à deux distincts}. Autrement dit : un arrangement, c'est une liste \textbf{ordonnée} de $p$ éléments \textbf{sans répétition} pris parmi $n$.
\end{definition}

\begin{propriete}[Nombre d'arrangements — démonstration exigible]
Le nombre d'arrangements de $p$ éléments parmi $n$, noté $A_n^p$ (ou $A(n,p)$), est
\[ A_n^p = n(n-1)(n-2)\cdots(n-p+1) = \frac{n!}{(n-p)!}. \]
En particulier, $A_n^n = \dfrac{n!}{0!} = n!$ : on retrouve les permutations.
\end{propriete}

\begin{experience}[Démonstration guidée du théorème]
Un arrangement est un $p$-uplet sans répétition : on le construit composante par composante.
\begin{enumerate}
\item Première composante : $n$ choix.
\item Deuxième composante : $n-1$ choix (l'élément déjà placé est interdit).
\item \ldots\ à la $p$-ième composante, $p-1$ éléments sont déjà utilisés : il reste $n-p+1$ choix.
\end{enumerate}
Par le principe multiplicatif :
\[ A_n^p = n(n-1)\cdots(n-p+1). \]
Pour obtenir l'écriture avec les factorielles, on multiplie et on divise par $(n-p)!$ :
\begin{align*}
A_n^p &= n(n-1)\cdots(n-p+1)\times \frac{(n-p)(n-p-1)\cdots 2\times 1}{(n-p)(n-p-1)\cdots 2\times 1} \\
&= \frac{n\times(n-1)\times\cdots\times 2\times 1}{(n-p)!} = \frac{n!}{(n-p)!}. \quad \blacksquare
\end{align*}
\end{experience}

\begin{attention}[Conditions d'application à vérifier avant de foncer]
Avant d'écrire $A_n^p$, pose-toi deux questions :
\begin{itemize}
\item \textbf{L'ordre compte-t-il ?} (podium, classement, mot de passe, bureau avec président/trésorier distincts : oui.)
\item \textbf{Y a-t-il répétition possible ?} Si oui, ce n'est \emph{pas} un arrangement mais une $p$-liste : le nombre est $n^p$ (section précédente).
\end{itemize}
Un arrangement exige : ordre important \emph{et} éléments distincts. Et bien sûr $p \leq n$ : on ne peut pas aligner 5 coureurs distincts sur un podium de 3 places si l'on n'a que 4 coureurs !
\end{attention}

\begin{exoresolu}[Podiums d'une finale d'athlétisme]
\textbf{Énoncé.} Huit athlètes participent à la finale du 100 mètres lors des jeux scolaires de Yaoundé. De combien de façons peut-on attribuer les trois médailles (or, argent, bronze), sans ex æquo ?
\tcblower
\textbf{Solution détaillée.} Attribuer les médailles, c'est choisir 3 athlètes parmi 8 \emph{en tenant compte de l'ordre} (l'or n'est pas le bronze) et sans répétition (un athlète ne peut pas avoir deux médailles). C'est un arrangement de 3 éléments parmi 8 :
\[ A_8^3 = \frac{8!}{(8-3)!} = \frac{8!}{5!} = 8\times 7\times 6 = 336. \]
Il y a donc $\boxed{336}$ podiums possibles. Remarque la simplification : inutile de calculer $8! = \num{40320}$ en entier, le quotient $\frac{8!}{5!}$ se réduit au produit $8\times 7\times 6$.
\end{exoresolu}

\begin{exoresolu}[Bureau d'association]
\textbf{Énoncé.} L'association des élèves d'une classe de Première C compte 12 membres. On doit élire un président, un secrétaire et un trésorier, tous distincts. Combien de bureaux différents peut-on former ?
\tcblower
\textbf{Solution détaillée.} Les trois postes sont \emph{différenciés} : le triplet (président, secrétaire, trésorier) est ordonné, et une même personne ne peut pas cumuler deux postes (pas de répétition). C'est un arrangement :
\[ A_{12}^3 = \frac{12!}{9!} = 12\times 11\times 10 = \num{1320}. \]
Il y a $\num{1320}$ bureaux possibles. Attention : si l'on avait demandé simplement « choisir 3 délégués sans distinction de rôle », l'ordre ne compterait plus et il faudrait utiliser les combinaisons (section suivante) — le résultat serait bien plus petit.
\end{exoresolu}

\begin{methode}[Reconnaître et calculer un arrangement]
\begin{enumerate}
\item Vérifier : ordre important ? éléments sans répétition ? $p \leq n$ ?
\item Identifier $n$ (taille de l'ensemble de départ) et $p$ (nombre de places à pourvoir).
\item Écrire $A_n^p = \dfrac{n!}{(n-p)!}$ puis simplifier : c'est le produit de $p$ facteurs consécutifs décroissants à partir de $n$.
\item Conclure par une phrase : « Il y a donc \ldots possibilités. »
\end{enumerate}
Astuce de vérification : $A_n^p$ est un produit d'\emph{exactement $p$ facteurs}. Ainsi $A_8^3 = 8\times 7\times 6$ : trois facteurs, pas quatre.
\end{methode}

\begin{savaistu}[Les arrangements autour de nous]
Les codes de la loterie, les classements du concours d'entrée à l'ENAM ou à l'ENS, la composition du bureau d'un club, l'ordre de passage des candidats à un oral : toutes ces situations camerounaises quotidiennes sont des arrangements. À l'inverse, un code PIN de téléphone MTN à 4 chiffres autorise les répétitions (par exemple 2244) : c'est une $4$-liste, avec $10^4 = \num{10000}$ possibilités, et non un arrangement. Savoir distinguer ces deux modèles est la compétence clé de ce chapitre.
\end{savaistu}

\begin{exercice}[Pour t'entraîner]
\begin{enumerate}
\item Calculer sans calculatrice : $\dfrac{9!}{7!}$ ; $\dfrac{6!}{3!\times 3!}$ ; $\dfrac{(n+2)!}{n!}$.
\item Combien de façons y a-t-il de ranger 6 élèves sur un banc de 6 places ?
\item Un entraîneur doit désigner, parmi 11 joueurs, un capitaine et un vice-capitaine. Combien de choix possibles ?
\item Combien d'anagrammes le mot CAMEROUN possède-t-il (les 8 lettres sont toutes distinctes) ?
\end{enumerate}
\end{exercice}

\begin{corrige}[Exercice 3]
\begin{enumerate}
\item $\dfrac{9!}{7!} = 9\times 8 = 72$ ; \quad $\dfrac{6!}{3!\times 3!} = \dfrac{720}{6\times 6} = 20$ ; \quad $\dfrac{(n+2)!}{n!} = (n+2)(n+1)$.
\item C'est une permutation de 6 élèves : $6! = 720$ rangements.
\item Ordre important (capitaine $\neq$ vice-capitaine), sans répétition : $A_{11}^2 = 11\times 10 = 110$ choix.
\item Permutations de 8 lettres distinctes : $8! = \num{40320}$ anagrammes.
\end{enumerate}
\end{corrige}

\begin{aretenir}[L'essentiel de la section]
\begin{itemize}
\item $n! = n\times(n-1)\times\cdots\times 2\times 1$ et $0! = 1$ ; relation clé : $(n+1)! = (n+1)\times n!$.
\item Le nombre de \cle{permutations} de $n$ éléments est $n!$.
\item Le nombre d'\cle{arrangements} de $p$ éléments parmi $n$ (ordre important, sans répétition) est
\[ A_n^p = \frac{n!}{(n-p)!} = \underbrace{n(n-1)\cdots(n-p+1)}_{p \text{ facteurs}}. \]
\item Cas particulier : $A_n^n = n!$ (les permutations sont les arrangements « complets »).
\item Simplifie toujours les quotients de factorielles en dépliant la plus grande : $\dfrac{10!}{8!} = 10\times 9$.
\end{itemize}
\end{aretenir}

\coursec{Combinaisons et parties d'un ensemble}

Dans la section précédente, nous avons maîtrisé le dénombrement des listes \textbf{ordonnées} : les $p$-uplets (avec ou sans répétition) et les arrangements. Mais dans de nombreuses situations concrètes — former une équipe, choisir des numéros au loto, constituer un comité — l'ordre de sélection \textbf{n'a aucune importance}. Choisir \{\textsc{Moussa}, \textsc{Aïcha}\} pour un binôme revient exactement au même que choisir \{\textsc{Aïcha}, \textsc{Moussa}\}. C'est ici qu'interviennent les \textbf{combinaisons}.

\courssub{De l'arrangement à la combinaison : l'ordre ne compte plus}

\paragraph{Intuition.}
Reprenons l'exemple des arrangements de $2$ éléments parmi $3$ : $\{a,b,c\}$.
Les arrangements $A(3,2)$ sont les $6$ couples ordonnés :
\[(a,b),\ (a,c),\ (b,a),\ (b,c),\ (c,a),\ (c,b).\]
Si l'ordre ne compte plus, les couples $(a,b)$ et $(b,a)$ représentent \textbf{une seule et même équipe} $\{a,b\}$. De même $(a,c)\sim(c,a)$ et $(b,c)\sim(c,b)$.
Les $6$ arrangements se regroupent donc en $3$ groupes de $2$ : $\{a,b\},\ \{a,c\},\ \{b,c\}$.
Chaque groupe correspond à un \textbf{sous-ensemble à $2$ éléments}. Il y a exactement $p! = 2!$ façons d'ordonner les éléments de chaque sous-ensemble.
C'est l'idée fondamentale : \textbf{Arrangement = Choix du sous-ensemble $\times$ Ordonnancement de ce sous-ensemble}.

\begin{popfigure}
\begin{tikzpicture}[scale=1.1, every node/.style={font=\small}]
% Niveau 1 : Ensemble de départ
\node[draw, rounded corners, fill=popblueL, thick, align=center] (E) at (0,3) {Ensemble $E = \{a,b,c\}$\\$n=3$};

% Niveau 2 : Choix des sous-ensembles (Combinaisons C(3,2))
\node[draw, rounded corners, fill=popgreenL, thick] (C1) at (-3,1) {$\{a,b\}$};
\node[draw, rounded corners, fill=popgreenL, thick] (C2) at (0,1) {$\{a,c\}$};
\node[draw, rounded corners, fill=popgreenL, thick] (C3) at (3,1) {$\{b,c\}$};

\draw[->, thick, popdark] (E.south) -- ++(0,-0.3) -| (C1.north) node[midway, left, popdark] {\footnotesize Choix 1};
\draw[->, thick, popdark] (E.south) -- (C2.north) node[midway, left, popdark] {\footnotesize Choix 2};
\draw[->, thick, popdark] (E.south) -- ++(0,-0.3) -| (C3.north) node[midway, right, popdark] {\footnotesize Choix 3};

% Niveau 3 : Ordonnancement (p! = 2!)
\node[draw, rounded corners, fill=poporangeL, thick] (A11) at (-4,-1) {$(a,b)$};
\node[draw, rounded corners, fill=poporangeL, thick] (A12) at (-2,-1) {$(b,a)$};
\node[draw, rounded corners, fill=poporangeL, thick] (A21) at (-1,-1) {$(a,c)$};
\node[draw, rounded corners, fill=poporangeL, thick] (A22) at (1,-1) {$(c,a)$};
\node[draw, rounded corners, fill=poporangeL, thick] (A31) at (2,-1) {$(b,c)$};
\node[draw, rounded corners, fill=poporangeL, thick] (A32) at (4,-1) {$(c,b)$};

\draw[->, poporange] (C1.south) -- (A11.north);
\draw[->, poporange] (C1.south) -- (A12.north);
\draw[->, poporange] (C2.south) -- (A21.north);
\draw[->, poporange] (C2.south) -- (A22.north);
\draw[->, poporange] (C3.south) -- (A31.north);
\draw[->, poporange] (C3.south) -- (A32.north);

% Annotations
\node[align=center, popgreen, font=\bfseries\small] at (-3,1.7) {$C(3,2)=3$\\sous-ensembles};
\node[align=center, poporange, font=\bfseries\small] at (0,-1.7) {$A(3,2)=6$\\arrangements};
\node[align=center, popblue, font=\bfseries\small] at (3,1.7) {Chaque sous-ensemble\\donne $p!=2$ arrangements};

% Relation finale (remplace l'accolade)
\draw[thick, popdark] (-4.3,-1.5) -- (4.3,-1.5);
\node[popdark, font=\bfseries\small] at (0,-2.0) {$A(n,p) = C(n,p) \times p!$};
\end{tikzpicture}
\legende{Décomposition d'un arrangement en choix d'un sous-ensemble (combinaison) puis ordonnancement (permutation).}
\end{popfigure}

\courssub{Définition et formule des combinaisons}

\begin{definition}[Combinaison de $p$ éléments parmi $n$]
Soit $E$ un ensemble fini de cardinal $n$ ($n \in \mathbb{N}^*$) et $p$ un entier tel que $0 \le p \le n$.
Une \textbf{combinaison} de $p$ éléments parmi $n$ est un \textbf{sous-ensemble} de $E$ à $p$ éléments.
L'ordre des éléments dans ce sous-ensemble \textbf{ne compte pas}.
On note $C_n^p$ (ou $\binom{n}{p}$ ou $C(n,p)$) le nombre de telles combinaisons.
\end{definition}

\begin{aretenir}[Notation]
Au Probatoire, les notations $C_n^p$ et $\binom{n}{p}$ sont toutes deux acceptées. Nous utiliserons indifféremment $C_n^p$ et $C(n,p)$ dans ce cours. La notation $\binom{n}{p}$ (lue « $n$ choisir $p$ ») est la standard international et apparaîtra naturellement avec le triangle de Pascal.
\end{aretenir}

\begin{propriete}[Formule fondamentale : $C_n^p = \dfrac{A_n^p}{p!} = \dfrac{n!}{p!\,(n-p)!}$]
\textbf{Démonstration (exigible au Probatoire).}
On compte le nombre d'arrangements $A_n^p$ de deux façons différentes (double dénombrement) :
\begin{enumerate}
    \item \textbf{Directement :} $A_n^p = n(n-1)\dots(n-p+1) = \dfrac{n!}{(n-p)!}$.
    \item \textbf{En deux étapes :}
    \begin{itemize}
        \item Étape 1 : Choisir le sous-ensemble de $p$ éléments qui va former l'arrangement. Il y a $C_n^p$ façons de faire.
        \item Étape 2 : Ordonner ces $p$ éléments choisis. Il y a $p!$ permutations possibles.
    \end{itemize}
    D'après le principe multiplicatif : $A_n^p = C_n^p \times p!$.
\end{enumerate}
En isolant $C_n^p$, on obtient :
\[C_n^p = \frac{A_n^p}{p!} = \frac{n!}{p!\,(n-p)!}.\]
\end{propriete}

\begin{attention}[Piège classique : Confondre $A_n^p$ et $C_n^p$]
\textbf{La question clé : « L'ordre compte-t-il ? »}
\begin{itemize}
    \item \textbf{OUI} (rang, code, mot de passe, podium, tirage \emph{séquentiel} sans remise) $\rightarrow$ \textbf{Arrangement} $A_n^p$.
    \item \textbf{NON} (équipe, comité, main de cartes, tirage \emph{simultané}, numéros loto) $\rightarrow$ \textbf{Combinaison} $C_n^p$.
\end{itemize}
\emph{Erreur fréquente :} Calculer $A_8^2 = 56$ pour le nombre de matchs entre 8 équipes. Comme l'ordre ne compte pas (Match A vs B = Match B vs A), c'est $C_8^2 = 28$.
\end{attention}

\begin{exemplebox}[Exemple : Tournoi inter-classes au Lycée de Nkolbisson]
Huit classes de Première (1ère C, 1ère D, 1ère E, 1ère TI, 1ère A, 1ère B, 1ère F4, 1ère F3) organisent un tournoi de football. Chaque classe rencontre chaque autre classe \textbf{exactement une fois}. Combien de matchs au total ?
\begin{itemize}
    \item Un match oppose \textbf{deux} classes distinctes. L'ordre n'a pas d'importance (1ère C vs 1ère D est le même match que 1ère D vs 1ère C).
    \item C'est un choix de $2$ classes parmi $8$ : $p=2$, $n=8$.
    \item Nombre de matchs $= C_8^2 = \dfrac{8!}{2!\,6!} = \dfrac{8 \times 7}{2} = 28$.
\end{itemize}
\textbf{Réponse :} Il y aura $28$ matchs au total.
\end{exemplebox}

\begin{exoresolu}[Comité de gestion d'une coopérative scolaire]
La coopérative du Lycée Technique de Douala-Bassa compte $12$ membres actifs. On doit former un comité de gestion composé de $5$ personnes : un président, un secrétaire, un trésorier et deux assesseurs (sans distinction entre les deux assesseurs).
Combien de comités différents peut-on former ?
\tcblower
\textbf{Analyse :} Il y a une nuance : les rôles de Président, Secrétaire, Trésorier sont \textbf{distincts} (ordonnés), alors que les deux assesseurs sont \textbf{indistincts} (non ordonnés).
\textbf{Méthode 1 (Choix séquentiel) :}
\begin{enumerate}
    \item Choisir le Président : $12$ choix.
    \item Choisir le Secrétaire : $11$ choix restants.
    \item Choisir le Trésorier : $10$ choix restants.
    \item Choisir les $2$ assesseurs parmi les $9$ restants : l'ordre ne compte pas $\rightarrow C_9^2 = \frac{9\times 8}{2} = 36$.
\end{enumerate}
Total : $12 \times 11 \times 10 \times 36 = 47\,520$ comités.

\textbf{Méthode 2 (Choix global + attribution des rôles) :}
\begin{enumerate}
    \item Choisir les $5$ membres du comité parmi les $12$ : $C_{12}^5 = \frac{12!}{5!7!} = 792$ façons.
    \item Parmi ces $5$ choisis, attribuer les 3 postes distincts (P, S, T) : c'est un arrangement de $3$ parmi $5$ $\rightarrow A_5^3 = 5 \times 4 \times 3 = 60$ façons. Les $2$ restants seront assesseurs.
\end{enumerate}
Total : $792 \times 60 = 47\,520$ comités. \textit{(Même résultat, vérification croisée réussie.)}
\end{exoresolu}

\courssub{Propriétés fondamentales des coefficients binomiaux}

Les nombres $C_n^p$ sont appelés \textbf{coefficients binomiaux}. Ils possèdent des propriétés symétriques et récursives essentielles.

\begin{propriete}[Valeurs particulières et Symétrie]
Pour tout $n \in \mathbb{N}^*$ et $p \in \llbracket 0, n \rrbracket$ :
\begin{enumerate}
    \item $C_n^0 = C_n^n = 1$ \quad (Choisir personne ou choisir tout le monde : une seule façon).
    \item $C_n^1 = C_n^{n-1} = n$ \quad (Choisir 1 élément parmi $n$, ou en exclure 1).
    \item \textbf{Symétrie :} $C_n^p = C_n^{n-p}$.
\end{enumerate}
\textbf{Preuve de la symétrie (combinatoire) :} Choisir un sous-ensemble de $p$ éléments pour le \textbf{garder} revient exactement à choisir le sous-ensemble complémentaire de $n-p$ éléments pour le \textbf{laisser}. Il y a donc une bijection entre les parties à $p$ éléments et les parties à $n-p$ éléments.
\textbf{Preuve (algébrique) :} $C_n^{n-p} = \dfrac{n!}{(n-p)!\,(n-(n-p))!} = \dfrac{n!}{(n-p)!\,p!} = C_n^p$.
\end{propriete}

\begin{propriete}[Relation de Pascal (Récurrence)]
Pour tout $n \ge 1$ et $1 \le p \le n-1$ :
\[C_n^p = C_{n-1}^{p-1} + C_{n-1}^p.\]
\end{propriete}

\begin{experience}[Démonstration guidée de la relation de Pascal]
\textbf{Objectif :} Prouver $C_n^p = C_{n-1}^{p-1} + C_{n-1}^p$ par un raisonnement combinatoire (double dénombrement).
\textbf{Énoncé :} Soit $E$ un ensemble à $n$ éléments. Fixons un élément particulier $x \in E$. On veut compter les parties de $E$ à $p$ éléments.
\begin{enumerate}
    \item \textbf{Partitionnez} l'ensemble des parties à $p$ éléments en deux catégories disjointes :
    \begin{itemize}
        \item Catégorie 1 : Les parties qui \textbf{contiennent} $x$.
        \item Catégorie 2 : Les parties qui \textbf{ne contiennent pas} $x$.
    \end{itemize}
    \item \textbf{Comptez la Catégorie 1 :} Une partie contenant $x$ et ayant $p$ éléments au total est formée de $x$ \textbf{plus} un choix de $p-1$ éléments parmi les $n-1$ éléments restants de $E \setminus \{x\}$.
    Nombre : $C_{n-1}^{p-1}$.
    \item \textbf{Comptez la Catégorie 2 :} Une partie ne contenant pas $x$ est simplement un choix de $p$ éléments parmi les $n-1$ éléments de $E \setminus \{x\}$.
    Nombre : $C_{n-1}^p$.
    \item \textbf{Concluez :} Par le principe additif (somme sur une partition), le total est la somme des deux.
\end{enumerate}
\textbf{Application numérique :} $C_5^3 = C_4^2 + C_4^3 = 6 + 4 = 10$. Vérifiez : $C_5^3 = \frac{5\times4\times3}{3\times2\times1} = 10$.
\end{experience}

\courssub{Le nombre de parties d'un ensemble fini}

\begin{propriete}[Cardinal de l'ensemble des parties]
Soit $E$ un ensemble fini de cardinal $n$. L'ensemble des parties de $E$, noté $\mathcal{P}(E)$, a pour cardinal :
\[\text{Card}(\mathcal{P}(E)) = 2^n.\]
\end{propriete}

\begin{experience}[Deux démonstrations du résultat $2^n$]
\textbf{Méthode 1 : Par la somme des combinaisons (Somme ligne du triangle de Pascal).}
Une partie de $E$ a une taille $k$ comprise entre $0$ et $n$. Pour chaque $k$, il y a $C_n^k$ parties de taille $k$. Ces ensembles de parties forment une partition de $\mathcal{P}(E)$.
\[\text{Card}(\mathcal{P}(E)) = \sum_{k=0}^n C_n^k.\]
On verra en Section 5 que cette somme vaut $2^n$ (formule du binôme avec $a=b=1$).

\textbf{Méthode 2 : Par arbre binaire (Construction constructive).}
On numérote les éléments de $E = \{x_1, x_2, \dots, x_n\}$.
Pour construire une partie $A \subseteq E$, on traite les éléments un par un. Pour chaque $x_i$, on a \textbf{deux choix exclusifs} : « $x_i \in A$ » (branche gauche) ou « $x_i \notin A$ » (branche droite).
C'est un arbre de choix à $n$ niveaux, chaque nœud ayant 2 fils.
Nombre de feuilles (issues finales = parties) $= \underbrace{2 \times 2 \times \dots \times 2}_{n \text{ fois}} = 2^n$.
\end{experience}

\begin{popfigure}
\begin{tikzpicture}[
    level distance=1.8cm,
    sibling distance=2.5cm,
    edge from parent/.style={draw, thick, ->, popdark},
    every node/.style={font=\small, draw, rounded corners, inner sep=2pt},
    level 1/.style={sibling distance=3.5cm},
    level 2/.style={sibling distance=1.8cm},
    level 3/.style={sibling distance=0.9cm}
]
% Racine
\node[fill=popblueL, thick] (root) {$E=\{a,b,c\}$}
child { node[fill=popgreenL] {$a \in A$}
    child { node[fill=poporangeL] {$b \in A$}
        child { node[fill=poppinkL] {$c \in A$} edge from parent node[left,font=\tiny] {$c\in$} }
        child { node[fill=poppinkL] {$c \notin A$} edge from parent node[right,font=\tiny] {$c\notin$} }
        edge from parent node[left,font=\tiny] {$b\in$}
    }
    child { node[fill=poporangeL] {$b \notin A$}
        child { node[fill=poppinkL] {$c \in A$} edge from parent node[left,font=\tiny] {$c\in$} }
        child { node[fill=poppinkL] {$c \notin A$} edge from parent node[right,font=\tiny] {$c\notin$} }
        edge from parent node[right,font=\tiny] {$b\notin$}
    }
    edge from parent node[left,font=\small] {$a\in$}
}
child { node[fill=popgreenL] {$a \notin A$}
    child { node[fill=poporangeL] {$b \in A$}
        child { node[fill=poppinkL] {$c \in A$} edge from parent node[left,font=\tiny] {$c\in$} }
        child { node[fill=poppinkL] {$c \notin A$} edge from parent node[right,font=\tiny] {$c\notin$} }
        edge from parent node[left,font=\tiny] {$b\in$}
    }
    child { node[fill=poporangeL] {$b \notin A$}
        child { node[fill=poppinkL] {$c \in A$} edge from parent node[left,font=\tiny] {$c\in$} }
        child { node[fill=poppinkL] {$c \notin A$} edge from parent node[right,font=\tiny] {$c\notin$} }
        edge from parent node[right,font=\tiny] {$b\notin$}
    }
    edge from parent node[right,font=\small] {$a\notin$}
};

% Étiquettes des feuilles (parties)
\node[anchor=west, font=\tiny, popdark] at (4.5, 3.6) {$\{a,b,c\}$};
\node[anchor=west, font=\tiny, popdark] at (4.5, 2.0) {$\{a,b\}$};
\node[anchor=west, font=\tiny, popdark] at (4.5, 0.4) {$\{a,c\}$};
\node[anchor=west, font=\tiny, popdark] at (4.5, -1.2) {$\{a\}$};
\node[anchor=west, font=\tiny, popdark] at (4.5, -2.8) {$\{b,c\}$};
\node[anchor=west, font=\tiny, popdark] at (4.5, -4.4) {$\{b\}$};
\node[anchor=west, font=\tiny, popdark] at (4.5, -6.0) {$\{c\}$};
\node[anchor=west, font=\tiny, popdark] at (4.5, -7.6) {$\varnothing$};

% Légende niveaux
\node[font=\bfseries\small, popblue] at (-4.5, 3) {Niveau 0};
\node[font=\bfseries\small, popgreen] at (-4.5, 1.2) {Niveau 1 ($a$)};
\node[font=\bfseries\small, poporange] at (-4.5, -0.6) {Niveau 2 ($b$)};
\node[font=\bfseries\small, poppink] at (-4.5, -2.4) {Niveau 3 ($c$)};
\node[font=\bfseries\small, popdark] at (-4.5, -4.5) {$2^3 = 8$ parties};
\end{tikzpicture}
\legende{Arbre binaire des parties de $\{a,b,c\}$. Chaque chemin de la racine à une feuille définit une partie unique.}
\end{popfigure}

\begin{exemplebox}[Application : Configuration d'un réseau local (LAN)]
Dans une salle informatique du Lycée de Maroua, $5$ postes de travail peuvent être connectés ou non à un serveur d'impression central. Chaque poste a deux états : \texttt{Connecté} ou \texttt{Déconnecté}.
Combien de configurations distinctes du réseau sont possibles ?
\textbf{Résolution :} Une configuration correspond à un sous-ensemble des $5$ postes connectés. C'est une partie de l'ensemble des $5$ postes.
Nombre de configurations $= 2^5 = 32$.
(On vérifie : $C_5^0 + C_5^1 + C_5^2 + C_5^3 + C_5^4 + C_5^5 = 1+5+10+10+5+1 = 32$).
\end{exemplebox}

\courssub{Bilan méthodologique : Choisir le bon outil}

Face à un problème de dénombrement, suivez cet algorithme décisionnel pour ne plus jamais vous tromper.

\begin{methode}[Arbre de décision pour le dénombrement (Sans répétition)]
\textbf{Question 1 :} Y a-t-il \textbf{répétition} autorisée des éléments choisis ?
\begin{itemize}
    \item \textbf{OUI} $\rightarrow$ $p$-uplets (listes) avec répétition : $\mathbf{n^p}$.
    \item \textbf{NON} (éléments distincts) $\rightarrow$ Aller à Question 2.
\end{itemize}
\textbf{Question 2 :} L'\textbf{ordre} de sélection/arrangement compte-t-il ?
\begin{itemize}
    \item \textbf{OUI} (séquence, rang, code, podium, tirage séquentiel) $\rightarrow$ \textbf{Arrangement} : $\mathbf{A_n^p = \frac{n!}{(n-p)!}}$.
    \item \textbf{NON} (équipe, comité, main, tirage simultané, numéros loto) $\rightarrow$ \textbf{Combinaison} : $\mathbf{C_n^p = \frac{n!}{p!(n-p)!}}$.
\end{itemize}
\textbf{Cas particulier :} Si $p = n$ (on utilise tous les éléments) et pas de répétition :
\begin{itemize}
    \item Ordre compte $\rightarrow$ \textbf{Permutation} : $\mathbf{n!}$.
    \item Ordre ne compte pas $\rightarrow$ Une seule partie (l'ensemble lui-même) : $\mathbf{1}$.
\end{itemize}
\end{methode}

\begin{center}
\begin{tabularx}{\linewidth}{|l|X|X|X|}\hline
\rowcolor{popblueL}
\textbf{Situation type} & \textbf{Ordre ?} & \textbf{Répétition ?} & \textbf{Formule} \\ \hline
Code PIN, Mot de passe, Numéro de téléphone & OUI & OUI & $n^p$ \\ \hline
Podium (1er, 2e, 3e), Rangement de livres, Tirage séquentiel sans remise & OUI & NON & $A_n^p = \frac{n!}{(n-p)!}$ \\ \hline
Permutation (anagramme, classement complet) & OUI & NON ($p=n$) & $n!$ \\ \hline
Équipe, Comité, Main de cartes, Loto, Matchs (A vs B), Parties d'ensemble & NON & NON & $C_n^p = \frac{n!}{p!(n-p)!}$ \\ \hline
Nombre total de sous-ensembles (toutes tailles) & NON & NON & $2^n$ \\ \hline
\end{tabularx}
\end{center}

\begin{attention}[Erreurs fatales à éviter au Probatoire]
\begin{enumerate}
    \item \textbf{Oublier la condition $0 \le p \le n$ pour $C_n^p$ et $A_n^p$.} Si $p > n$, le nombre est $0$ (impossible de choisir plus d'éléments qu'il n'y en a).
    \item \textbf{Confondre « tirage simultané » et « tirage successif sans remise ».}
    \begin{itemize}
        \item « On tire 3 boules \textbf{simultanément} » $\rightarrow$ Combinaison $C_n^3$.
        \item « On tire 3 boules \textbf{l'une après l'autre sans remise} » $\rightarrow$ Si on note l'ordre de sortie : Arrangement $A_n^3$. Si on ne note que le contenu final : Combinaison $C_n^3$. \textbf{Lisez l'énoncé : « Combien de résultats différents ? » (souvent = combinaisons) vs « Combien de façons de tirer ? » (parfois = arrangements).}
    \end{itemize}
    \item \textbf{Calculer $C_n^p$ sans simplifier les factorielles.} Toujours simplifier : $C_{100}^2 = \frac{100 \times 99}{2} = 4950$, pas $\frac{100!}{2!98!}$.
    \item \textbf{Oublier les cas $p=0$ ou $p=n$.} $C_n^0 = 1$ (la partie vide) et $C_n^n = 1$ (l'ensemble tout entier). Cela compte dans la somme $2^n$.
\end{enumerate}
\end{attention}

\begin{savaistu}[Le Loto Camerounais (Lonaci / PMUC) et les combinaisons]
Au Cameroun, le jeu de loto classique « 5/90 » (ou variantes 5/36, 6/49 selon les opérateurs comme la LONACI) consiste à cocher $5$ numéros sur une grille de $90$.
Le nombre de grilles possibles (combinaisons simples) est $C_{90}^5 = \frac{90 \times 89 \times 88 \times 87 \times 86}{5 \times 4 \times 3 \times 2 \times 1} = 43\,949\,268$.
C'est pourquoi gagner au jackpot est un événement rare (probabilité $\approx 2,27 \times 10^{-8}$) !
Les « systèmes » (cocher 6, 7, 8... numéros) correspondent à jouer toutes les combinaisons de $5$ parmi les $k$ numéros cochés : $C_k^5$ grilles. Le prix de la mise est proportionnel à $C_k^5$.
\end{savaistu}

\begin{exercice}[Entraînement : Le bon outil]
Pour chaque situation, identifiez l'outil (Puissance $n^p$, Arrangement $A_n^p$, Combinaison $C_n^p$, Permutation $n!$, Parties $2^n$) et donnez l'expression littérale (ne pas calculer numériquement).
\begin{enumerate}
    \item Un code de carte SIM MTN à 4 chiffres (chiffres de 0 à 9).
    \item Un jury de 3 professeurs à choisir parmi 10 collègues.
    \item Classement des 5 premières places d'un marathon avec 50 coureurs.
    \item Nombre de diagonales d'un polygone convexe à 12 sommets.
    \item Un élève doit répondre à 7 questions sur 10 proposées à l'examen.
    \item Nombre de sous-ensembles d'un ensemble à 8 éléments qui contiennent au moins 1 élément.
\end{enumerate}
\end{exercice}

\begin{corrige}[Exercice : Le bon outil]
\begin{enumerate}
    \item \textbf{Puissance $10^4$}. Ordre OUI (code), Répétition OUI (chiffres peuvent se répéter). $n=10, p=4$.
    \item \textbf{Combinaison $C_{10}^3$}. Ordre NON (équipe/jury), Répétition NON. $n=10, p=3$.
    \item \textbf{Arrangement $A_{50}^5$}. Ordre OUI (classement 1er, 2e...), Répétition NON. $n=50, p=5$.
    \item \textbf{Combinaison $C_{12}^2 - 12$}. Une diagonale relie 2 sommets non consécutifs. Nombre de paires de sommets = $C_{12}^2$. On retire les 12 côtés. Résultat : $C_{12}^2 - 12 = 66 - 12 = 54$.
    \item \textbf{Combinaison $C_{10}^7$}. Ordre NON (choix des questions), Répétition NON. $n=10, p=7$. Note : $C_{10}^7 = C_{10}^3 = 120$.
    \item \textbf{Parties non vides : $2^8 - 1$}. Total parties $2^8$. On retire la partie vide ($\varnothing$). $256 - 1 = 255$.
\end{enumerate}
\end{corrige}

\begin{aretenir}[Synthèse de la section : Combinaisons et Parties]
\begin{itemize}
    \item \textbf{Combinaison $C_n^p$ :} Choix de $p$ éléments parmi $n$, \textbf{sans ordre}, \textbf{sans répétition}. Formule : $\dfrac{n!}{p!(n-p)!}$.
    \item \textbf{Lien Arrangement/Combinaison :} $A_n^p = C_n^p \times p!$. L'arrangement surcompte les combinaisons par un facteur $p!$ (les permutations internes).
    \item \textbf{Propriétés clés :} $C_n^0=C_n^n=1$, $C_n^1=n$, \textbf{Symétrie} $C_n^p=C_n^{n-p}$, \textbf{Relation de Pascal} $C_n^p=C_{n-1}^{p-1}+C_{n-1}^p$.
    \item \textbf{Parties d'un ensemble :} $\text{Card}(\mathcal{P}(E)) = 2^n = \sum_{k=0}^n C_n^k$. Démonstration par arbre binaire (choix binaire pour chaque élément) ou par somme des combinaisons.
    \item \textbf{Mots-clés pour repérer une combinaison :} « Équipe », « Comité », « Groupe », « Sous-ensemble », « Tirage simultané », « Main », « Loto », « Matchs », « Diagonales », « Poignées de main ».
\end{itemize}
\end{aretenir}

\coursec{Triangle de Pascal et formule du binôme}

Dans la section précédente, nous avons appris à compter les parties d'un ensemble : le nombre de façons de choisir $p$ éléments parmi $n$ est le coefficient binomial $\binom{n}{p} = \dfrac{n!}{p!(n-p)!}$. Nous avons aussi rencontré la \cle{relation de Pascal} :
\[ \binom{n}{p} = \binom{n-1}{p-1} + \binom{n-1}{p}. \]
Cette petite formule, en apparence anodine, est une véritable machine à calculer : elle permet de fabriquer \emph{tous} les coefficients binomiaux, de proche en proche, par de simples additions. C'est le fameux \cle{triangle de Pascal}. Et surprise magnifique : ces nombres, nés d'un problème de dénombrement, apparaissent exactement quand on développe $(a+b)^n$. C'est l'objet de cette section.

\courssub{Construction du triangle de Pascal}

\textbf{Intuition.} Imagine que tu dois calculer $\binom{7}{3}$. La formule avec les factorielles fonctionne, mais elle est lourde. La relation de Pascal dit autre chose : pour choisir 3 élèves parmi 7, distingue un élève particulier, disons Amina. Ou bien Amina est choisie (il reste 2 élèves à choisir parmi les 6 autres : $\binom{6}{2}$ façons), ou bien elle ne l'est pas (il faut choisir les 3 parmi les 6 autres : $\binom{6}{3}$ façons). Donc $\binom{7}{3} = \binom{6}{2} + \binom{6}{3} = 15 + 20 = 35$. Chaque coefficient est la \emph{somme des deux coefficients situés juste au-dessus de lui}.

\begin{propriete}[Relation de Pascal (rappel)]
Pour tous entiers $n \geq 1$ et $1 \leq p \leq n-1$ :
\[ \binom{n}{p} = \binom{n-1}{p-1} + \binom{n-1}{p}, \qquad \text{avec} \quad \binom{n}{0} = \binom{n}{n} = 1. \]
\end{propriete}

\begin{definition}[Triangle de Pascal]
Le \cle{triangle de Pascal} est le tableau triangulaire dont la ligne numérotée $n$ contient, de gauche à droite, les coefficients $\binom{n}{0}, \binom{n}{1}, \ldots, \binom{n}{n}$. On le construit ligne par ligne :
\begin{itemize}
\item on place des $1$ sur les deux bords (car $\binom{n}{0} = \binom{n}{n} = 1$) ;
\item chaque nombre intérieur est la somme des deux nombres situés au-dessus de lui.
\end{itemize}
\end{definition}

\begin{popfigure}
\begin{center}
\begin{tikzpicture}[scale=0.92, every node/.style={font=\small}]
% Lignes du triangle : n = 0 a 6
\foreach \n/\row in {0/{1}, 1/{1,1}, 2/{1,2,1}, 3/{1,3,3,1}, 4/{1,4,6,4,1}, 5/{1,5,10,10,5,1}, 6/{1,6,15,20,15,6,1}} {
  \foreach [count=\k from 0] \val in \row {
    \ifnum\n=4
      \node[draw=poporange, fill=poporangeL, circle, inner sep=1.5pt, minimum size=16pt] (n\n-\k) at (\k-\n/2, -\n) {\val};
    \else
      \node[fill=popblueL, circle, inner sep=1.5pt, minimum size=16pt] (n\n-\k) at (\k-\n/2, -\n) {\val};
    \fi
  }
}
% Etiquettes des lignes
\foreach \n in {0,1,2,3,4,5,6} {
  \node[left, color=popmuted] at (-\n/2-0.9, -\n) {$n=\n$};
}
% Fleches illustrant la relation de Pascal : 10 = 4 + 6 (ligne 5, k=1)
\draw[-{Stealth}, thick, popgreen] (n4-0.south west) -- (n5-1.north west);
\draw[-{Stealth}, thick, popgreen] (n4-1.south) -- (n5-1.north east);
\node[color=popgreen, font=\small\bfseries] at (3.4,-4.6) {$10 = 4 + 6$};
% Fleches illustrant 20 = 10 + 10
\draw[-{Stealth}, thick, poppurple] (n5-2.south) -- (n6-3.north west);
\draw[-{Stealth}, thick, poppurple] (n5-3.south west) -- (n6-3.north east);
\node[color=poppurple, font=\small\bfseries] at (3.4,-5.7) {$20 = 10 + 10$};
\end{tikzpicture}
\end{center}
\legende{Triangle de Pascal jusqu'à la ligne $n=6$. Chaque nombre est la somme des deux nombres au-dessus de lui (flèches). La ligne $n=4$ est mise en évidence : ce sont les coefficients de $(a+b)^4$.}
\end{popfigure}

\begin{exemplebox}[Lecture du triangle]
Pour lire $\binom{5}{2}$, on se place sur la ligne $n=5$ et on prend le 3\textsuperscript{e} nombre (car on commence à $k=0$) : la ligne est $1 \quad 5 \quad 10 \quad 10 \quad 5 \quad 1$, donc $\binom{5}{2} = 10$. Vérification : $\binom{5}{2} = \dfrac{5!}{2!\,3!} = \dfrac{120}{2 \times 6} = 10$.
\end{exemplebox}

\begin{attention}[Bien numéroter à partir de zéro]
La ligne « $1 \quad 4 \quad 6 \quad 4 \quad 1$ » est la ligne $n=4$, pas la ligne $n=5$ ! La toute première ligne (le $1$ solitaire au sommet) est la ligne $n=0$. De même, dans une ligne, les positions sont numérotées $k=0, 1, 2, \ldots$ : sur la ligne $n=4$, le nombre $6$ correspond à $\binom{4}{2}$. Une erreur de décalage fausse tout le développement.
\end{attention}

\courssub{Formule du binôme de Newton}

\textbf{Intuition.} Développons à la main $(a+b)^3 = (a+b)(a+b)(a+b)$. Pour former chaque terme, on choisit dans \emph{chacune} des 3 parenthèses soit $a$, soit $b$. Obtenir $a^2 b$, c'est choisir $b$ dans exactement \emph{une} parenthèse parmi 3 (et $a$ dans les deux autres) : il y a $\binom{3}{1} = 3$ façons, d'où le terme $3a^2b$. De même, $ab^2$ apparaît $\binom{3}{2} = 3$ fois. On retrouve :
\[ (a+b)^3 = a^3 + 3a^2b + 3ab^2 + b^3. \]
Les coefficients sont exactement la ligne $n=3$ du triangle ! Ce n'est pas un hasard : c'est le théorème central de cette leçon.

\begin{propriete}[Formule du binôme de Newton — admise]
Pour tous nombres $a$ et $b$ et tout entier naturel $n \geq 1$ :
\[ (a+b)^n = \sum_{k=0}^{n} \binom{n}{k}\, a^{\,n-k}\, b^{k} = \binom{n}{0}a^n + \binom{n}{1}a^{n-1}b + \cdots + \binom{n}{k}a^{n-k}b^k + \cdots + \binom{n}{n}b^n. \]
Les coefficients $\binom{n}{k}$ sont précisément ceux de la ligne $n$ du triangle de Pascal. \emph{Cette formule est admise au niveau Première ; sa démonstration par récurrence est proposée ci-dessous en démonstration guidée pour les élèves rapides.}
\end{propriete}

\begin{experience}[Démonstration guidée par récurrence (pour aller plus loin)]
On veut prouver que pour tout $n \geq 1$, $(a+b)^n = \sum_{k=0}^{n}\binom{n}{k}a^{n-k}b^k$.
\begin{enumerate}
\item \textbf{Initialisation.} Pour $n=1$ : $(a+b)^1 = a+b$ et $\binom{1}{0}a + \binom{1}{1}b = a+b$. La formule est vraie.
\item \textbf{Hérédité.} Suppose la formule vraie au rang $n$. Alors $(a+b)^{n+1} = (a+b)(a+b)^n$. En développant avec l'hypothèse de récurrence :
\[ (a+b)^{n+1} = a\sum_{k=0}^{n}\binom{n}{k}a^{n-k}b^k + b\sum_{k=0}^{n}\binom{n}{k}a^{n-k}b^k. \]
Le premier morceau donne les termes en $a^{n+1-k}b^k$ avec coefficient $\binom{n}{k}$ ; le second, après le changement d'indice $k \to k-1$, donne les termes en $a^{n+1-k}b^k$ avec coefficient $\binom{n}{k-1}$.
\item \textbf{Conclusion de l'hérédité.} Le coefficient de $a^{n+1-k}b^k$ est donc $\binom{n}{k} + \binom{n}{k-1} = \binom{n+1}{k}$ par la relation de Pascal. Les termes extrêmes $a^{n+1}$ et $b^{n+1}$ ont pour coefficient $1 = \binom{n+1}{0} = \binom{n+1}{n+1}$. La formule est donc vraie au rang $n+1$.
\item \textbf{Conclusion.} Par récurrence, la formule est vraie pour tout $n \geq 1$.
\end{enumerate}
Retiens l'idée clé : \emph{la relation de Pascal est exactement ce qui fait fonctionner la récurrence}. Triangle et binôme sont les deux visages d'un même objet.
\end{experience}

\begin{methode}[Développer $(a+b)^n$ en pratique]
\begin{enumerate}
\item Écris la \textbf{ligne $n$} du triangle de Pascal (par exemple, pour $n=4$ : $1, 4, 6, 4, 1$).
\item Écris les \textbf{puissances de $a$ en décroissant} : $a^n, a^{n-1}, \ldots, a, 1$.
\item Écris les \textbf{puissances de $b$ en croissant} : $1, b, b^2, \ldots, b^n$.
\item Multiplie colonne par colonne et additionne : chaque terme est $\binom{n}{k} a^{n-k} b^k$.
\item Vérifie : il doit y avoir $n+1$ termes, et dans chaque terme la somme des exposants vaut $n$.
\end{enumerate}
\end{methode}

\begin{exemplebox}[Les développements à connaître]
\begin{align*}
(a+b)^2 &= a^2 + 2ab + b^2 \\
(a+b)^3 &= a^3 + 3a^2b + 3ab^2 + b^3 \\
(a+b)^4 &= a^4 + 4a^3b + 6a^2b^2 + 4ab^3 + b^4 \\
(a+b)^5 &= a^5 + 5a^4b + 10a^3b^2 + 10a^2b^3 + 5ab^4 + b^5
\end{align*}
Observe la symétrie des coefficients ($\binom{n}{k} = \binom{n}{n-k}$) : elle se lit directement sur le triangle.
\end{exemplebox}

\begin{popfigure}
\begin{center}
\begin{tikzpicture}[scale=1]
% Ligne n=4 du triangle
\foreach [count=\k from 0] \val in {1,4,6,4,1} {
  \node[draw=popblue, fill=popblueL, circle, minimum size=20pt, font=\bfseries] (c\k) at (1.6*\k, 0) {\val};
}
% Puissances de a
\foreach [count=\k from 0] \pa in {a^4, a^3, a^2, a, 1} {
  \node[color=popdark] (pa\k) at (1.6*\k, -1) {$\pa$};
}
% Puissances de b
\foreach [count=\k from 0] \pb in {1, b, b^2, b^3, b^4} {
  \node[color=popdark] (pb\k) at (1.6*\k, -1.8) {$\pb$};
}
% Fleches
\foreach \k in {0,1,2,3,4} {
  \draw[-{Stealth}, poporange, thick] (c\k.south) -- (pa\k.north);
  \draw[-{Stealth}, popgreen, thick] (pa\k.south) -- (pb\k.north);
}
\node[color=popmuted, left] at (-1.2,0) {coefficients};
\node[color=popmuted, left] at (-1.2,-1) {puissances de $a$};
\node[color=popmuted, left] at (-1.2,-1.8) {puissances de $b$};
\node[align=center, font=\small] at (3.2,-2.9) {$(a+b)^4 = a^4 + 4a^3b + 6a^2b^2 + 4ab^3 + b^4$};
\end{tikzpicture}
\end{center}
\legende{Fabriquer le développement de $(a+b)^4$ : ligne $n=4$ du triangle, puissances de $a$ décroissantes, puissances de $b$ croissantes.}
\end{popfigure}

\courssub{Applications : développements concrets et coefficients particuliers}

\begin{exoresolu}[Développement complet de $(x+2)^4$]
Développer et réduire $(x+2)^4$.
\tcblower
On applique la formule du binôme avec $a = x$, $b = 2$ et la ligne $n=4$ : $1, 4, 6, 4, 1$.
\begin{align*}
(x+2)^4 &= \binom{4}{0}x^4 \cdot 2^0 + \binom{4}{1}x^3 \cdot 2^1 + \binom{4}{2}x^2 \cdot 2^2 + \binom{4}{3}x \cdot 2^3 + \binom{4}{4} \cdot 2^4 \\
&= 1 \cdot x^4 \cdot 1 + 4 \cdot x^3 \cdot 2 + 6 \cdot x^2 \cdot 4 + 4 \cdot x \cdot 8 + 1 \cdot 16 \\
&= x^4 + 8x^3 + 24x^2 + 32x + 16.
\end{align*}
Vérification rapide : pour $x=0$, $(0+2)^4 = 16$ et le développement donne bien $16$. Pour $x=1$ : $3^4 = 81$ et $1+8+24+32+16 = 81$.
\end{exoresolu}

\begin{exoresolu}[Développement de $(x-2)^4$ et signes alternés]
Développer $(x-2)^4$.
\tcblower
On écrit $(x-2)^4 = (x + (-2))^4$ et on applique le binôme avec $b = -2$ :
\begin{align*}
(x-2)^4 &= x^4 + 4x^3(-2) + 6x^2(-2)^2 + 4x(-2)^3 + (-2)^4 \\
&= x^4 - 8x^3 + 24x^2 - 32x + 16.
\end{align*}
Les signes alternent : $+, -, +, -, +$, car les puissances \emph{impaires} de $(-2)$ sont négatives.
\end{exoresolu}

\begin{exoresolu}[Coefficient d'un terme particulier]
Quel est le coefficient de $x^5$ dans le développement de $(2x+1)^8$ ?
\tcblower
Le terme général du développement est $\binom{8}{k}(2x)^{8-k} \cdot 1^k = \binom{8}{k} 2^{8-k} x^{8-k}$.
On veut $x^5$, donc $8 - k = 5$, c'est-à-dire $k = 3$. Le coefficient cherché est :
\[ \binom{8}{3} \times 2^5 = 56 \times 32 = 1792. \]
Le terme correspondant est $1792\,x^5$. Remarque la méthode : on n'a pas eu besoin de développer entièrement, il suffit de repérer la bonne valeur de $k$.
\end{exoresolu}

\begin{attention}[Les deux pièges classiques du Probatoire]
\begin{itemize}
\item \textbf{Oublier les signes dans $(a-b)^n$.} Écris toujours $(a-b)^n = (a+(-b))^n$ : les termes contenant une puissance impaire de $b$ portent un signe moins. Ainsi $(a-b)^3 = a^3 - 3a^2b + 3ab^2 - b^3$, et non $a^3 - 3a^2b - 3ab^2 - b^3$.
\item \textbf{Oublier d'élever les coefficients aux bonnes puissances.} Dans $(2x+1)^n$, le terme $\binom{n}{k}(2x)^{n-k}$ contient $2^{\,n-k}$ : le $2$ est lui aussi élevé à la puissance $n-k$ ! Écrire $(x+2)^4 = x^4 + 4x^3 + 6x^2 + 4x + 1$ est faux : les puissances de $2$ ont disparu.
\end{itemize}
\end{attention}

\courssub{Deux conséquences remarquables}

\begin{propriete}[Somme des coefficients binomiaux]
Pour tout entier naturel $n$ :
\[ \sum_{k=0}^{n} \binom{n}{k} = \binom{n}{0} + \binom{n}{1} + \cdots + \binom{n}{n} = 2^n. \]
\end{propriete}

\begin{experience}[Démonstration]
Applique la formule du binôme avec $a = 1$ et $b = 1$ :
\[ (1+1)^n = \sum_{k=0}^{n} \binom{n}{k} \, 1^{n-k} \, 1^{k} = \sum_{k=0}^{n} \binom{n}{k}. \]
Or $(1+1)^n = 2^n$, d'où le résultat.
\end{experience}

\begin{aretenir}[Lien avec le dénombrement des parties]
Ce résultat redonne élégamment ce que nous avions établi dans la section précédente : un ensemble à $n$ éléments possède $2^n$ parties. En effet, $\binom{n}{k}$ compte les parties à $k$ éléments, et en sommant sur $k = 0, 1, \ldots, n$ on compte \emph{toutes} les parties. Par exemple, la ligne $n=4$ du triangle donne $1+4+6+4+1 = 16 = 2^4$.
\end{aretenir}

\begin{propriete}[Somme alternée]
Pour tout entier $n \geq 1$ :
\[ \sum_{k=0}^{n} (-1)^k \binom{n}{k} = \binom{n}{0} - \binom{n}{1} + \binom{n}{2} - \cdots + (-1)^n \binom{n}{n} = 0. \]
\end{propriete}

\begin{experience}[Démonstration]
Applique la formule du binôme avec $a = 1$ et $b = -1$ :
\[ (1 + (-1))^n = \sum_{k=0}^{n} \binom{n}{k} \, 1^{n-k} \, (-1)^k = \sum_{k=0}^{n} (-1)^k \binom{n}{k}. \]
Or $(1-1)^n = 0^n = 0$ pour $n \geq 1$, d'où le résultat. Vérification sur la ligne $n=4$ : $1 - 4 + 6 - 4 + 1 = 0$.
\end{experience}

\begin{exemplebox}[Interprétation combinatoire de la somme alternée]
L'égalité $\sum_{k=0}^{n} (-1)^k \binom{n}{k} = 0$ se réécrit $\binom{n}{0} + \binom{n}{2} + \cdots = \binom{n}{1} + \binom{n}{3} + \cdots$ : un ensemble à $n$ éléments ($n \geq 1$) a \emph{autant} de parties à nombre pair d'éléments que de parties à nombre impair d'éléments, à savoir $2^{n-1}$ de chaque sorte. Pour $n=3$ : parties paires $\binom{3}{0}+\binom{3}{2} = 1+3 = 4$, parties impaires $\binom{3}{1}+\binom{3}{3} = 3+1 = 4$.
\end{exemplebox}

\begin{savaistu}[Un triangle bien plus vieux que Pascal]
Blaise Pascal publie son \emph{Traité du triangle arithmétique} en 1654, mais ce triangle était connu bien plus tôt : le mathématicien persan \textbf{Al-Karaji} l'utilisait vers l'an 1000 pour développer $(a+b)^n$, et le Chinois \textbf{Yang Hui} le présentait en 1261 dans un ouvrage où il attribue sa découverte à Jia Xian (XI\textsuperscript{e} siècle). En Chine, on l'appelle d'ailleurs « triangle de Yang Hui ». Le génie de Pascal fut d'en explorer systématiquement les propriétés et de l'appliquer au calcul des probabilités — un domaine que tu rencontreras en Terminale.
\end{savaistu}

\begin{savaistu}[Vers la Terminale : le binôme et les probabilités (hors exigibilité en Première)]
Les coefficients binomiaux reviendront en force en Terminale avec la \cle{loi binomiale}. Par exemple, si tu lances 5 fois une pièce équilibrée, la probabilité d'obtenir exactement 3 « Face » est $\binom{5}{3}\left(\tfrac{1}{2}\right)^3\left(\tfrac{1}{2}\right)^2 = \dfrac{10}{32}$ : le coefficient $\binom{5}{3}$ compte les positions possibles des 3 « Face » parmi les 5 lancers. La formule du binôme garantit d'ailleurs que la somme de toutes ces probabilités vaut $\left(\tfrac12 + \tfrac12\right)^5 = 1$. Tout se tient !
\end{savaistu}

\begin{exercice}[Pour t'entraîner]
\begin{enumerate}
\item Construire la ligne $n=7$ du triangle de Pascal.
\item Développer et réduire $(x+1)^5$ puis $(2x-3)^3$.
\item Déterminer le coefficient de $x^3$ dans $(x+3)^6$.
\item Calculer sans calculatrice : $\binom{6}{0} + \binom{6}{1} + \binom{6}{2} + \binom{6}{3} + \binom{6}{4} + \binom{6}{5} + \binom{6}{6}$.
\end{enumerate}
\end{exercice}

\begin{corrige}[Exercice]
\begin{enumerate}
\item On part de la ligne $n=6$ : $1, 6, 15, 20, 15, 6, 1$. Chaque nombre de la ligne $7$ est la somme de deux voisins : $1$, $1+6=7$, $6+15=21$, $15+20=35$, $20+15=35$, $15+6=21$, $6+1=7$, $1$. Ligne $n=7$ : $1, 7, 21, 35, 35, 21, 7, 1$.
\item $(x+1)^5 = x^5 + 5x^4 + 10x^3 + 10x^2 + 5x + 1$ (ligne $n=5$, toutes les puissances de $1$ valent $1$).
Pour $(2x-3)^3$, on utilise la ligne $n=3$ avec $a=2x$ et $b=-3$ :
\begin{align*}
(2x-3)^3 &= (2x)^3 + 3(2x)^2(-3) + 3(2x)(-3)^2 + (-3)^3 \\
&= 8x^3 - 36x^2 + 54x - 27.
\end{align*}
\item Terme général : $\binom{6}{k} x^{6-k} 3^k$. On veut $6-k = 3$, donc $k=3$ : coefficient $\binom{6}{3} \times 3^3 = 20 \times 27 = 540$.
\item C'est la somme des coefficients de la ligne $n=6$, qui vaut $2^6 = 64$.
\end{enumerate}
\end{corrige}

\begin{aretenir}[L'essentiel de la section]
\begin{itemize}
\item Le \textbf{triangle de Pascal} se construit par additions successives grâce à $\binom{n}{p} = \binom{n-1}{p-1} + \binom{n-1}{p}$, avec des $1$ sur les bords.
\item \textbf{Formule du binôme} : $(a+b)^n = \displaystyle\sum_{k=0}^{n} \binom{n}{k} a^{n-k} b^k$ ; les coefficients sont la ligne $n$ du triangle.
\item Puissances de $a$ décroissantes, puissances de $b$ croissantes ; somme des exposants toujours égale à $n$ ; $n+1$ termes.
\item Dans $(a-b)^n$, les signes alternent ; dans $(2x+1)^n$, n'oublie pas les puissances de $2$.
\item En posant $a=b=1$ : $\sum_{k=0}^{n}\binom{n}{k} = 2^n$ (nombre total de parties). En posant $a=1$, $b=-1$ : $\sum_{k=0}^{n}(-1)^k\binom{n}{k} = 0$.
\end{itemize}
\end{aretenir}

\coursec{Synthèse : choisir le bon outil de dénombrement}

Dans la section précédente, nous avons construit le triangle de Pascal et la formule du binôme $(a+b)^n = \sum_{p=0}^{n} C_n^p\, a^{n-p}b^p$, qui couronne l'étude des combinaisons. Nous disposons désormais de toute la panoplie du dénombrement : $n^p$, $A_n^p$, $C_n^p$, $n!$. Mais posséder quatre outils ne suffit pas : au Probatoire comme dans la vie, la vraie compétence consiste à \cle{choisir le bon outil} face à une situation concrète. C'est l'objet de cette dernière section, qui est aussi la plus stratégique de la leçon.

\courssub{Les deux questions qui décident de tout}

Toute situation de dénombrement se ramène à un schéma de choix : on constitue un résultat en choisissant $p$ objets (ou places) parmi $n$ objets disponibles. Pour savoir quel outil utiliser, il suffit de se poser, dans l'ordre, deux questions.

\begin{definition}[Les deux critères fondamentaux]
Face à une situation de dénombrement, on se demande :
\begin{enumerate}
\item \textbf{L'\cle{ordre} compte-t-il ?} Deux résultats formés des mêmes éléments mais dans un ordre différent sont-ils considérés comme distincts ? (Podium : oui. Équipe : non.)
\item \textbf{La \cle{répétition} est-elle possible ?} Un même élément peut-il être choisi plusieurs fois ? (Chiffres d'un code : oui. Joueurs d'une équipe : non, on ne peut pas être titulaire deux fois.)
\end{enumerate}
\end{definition}

\begin{exemplebox}[Intuition : quatre situations, quatre outils]
Une urne contient 5 boules numérotées de 1 à 5 ; on en extrait 3.
\begin{itemize}
\item Tirage \emph{successif avec remise}, on note la suite obtenue : ordre oui, répétition oui $\Rightarrow$ $5^3 = 125$ résultats.
\item Tirage \emph{successif sans remise}, on note la suite : ordre oui, répétition non $\Rightarrow$ $A_5^3 = 5 \times 4 \times 3 = 60$.
\item Tirage \emph{simultané} de 3 boules : ordre non, répétition non $\Rightarrow$ $C_5^3 = \dfrac{5!}{3!\,2!} = 10$.
\item On range les 5 boules en file : ordre oui, tous les éléments $\Rightarrow$ $5! = 120$.
\end{itemize}
\end{exemplebox}

\begin{attention}[Mal lire l'énoncé : l'erreur la plus coûteuse]
La majorité des erreurs au Probatoire vient d'une lecture trop rapide. Retiens les mots-clés :
\begin{itemize}
\item \emph{« successivement avec remise »} $\Rightarrow$ répétition possible : $n^p$.
\item \emph{« successivement sans remise »} $\Rightarrow$ arrangement : $A_n^p$.
\item \emph{« simultanément »}, \emph{« en une seule fois »}, \emph{« une équipe de »}, \emph{« une délégation de »} $\Rightarrow$ combinaison : $C_n^p$.
\item \emph{« classement »}, \emph{« podium »}, \emph{« code »}, \emph{« numéro »} $\Rightarrow$ l'ordre compte.
\item \emph{« bureau »} (président, secrétaire, trésorier) : les postes sont \emph{distincts}, donc l'ordre compte $\Rightarrow$ arrangement, même si le mot « ordre » n'apparaît pas !
\end{itemize}
\end{attention}

\courssub{Tableau de synthèse et organigramme de décision}

\begin{aretenir}[Les quatre outils du dénombrement]
\begin{center}
\begin{tabularx}{\linewidth}{|l|X|X|X|}
\hline
\rowcolor{popblueL} \textbf{Situation} & \textbf{Ordre ?} & \textbf{Répétition ?} & \textbf{Nombre} \\
\hline
$p$-uplets (listes) & Oui & Oui & $n^p$ \\
\hline
Arrangements & Oui & Non & $A_n^p = \dfrac{n!}{(n-p)!}$ \\
\hline
Combinaisons (parties) & Non & Non & $C_n^p = \dfrac{n!}{p!(n-p)!}$ \\
\hline
Permutations & Oui (tous les éléments) & Non & $n!$ \\
\hline
\end{tabularx}
\end{center}
On a toujours la relation : $A_n^p = C_n^p \times p!$ (on choisit les éléments, puis on les ordonne). En particulier $C_n^p = \dfrac{A_n^p}{p!}$, ce qui explique pourquoi une combinaison « coûte » $p!$ fois moins cher qu'un arrangement.
\end{aretenir}

\begin{methode}[L'organigramme de choix de l'outil]
\begin{enumerate}
\item \textbf{Modéliser} : identifier $n$ (objets disponibles) et $p$ (choix à faire).
\item \textbf{Question 1} : l'ordre compte-t-il ?
\item \textbf{Question 2} : la répétition est-elle possible ?
\item \textbf{Appliquer} la formule correspondante (voir organigramme ci-dessous).
\item \textbf{Vérifier} l'ordre de grandeur et la cohérence (un nombre de délégations ne peut pas dépasser le nombre de listes ordonnées correspondantes).
\end{enumerate}
\end{methode}

\begin{popfigure}
\begin{center}
\begin{tikzpicture}[
  node distance=0.9cm and 1.1cm,
  question/.style={rectangle, rounded corners, draw=popblue, fill=popblueL, thick, text width=3.6cm, align=center, font=\small\bfseries},
  outil/.style={rectangle, rounded corners, draw=popgreen, fill=popgreenL, thick, text width=3.4cm, align=center, font=\small},
  lab/.style={font=\small\itshape, text=popdark}
]
\node[question] (q1) {L'ordre compte-t-il ?};
\node[question, below left=1.1cm and 0.2cm of q1] (q2) {Répétition possible ?};
\node[outil, below right=1.1cm and 0.2cm of q1, align=center] (comb) {Combinaisons\\ $C_n^p$};
\node[outil, below left=1.0cm and 0.1cm of q2, align=center] (pupl) {$p$-uplets\\ $n^p$};
\node[outil, below right=1.0cm and 0.1cm of q2, align=center] (arr) {Arrangements\\ $A_n^p$};
\node[outil, right=0.8cm of arr, align=center] (perm) {Cas $p=n$ :\\ permutations $n!$};
\draw[-{Stealth}, thick, popblue] (q1) -- node[lab, above left]{Oui} (q2);
\draw[-{Stealth}, thick, popblue] (q1) -- node[lab, above right]{Non} (comb);
\draw[-{Stealth}, thick, popblue] (q2) -- node[lab, above left]{Oui} (pupl);
\draw[-{Stealth}, thick, popblue] (q2) -- node[lab, above right]{Non} (arr);
\draw[-{Stealth}, thick, popgreen, dashed] (arr) -- (perm);
\end{tikzpicture}
\end{center}
\legende{Organigramme de décision : deux questions suffisent pour choisir l'outil.}
\end{popfigure}

\courssub{Problèmes mixtes : combiner les principes}

Les vrais problèmes (et les sujets de Probatoire !) exigent souvent de \cle{décomposer} la situation en étapes, puis d'appliquer le \cle{principe multiplicatif} entre les étapes, chaque étape utilisant l'un des quatre outils.

\begin{methode}[Stratégie générale face à un problème mixte]
\begin{enumerate}
\item Découper le résultat final en \emph{étapes indépendantes} (« choisir ceci, \emph{puis} choisir cela »).
\item Pour chaque étape, appliquer l'organigramme de décision.
\item \textbf{Multiplier} les nombres obtenus (principe multiplicatif) si les étapes se succèdent ; \textbf{additionner} si l'on traite des cas disjoints (principe additif : « cas 1 \emph{ou} cas 2 »).
\item Vérifier : le résultat doit être un entier, inférieur au nombre total sans contrainte.
\end{enumerate}
\end{methode}

\begin{exoresolu}[Délégation avec contrainte]
Dans une classe de Première C de 35 élèves (15 filles et 20 garçons), combien peut-on former de délégations de 3 élèves comportant \emph{exactement une fille} ?
\tcblower
\textbf{Modélisation.} Former une telle délégation, c'est : choisir 1 fille parmi 15 (étape 1), \emph{puis} choisir 2 garçons parmi 20 (étape 2). Dans chaque étape, l'ordre ne compte pas (une délégation n'est pas ordonnée) et il n'y a pas de répétition : ce sont des combinaisons.

\textbf{Calcul.}
\begin{align*}
N &= C_{15}^1 \times C_{20}^2 = 15 \times \frac{20 \times 19}{2} = 15 \times 190 = 2850.
\end{align*}
\textbf{Vérification.} Le nombre total de délégations sans contrainte est $C_{35}^3 = \dfrac{35 \times 34 \times 33}{6} = 6545$ : notre résultat $2850 < 6545$, ce qui est cohérent. Il y a donc \textbf{2850 délégations} possibles.
\end{exoresolu}

\begin{exoresolu}[Plaques d'immatriculation]
Une plaque provisoire est formée de 2 lettres distinctes choisies parmi les 26 lettres de l'alphabet, suivies de 3 chiffres (répétition autorisée, le premier chiffre pouvant être 0). Combien de plaques différentes peut-on produire ?
\tcblower
\textbf{Étape 1 — les lettres.} Deux lettres \emph{distinctes} et \emph{ordonnées} (AB $\neq$ BA) : arrangement, $A_{26}^2 = 26 \times 25 = 650$.

\textbf{Étape 2 — les chiffres.} Trois chiffres ordonnés avec répétition possible : $10^3 = 1000$.

\textbf{Principe multiplicatif.}
\[ N = 650 \times 1000 = 650\,000 \text{ plaques.} \]
De quoi immatriculer largement tous les véhicules neufs d'une région !
\end{exoresolu}

\begin{exemplebox}[Menu au restaurant et tournoi de football]
\textbf{Menu.} Un restaurant de Bonanjo propose 3 entrées, 4 plats et 2 desserts. Un menu complet (une entrée, un plat, un dessert) : $3 \times 4 \times 2 = 24$ menus (principe multiplicatif pur).

\textbf{Tournoi.} 8 équipes de quartiers de Yaoundé participent à un tournoi où chaque équipe rencontre chaque autre une seule fois. Un match est une paire non ordonnée : $C_8^2 = \dfrac{8 \times 7}{2} = 28$ matchs. Si l'on demande le nombre de classements possibles des 8 équipes : $8! = 40\,320$.
\end{exemplebox}

\begin{attention}[Pièges classiques des problèmes mixtes]
\begin{itemize}
\item \emph{« au moins une fille »} : passer par le \textbf{complémentaire} : total $-$ (aucune fille). Exemple : $C_{35}^3 - C_{20}^3 = 6545 - 1140 = 5405$ délégations avec au moins une fille.
\item Ne jamais additionner des résultats d'étapes successives : étapes successives $\Rightarrow$ on \textbf{multiplie} ; cas exclusifs $\Rightarrow$ on \textbf{additionne}.
\item Dans un « bureau » (président, vice-président, trésorier), les rôles sont distincts : c'est un arrangement, pas une combinaison.
\item Attention aux contraintes de position (« le président est un garçon ») : traiter d'abord la case contrainte, puis les autres.
\end{itemize}
\end{attention}

\begin{savaistu}[Pascal, Fermat et la naissance des probabilités]
C'est à propos de jeux de hasard que le dénombrement est devenu une science. En 1654, le chevalier de Méré, joueur invétéré, soumet à Blaise \cle{Pascal} un problème de partage de gains dans une partie interrompue. Pascal échange alors avec Pierre de \cle{Fermat} : de cette correspondance naît le calcul des probabilités, dont le dénombrement est le socle (une probabilité, c'est souvent un quotient de deux dénombrements !). Aujourd'hui, ces outils servent partout : sécurité des transactions Mobile Money (combien de codes PIN à 4 chiffres ? $10^4 = 10\,000$), tirages de loterie, tests médicaux, et même la conception des réseaux de téléphonie de MTN ou Orange.
\end{savaistu}

\begin{exercice}[Pour s'entraîner]
\begin{enumerate}
\item Un code de cadenas comporte 3 chiffres. Combien de codes possibles ? Même question si les chiffres sont deux à deux distincts.
\item Le CEP d'un lycée de Douala veut élire un bureau de 3 membres (président, secrétaire, trésorier) parmi 10 candidats. Combien de bureaux possibles ?
\item Combien de comités de 4 personnes, comportant au moins une femme, peut-on former parmi 6 femmes et 5 hommes ?
\end{enumerate}
\end{exercice}

\begin{corrige}[Exercice]
\begin{enumerate}
\item Avec répétition : $10^3 = 1000$ codes. Chiffres distincts : $A_{10}^3 = 10 \times 9 \times 8 = 720$.
\item Les postes sont distincts, l'ordre compte : $A_{10}^3 = 10 \times 9 \times 8 = 720$ bureaux.
\item Par le complémentaire : $C_{11}^4 - C_{5}^4 = 330 - 5 = 325$ comités.
\end{enumerate}
\end{corrige}

\begin{aretenir}[L'essentiel de la leçon]
\begin{itemize}
\item Deux questions guident tout dénombrement : \emph{l'ordre compte-t-il ?} et \emph{y a-t-il répétition ?}
\item Ordre + répétition : $n^p$ ; ordre sans répétition : $A_n^p$ ; ni ordre ni répétition : $C_n^p$ ; tous les éléments ordonnés : $n!$.
\item Relation charnière : $A_n^p = C_n^p \times p!$.
\item Problèmes mixtes : décomposer en étapes, multiplier entre étapes, additionner entre cas disjoints, et penser au complémentaire pour « au moins un ».
\end{itemize}
\end{aretenir}

\newpage

\coursec{Activité d'intégration}

\courssub{Objectifs de l'activité}
\noindent Cette activité d'intégration vise à mobiliser l'ensemble des outils de dénombrement étudiés dans la leçon (diagrammes de Venn, arbres de choix, $p$-uplets, arrangements, combinaisons, triangle de Pascal, formule du binôme) pour résoudre une situation complexe, authentique et contextualisée. À l'issue de ce travail, tu seras capable de :
\begin{itemize}
    \item Analyser un problème concret pour identifier les sous-problèmes de dénombrement qu'il contient.
    \item Choisir et justifier l'outil adapté (Venn, arbre, $A_n^p$, $C_n^p$, binôme) pour chaque sous-problème.
    \item Articuler plusieurs raisonnements de dénombrement au sein d'une même résolution.
    \item Rédiger un rapport mathématique structuré, clair et rigoureux, tel qu'exigé au Probatoire.
\end{itemize}

\courssub{Situation : La semaine sportive et culturelle du Lycée Bilingue de Nkolbisson}
\noindent Le Lycée Bilingue de Nkolbisson (Yaoundé) organise sa semaine sportive et culturelle annuelle. Le comité d'organisation, dont tu fais partie, doit finaliser plusieurs points logistiques avant l'ouverture officielle. Voici les données chiffrées transmises par l'intendance et les délégués de classes :

\begin{enumerate}
    \item \textbf{Tournoi de football (phase de poules) :} Huit classes participent : $6^{\text{ème}}A$, $6^{\text{ème}}B$, $5^{\text{ème}}A$, $5^{\text{ème}}B$, $4^{\text{ème}}A$, $4^{\text{ème}}B$, $3^{\text{ème}}A$, $3^{\text{ème}}B$. Chaque classe doit rencontrer \emph{exactement une fois} chacune des sept autres classes. L'arbitre central, M. Essomba, doit établir le calendrier complet des matchs.
    \item \textbf{Phase finale (podium) :} À l'issue des poules, les 8 classes sont classées. Pour la cérémonie de clôture, seul le podium (1\textsuperscript{er}, 2\textsuperscript{ème}, 3\textsuperscript{ème}) sera récompensé. Le comité veut connaître le nombre de configurations possibles pour ce podium afin de prévoir les médailles (or, argent, bronze) gravées à l'avance.
    \item \textbf{Restauration (Cantine scolaire) :} La cantine propose un menu unique « Spécial Semaine Culturelle » composé d'un plat principal, d'un accompagnement et d'une boisson. Les stocks disponibles ce jour-là sont :
    \begin{itemize}
        \item Plats principaux (3) : Ndolè, Poulet DG, Poisson braisé.
        \item Accompagnements (2) : Bâton de manioc, Riz blanc.
        \item Boissons (2) : Jus de gingembre (Zom), Eau minérale (Tangui).
    \end{itemize}
    La gérante, Mme Fouda, veut afficher le nombre total de menus différents possibles pour que les élèves choisissent rapidement.
    \item \textbf{Sondage « Préférences culinaires » :} Un sondage a été réalisé auprès de 60 élèves de Terminale tirés au sort. Résultats :
    \begin{itemize}
        \item 40 élèves aiment le Ndolè.
        \item 35 élèves aiment le Poulet DG.
        \item 18 élèves aiment \emph{les deux} (Ndolè et Poulet DG).
    \end{itemize}
    Le comité souhaite savoir combien d'élèves n'aiment \emph{ni} le Ndolè \emph{ni} le Poulet DG, pour prévoir une alternative (spaghettis).
    \item \textbf{Équipes de soutien (Animation) :} Chaque classe doit présenter une équipe de 4 élèves pour l'animation culturelle (chants, danses, sketchs). Dans la classe de $1^{\text{ère}}D$, 10 élèves se portent volontaires. Le délégué de classe doit former l'équipe. L'ordre de passage sur scène n'est pas encore fixé ; il s'agit seulement de choisir les 4 membres.
    \item \textbf{Bonus (Vérification théorique) :} Le président du comité, passionné de maths, propose de développer $(x+1)^5$ à l'aide du triangle de Pascal pour retrouver le nombre total de parties d'un ensemble à 5 éléments (lien avec le nombre de sous-équipes possibles parmi 5 volontaires).
\end{enumerate}

\courssub{Consignes}
\noindent Travaille en groupe de 3 ou 4 élèves. Pour chaque point de la situation, tu dois :
\begin{enumerate}
    \item Identifier clairement l'outil de dénombrement nécessaire (diagramme de Venn, arbre, $p$-uplets, arrangement, combinaison, binôme).
    \item Justifier ton choix en explicitant les critères : \emph{ordre important ou non ?}, \emph{répétition autorisée ou non ?}, \emph{choix simultanés ou successifs ?}, \emph{intersections d'ensembles ?}.
    \item Effectuer le calcul complet en détaillant les étapes (formule, simplification des factorielles, calcul numérique final).
    \item Pour le point 4, réaliser un diagramme de Venn clair et légendé.
    \item Pour le point 6, construire le triangle de Pascal jusqu'au rang 5 et écrire le développement.
    \item Rédiger collectivement un rapport final structuré : \textbf{Introduction} (contexte), \textbf{Développement} (6 parties numérotées avec outil, justification, calcul, résultat), \textbf{Conclusion} (synthèse des outils utilisés et difficultés rencontrées).
\end{enumerate}

\courssub{Ressources à mobiliser}
\noindent Rappel des formules et propriétés clés du cours.

\begin{propriete}[Principes de base]
\textbf{Principe de l'addition :} Si un événement $A$ peut se réaliser de $a$ façons et un événement $B$ \emph{incompatible} avec $A$ de $b$ façons, alors « $A$ ou $B$ » se réalise de $a+b$ façons.\\
\textbf{Principe de la multiplication :} Si un processus se décompose en $k$ étapes successives, avec $n_1, n_2, \dots, n_k$ choix possibles à chaque étape (le nombre de choix à une étape ne dépendant pas des choix précédents), le nombre total de façons est $n_1 \times n_2 \times \dots \times n_k$.
\end{propriete}

\begin{propriete}[$p$-uplets (listes ordonnées)]
Nombre de $p$-uplets (ou $p$-listes) de $p$ éléments choisis parmi $n$ éléments distincts :
\begin{itemize}
    \item \textbf{Avec répétition :} $n^p$ (choix indépendants, ordre important, un même élément peut être repris).
    \item \textbf{Sans répétition :} $A_n^p = n \times (n-1) \times \dots \times (n-p+1) = \dfrac{n!}{(n-p)!}$ (arrangements), pour $1 \le p \le n$.
\end{itemize}
\end{propriete}

\begin{propriete}[Permutations]
Cas particulier $p=n$ : le nombre de permutations de $n$ éléments distincts est
\[
A_n^n = n! = n \times (n-1) \times \dots \times 2 \times 1.
\]
Par convention $0! = 1$.
\end{propriete}

\begin{propriete}[Combinaisons]
Le nombre de parties à $p$ éléments (sous-ensembles) d'un ensemble à $n$ éléments, avec $0 \le p \le n$, est :
\[
C_n^p = \binom{n}{p} = \frac{A_n^p}{p!} = \frac{n!}{p!\,(n-p)!}.
\]
Ici l'ordre n'est \textbf{pas} important et il n'y a \textbf{pas} de répétition.\\
Propriétés : $C_n^p = C_n^{\,n-p}$ (symétrie) ; relation de Pascal : $C_n^p = C_{n-1}^{\,p-1} + C_{n-1}^{\,p}$ pour $1 \le p \le n-1$.
\end{propriete}

\begin{propriete}[Diagramme de Venn et cardinal]
Pour deux ensembles finis $A$ et $B$ inclus dans un univers fini $U$ :
\[
\mathrm{Card}(A \cup B) = \mathrm{Card}(A) + \mathrm{Card}(B) - \mathrm{Card}(A \cap B),
\qquad
\mathrm{Card}(\overline{A \cup B}) = \mathrm{Card}(U) - \mathrm{Card}(A \cup B).
\]
\end{propriete}

\begin{propriete}[Triangle de Pascal et formule du binôme (Newton)]
Pour tous réels $a$ et $b$ et tout entier naturel $n$ :
\[
(a+b)^n = \sum_{k=0}^{n} \binom{n}{k} a^{\,n-k} b^{\,k}.
\]
En posant $a=b=1$ : $\displaystyle\sum_{k=0}^{n} \binom{n}{k} = 2^n$, qui est le nombre de parties d'un ensemble à $n$ éléments.
\end{propriete}

\courssub{Démarche et corrigé commenté}
\begin{exoresolu}[Rapport de l'activité : Organisation de la semaine culturelle]
\noindent \textbf{Introduction}\\
Notre groupe a pour mission d'assister le comité d'organisation du Lycée Bilingue de Nkolbisson. Nous allons traiter les six points logistiques en identifiant pour chacun la structure de dénombrement sous-jacente, puis en appliquant l'outil adapté avec un calcul détaillé.

\tcblower

\noindent \textbf{Point 1 : Nombre de matchs de la phase de poules (8 classes)}\\
\textbf{Analyse :} Un match oppose \textbf{deux classes distinctes}. L'ordre n'a pas d'importance : le match « $6^{\text{ème}}A$ contre $6^{\text{ème}}B$ » est le même que « $6^{\text{ème}}B$ contre $6^{\text{ème}}A$ ». Une classe ne joue pas contre elle-même (pas de répétition). C'est un choix \textbf{simultané} de 2 éléments parmi 8, sans ordre.\\
\textbf{Outil :} \textbf{Combinaison} $C_8^2$.\\
\textbf{Calcul :}
\[
C_8^2 = \binom{8}{2} = \frac{8!}{2!\,(8-2)!} = \frac{8 \times 7 \times 6!}{2 \times 1 \times 6!} = \frac{8 \times 7}{2} = \frac{56}{2} = 28.
\]
\textbf{Résultat :} \textbf{28 matchs} au total. M. Essomba devra planifier 28 créneaux horaires.

\noindent \textbf{Point 2 : Nombre de podiums possibles (3 premières places sur 8 classes)}\\
\textbf{Analyse :} Il s'agit de classer les 3 premières classes. L'ordre est \textbf{capital} : être 1\textsuperscript{er} (médaille d'or) est différent d'être 2\textsuperscript{ème} (argent) ou 3\textsuperscript{ème} (bronze). On choisit 3 classes distinctes parmi 8 \textbf{en les rangeant}. C'est un choix \textbf{ordonné} sans répétition.\\
\textbf{Outil :} \textbf{Arrangement} $A_8^3$.\\
\textbf{Calcul :}
\[
A_8^3 = \frac{8!}{(8-3)!} = \frac{8!}{5!} = 8 \times 7 \times 6 = 336.
\]
\textbf{Résultat :} \textbf{336 podiums} différents possibles. Le comité doit donc prévoir 336 configurations potentielles (ou commander des médailles génériques non gravées au nom des classes).

\noindent \textbf{Point 3 : Nombre de menus à la cantine (arbre de choix)}\\
\textbf{Analyse :} Un menu se construit par \textbf{étapes successives} et \textbf{indépendantes} : Étape 1 : choix du plat (3 options). Étape 2 : choix de l'accompagnement (2 options). Étape 3 : choix de la boisson (2 options). Le principe de la multiplication s'applique directement, et l'arbre de choix le visualise.\\
\textbf{Outil :} \textbf{Arbre de choix} / \textbf{principe de la multiplication} (un menu est un triplet (plat ; accompagnement ; boisson), c'est-à-dire un $3$-uplet d'éléments pris dans trois ensembles différents).
\begin{popfigure}
\begin{center}
\begin{tikzpicture}[level distance=1.4cm,
  level 1/.style={sibling distance=3.2cm},
  level 2/.style={sibling distance=1.6cm},
  level 3/.style={sibling distance=0.85cm},
  scale=0.72, transform shape]
\node[circle, draw=popblue, fill=popblueL, thick] {Départ}
    child { node[circle, draw=poporange, fill=poporangeL, thick] {Ndolè}
        child { node[circle, draw=popgreen, fill=popgreenL, thick] {Bâton}
            child { node[rectangle, draw=poppink, fill=poppinkL] {Zom} }
            child { node[rectangle, draw=poppink, fill=poppinkL] {Tangui} }
        }
        child { node[circle, draw=popgreen, fill=popgreenL, thick] {Riz}
            child { node[rectangle, draw=poppink, fill=poppinkL] {Zom} }
            child { node[rectangle, draw=poppink, fill=poppinkL] {Tangui} }
        }
    }
    child { node[circle, draw=poporange, fill=poporangeL, thick] {Poulet DG}
        child { node[circle, draw=popgreen, fill=popgreenL, thick] {Bâton}
            child { node[rectangle, draw=poppink, fill=poppinkL] {Zom} }
            child { node[rectangle, draw=poppink, fill=poppinkL] {Tangui} }
        }
        child { node[circle, draw=popgreen, fill=popgreenL, thick] {Riz}
            child { node[rectangle, draw=poppink, fill=poppinkL] {Zom} }
            child { node[rectangle, draw=poppink, fill=poppinkL] {Tangui} }
        }
    }
    child { node[circle, draw=poporange, fill=poporangeL, thick] {Poisson}
        child { node[circle, draw=popgreen, fill=popgreenL, thick] {Bâton}
            child { node[rectangle, draw=poppink, fill=poppinkL] {Zom} }
            child { node[rectangle, draw=poppink, fill=poppinkL] {Tangui} }
        }
        child { node[circle, draw=popgreen, fill=popgreenL, thick] {Riz}
            child { node[rectangle, draw=poppink, fill=poppinkL] {Zom} }
            child { node[rectangle, draw=poppink, fill=poppinkL] {Tangui} }
        }
    };
\end{tikzpicture}
\end{center}
\legende{Arbre des menus possibles : 3 plats, puis 2 accompagnements, puis 2 boissons ; chaque chemin de la racine à une feuille est un menu.}
\end{popfigure}
\textbf{Calcul :} $3 \times 2 \times 2 = 12$. L'arbre possède bien 12 feuilles.\\
\textbf{Résultat :} \textbf{12 menus différents} possibles. Mme Fouda affichera 12 formules.

\noindent \textbf{Point 4 : Élèves n'aimant ni Ndolè ni Poulet DG (diagramme de Venn)}\\
\textbf{Analyse :} On considère deux ensembles : $N$ (élèves aimant le Ndolè) et $P$ (élèves aimant le Poulet DG), dans l'univers $U$ des 60 élèves sondés. On connaît $\mathrm{Card}(N)=40$, $\mathrm{Card}(P)=35$, $\mathrm{Card}(N \cap P)=18$. On cherche $\mathrm{Card}(\overline{N \cup P})$.\\
\textbf{Outil :} \textbf{Diagramme de Venn} + \textbf{formule d'inclusion-exclusion}.\\
\textbf{Remplissage du diagramme (de l'intersection vers l'extérieur) :}
\begin{itemize}
    \item Intersection $N \cap P$ : 18 élèves.
    \item Ndolè seulement ($N \setminus P$) : $40 - 18 = 22$ élèves.
    \item Poulet DG seulement ($P \setminus N$) : $35 - 18 = 17$ élèves.
    \item Total aimant au moins l'un des deux ($N \cup P$) : $22 + 18 + 17 = 57$ élèves.
\end{itemize}
\begin{popfigure}
\begin{center}
\begin{tikzpicture}[scale=0.85]
    \draw[thick, popdark] (-4,-3.6) rectangle (7,2.6);
    \node[popdark] at (6.3,2.1) {$U$};
    \begin{scope}[blend group=soft light]
        \fill[popblueL] (0,0) circle (2.5cm);
        \fill[poporangeL] (3,0) circle (2.5cm);
    \end{scope}
    \draw[thick, popblue] (0,0) circle (2.5cm);
    \draw[thick, poporange] (3,0) circle (2.5cm);
    \node[popblue] at (-2.2, 1.6) {$N$};
    \node[poporange] at (5.2, 1.6) {$P$};
    \node at (1.5, 0) {\textbf{18}};
    \node at (-1.2, -0.6) {\textbf{22}};
    \node at (4.2, -0.6) {\textbf{17}};
    \node at (1.5, -3.1) {Hors $N \cup P$ : \textbf{3} élèves};
\end{tikzpicture}
\end{center}
\legende{Diagramme de Venn des préférences culinaires : $N$ (Ndolè), $P$ (Poulet DG), univers $U$ de 60 élèves.}
\end{popfigure}
\textbf{Calcul formel :}
\[
\mathrm{Card}(N \cup P) = \mathrm{Card}(N) + \mathrm{Card}(P) - \mathrm{Card}(N \cap P) = 40 + 35 - 18 = 57.
\]
\[
\mathrm{Card}(\overline{N \cup P}) = \mathrm{Card}(U) - \mathrm{Card}(N \cup P) = 60 - 57 = 3.
\]
\textbf{Résultat :} \textbf{3 élèves} n'aiment ni le Ndolè ni le Poulet DG. Il faut prévoir 3 portions de spaghettis.

\noindent \textbf{Point 5 : Choix de l'équipe d'animation (4 élèves sur 10 volontaires)}\\
\textbf{Analyse :} On doit former un \textbf{groupe} de 4 élèves. L'ordre de sélection n'importe pas : l'équipe $\{$Amina, Bruno, Carine, David$\}$ est la même que $\{$David, Carine, Bruno, Amina$\}$. Pas de répétition (un élève ne peut pas être choisi deux fois). Choix simultané.\\
\textbf{Outil :} \textbf{Combinaison} $C_{10}^4$.\\
\textbf{Calcul :}
\[
C_{10}^4 = \binom{10}{4} = \frac{10!}{4!\,6!} = \frac{10 \times 9 \times 8 \times 7}{4 \times 3 \times 2 \times 1} = \frac{5040}{24} = 210.
\]
Détail de la simplification : $\dfrac{10 \times 9 \times 8 \times 7}{4 \times 3 \times 2 \times 1} = \dfrac{10}{2} \times \dfrac{9}{3} \times \dfrac{8}{4} \times 7 = 5 \times 3 \times 2 \times 7 = 210$.\\
\textbf{Résultat :} \textbf{210 équipes} différentes possibles pour la $1^{\text{ère}}D$.

\noindent \textbf{Point 6 : Bonus — Développement de $(x+1)^5$ et parties d'un ensemble à 5 éléments}\\
\textbf{Analyse :} On utilise le triangle de Pascal pour obtenir les coefficients $\binom{5}{k}$. Le nombre de parties d'un ensemble à $n$ éléments est $2^n = \sum_{k=0}^n \binom{n}{k}$. Pour $n=5$, on attend $2^5 = 32$.\\
\textbf{Outil :} \textbf{Triangle de Pascal} et \textbf{formule du binôme}.\\
\textbf{Construction du triangle (jusqu'au rang 5) :} chaque nombre est la somme des deux nombres situés au-dessus de lui (relation de Pascal).
\begin{center}
\begin{tabular}{c c c c c c c c c c cc}
Rang 0 : & & & & & 1 & & & & & \\
Rang 1 : & & & & 1 & & 1 & & & & \\
Rang 2 : & & & 1 & & 2 & & 1 & & & \\
Rang 3 : & & 1 & & 3 & & 3 & & 1 & & \\
Rang 4 : & 1 & & 4 & & 6 & & 4 & & 1 & \\
Rang 5 : & 1 & & 5 & & 10 & & 10 & & 5 & & 1
\end{tabular}
\end{center}
\textbf{Développement :}
\begin{align*}
(x+1)^5 &= \binom{5}{0}x^5 + \binom{5}{1}x^4 + \binom{5}{2}x^3 + \binom{5}{3}x^2 + \binom{5}{4}x + \binom{5}{5} \\
&= x^5 + 5x^4 + 10x^3 + 10x^2 + 5x + 1.
\end{align*}
\textbf{Vérification :} en posant $x=1$, la somme des coefficients vaut $1+5+10+10+5+1 = 32 = 2^5$.\\
Cela correspond bien au nombre de parties d'un ensemble à 5 éléments : chaque coefficient $\binom{5}{k}$ compte les parties à $k$ éléments, et leur somme compte toutes les parties.\\
\textbf{Résultat :} développement validé, lien combinaisons/parties confirmé.

\noindent \textbf{Conclusion}\\
Cette activité a permis de mettre en pratique l'ensemble de la boîte à outils du dénombrement :
\begin{itemize}
    \item \textbf{Combinaisons} (Points 1 et 5) pour des choix simultanés sans ordre (matchs, équipe).
    \item \textbf{Arrangements} (Point 2) pour des choix ordonnés sans répétition (podium).
    \item \textbf{Arbre / principe de la multiplication} (Point 3) pour des processus séquentiels à étapes indépendantes (menu).
    \item \textbf{Diagramme de Venn} (Point 4) pour gérer intersections et complémentaires dans un univers fini.
    \item \textbf{Triangle de Pascal / binôme} (Point 6) pour le lien structural entre coefficients binomiaux, combinaisons et cardinal de l'ensemble des parties ($2^n$).
\end{itemize}
La difficulté principale a été de \textbf{ne pas confondre arrangement et combinaison} (Points 1 et 2) et de bien remplir le diagramme de Venn en partant de l'intersection (Point 4). La rédaction du rapport impose une rigueur dans la justification du choix de l'outil, point crucial pour le Probatoire.
\end{exoresolu}

\courssub{Bilan de l'activité}
\noindent Cette situation d'intégration a mis en évidence les compétences suivantes, conformes au programme officiel de Première (séries C, D, E, TI) :

\begin{itemize}
    \item \textbf{Modélisation :} traduire un énoncé réel (tournoi, cantine, sondage) en objets mathématiques (ensembles, $p$-uplets, arrangements, combinaisons).
    \item \textbf{Choix raisonné de l'outil :} la discrimination entre \emph{ordre important} (arrangement, $p$-uplet) et \emph{ordre non important} (combinaison, Venn) est le pivot central. L'arbre s'impose naturellement pour les choix séquentiels à étapes indépendantes.
    \item \textbf{Maîtrise des calculs avec factorielles :} simplification de $\dfrac{n!}{(n-p)!}$ et de $\dfrac{n!}{p!\,(n-p)!}$ sans calculatrice pour des nombres raisonnables, utilisation efficace de la symétrie $C_n^p = C_n^{\,n-p}$.
    \item \textbf{Représentation graphique :} construction et lecture d'un diagramme de Venn à deux ensembles (remplissage de l'intersection vers l'extérieur) et d'un arbre de choix (niveaux = étapes, feuilles = issues).
    \item \textbf{Lien algébre-combinatoire :} le bonus ancre la compréhension du triangle de Pascal comme table des combinaisons et comme générateur des coefficients binomiaux, reliant le dénombrement des parties ($2^n$) au développement de $(1+1)^n$.
    \item \textbf{Communication mathématique :} la rédaction du rapport final (Introduction, Développement structuré, Conclusion) entraîne à l'argumentation écrite exigée à l'épreuve écrite de Mathématiques du Probatoire.
\end{itemize}

\begin{attention}[Pièges à éviter pour le Probatoire]
\begin{itemize}
    \item \textbf{Confondre $A_n^p$ et $C_n^p$ :} demande-toi toujours : « Si j'échange deux éléments choisis, obtient-on une situation différente ? » Oui $\rightarrow$ arrangement / $p$-uplet. Non $\rightarrow$ combinaison.
    \item \textbf{Oublier l'intersection dans un diagramme de Venn :} ne jamais écrire $\mathrm{Card}(A \cup B) = \mathrm{Card}(A) + \mathrm{Card}(B)$ en général. Toujours soustraire $\mathrm{Card}(A \cap B)$, sauf si les ensembles sont disjoints.
    \item \textbf{Mauvaise lecture de l'arbre :} le nombre total de chemins (feuilles) est le \emph{produit} des nombres de branches à chaque niveau, pas leur somme.
    \item \textbf{Notation factorielle :} $0! = 1$ (piège classique dans $C_n^n = C_n^0 = 1$). Écris les étapes de simplification ($n! = n \times (n-1)!$) pour éviter les erreurs de calcul.
    \item \textbf{Contexte « avec / sans répétition » :} par défaut en dénombrement classique (équipes, matchs, podiums), le choix se fait \textbf{sans répétition} (un objet ne peut pas être choisi deux fois). La répétition ne concerne que les codes, mots de passe, ou tirages avec remise.
    \item \textbf{Rédaction :} annonce toujours l'outil utilisé et justifie-le (ordre ? répétition ?) \emph{avant} d'écrire la formule : c'est cette justification qui est évaluée au Probatoire.
\end{itemize}
\end{attention}

\newpage

\coursec{Méthodes et exercices résolus}

\coursec{Dénombrement}

\courssub{Cardinal d'un ensemble fini et premiers outils}

\begin{definition}[Cardinal d'un ensemble fini]
Le \cle{cardinal} d'un ensemble fini $E$, noté $\mathrm{card}(E)$ ou $|E|$, est le nombre d'éléments de $E$. Par convention, $\mathrm{card}(\varnothing)=0$.
\end{definition}

\begin{definition}[Diagramme de Venn et tableau à double entrée]
Un \cle{diagramme de Venn} représente les ensembles par des disques à l'intérieur d'un rectangle (l'univers $\Omega$). Les intersections, réunions et complémentaires se visualisent par les zones communes ou disjointes.
Un \cle{tableau à double entrée} (ou tableau de contingence) croise deux caractères : les lignes pour un caractère, les colonnes pour l'autre, les cases centrales pour les effectifs croisés, les marges pour les totaux partiels et le total général.
\end{definition}

\begin{propriete}[Formule du crible (Inclusion-Exclusion pour deux ensembles)]
Pour tous ensembles finis $A$ et $B$ :
\[ \mathrm{card}(A\cup B)=\mathrm{card}(A)+\mathrm{card}(B)-\mathrm{card}(A\cap B). \]
\emph{Preuve exigible au Probatoire :} On compte les éléments de $A$, puis ceux de $B$. Les éléments de $A\cap B$ ont été comptés deux fois, on les retire une fois.
\end{propriete}

\begin{methode}[Dénombrer avec un diagramme de Venn ou un tableau à double entrée]
Pour dénombrer les éléments d'une population décrite par deux caractères :
\begin{enumerate}
\item \cle{Identifier} les ensembles : $\Omega$ (total), $A$ et $B$ (sous-ensembles). Traduire chaque donnée en cardinal : $\mathrm{card}(A)$, $\mathrm{card}(B)$, $\mathrm{card}(A\cap B)$, etc.
\item \cle{Commencer par l'intersection} : dans le diagramme, remplir $A\cap B$ en premier, puis « $A$ seulement » ($\mathrm{card}(A)-\mathrm{card}(A\cap B)$) et « $B$ seulement ».
\item \cle{Appliquer la formule} : $\mathrm{card}(A\cup B)=\mathrm{card}(A)+\mathrm{card}(B)-\mathrm{card}(A\cap B)$.
\item \cle{Vérifier} que la somme de toutes les zones (y compris l'extérieur $\overline{A\cup B}$) redonne $\mathrm{card}(\Omega)$.
\end{enumerate}
\textbf{Réflexe} : « seulement » = on retire l'intersection ; « au moins un » = réunion ; « ni... ni... » = complémentaire de la réunion.
\end{methode}

\begin{popfigure}
\begin{tikzpicture}[scale=0.9]
  \def\R{2.5}
  \def\sep{1.5}
  \begin{scope}[fill opacity=0.25]
    \fill[popblue] (-\sep,0) circle (\R);
    \fill[poporange] (\sep,0) circle (\R);
  \end{scope}
  \draw[thick, popblue] (-\sep,0) circle (\R) node[left=3.8cm] {$M$ (MTN)};
  \draw[thick, poporange] (\sep,0) circle (\R) node[right=3.8cm] {$O$ (Orange)};
  \draw[thick] (-5.5,-3.5) rectangle (5.5,3.5) node[above left] {$\Omega$ (60 élèves)};
  \node[font=\bfseries\large] at (-\sep,0.8) {23};
  \node[font=\bfseries\large] at (\sep,0.8) {16};
  \node[font=\bfseries\large] at (0,0) {12};
  \node[font=\bfseries] at (0,-2.8) {9 (aucune)};
\end{tikzpicture}
\legende{Diagramme de Venn : répartition des 60 élèves selon la possession de puces MTN et/or Orange}
\end{popfigure}

\begin{exoresolu}[Abonnés MTN et Orange dans une classe]
Dans une classe de Première C de 60 élèves, 35 élèves possèdent une puce MTN, 28 possèdent une puce Orange et 12 possèdent les deux puces.
\begin{enumerate}
\item Combien d'élèves possèdent au moins une des deux puces ?
\item Combien d'élèves ne possèdent aucune des deux puces ?
\item Combien d'élèves possèdent uniquement une puce MTN ?
\end{enumerate}
\tcblower
Notons $\Omega$ l'ensemble des 60 élèves, $M$ l'ensemble des élèves possédant une puce MTN et $O$ celui des élèves possédant une puce Orange. Les données se traduisent par :
\[ \mathrm{card}(\Omega)=60,\quad \mathrm{card}(M)=35,\quad \mathrm{card}(O)=28,\quad \mathrm{card}(M\cap O)=12. \]

\textbf{1.} « Au moins une des deux puces » correspond à la réunion $M\cup O$. D'après la formule du crible :
\[ \mathrm{card}(M\cup O)=\mathrm{card}(M)+\mathrm{card}(O)-\mathrm{card}(M\cap O)=35+28-12=51. \]
On soustrait 12 car les élèves aux deux puces auraient été comptés deux fois. Donc \textbf{51 élèves} possèdent au moins une des deux puces.

\textbf{2.} « Aucune des deux puces » est le complémentaire de $M\cup O$ dans $\Omega$ :
\[ \mathrm{card}\left(\overline{M\cup O}\right)=\mathrm{card}(\Omega)-\mathrm{card}(M\cup O)=60-51=9. \]
Donc \textbf{9 élèves} ne possèdent aucune des deux puces.

\textbf{3.} « Uniquement MTN » désigne la zone $M$ privée de l'intersection :
\[ \mathrm{card}(M)-\mathrm{card}(M\cap O)=35-12=23. \]
Donc \textbf{23 élèves} possèdent uniquement une puce MTN.

\textbf{Vérification} : zones du diagramme : MTN seulement $23$, Orange seulement $28-12=16$, les deux $12$, aucune $9$. Total : $23+16+12+9=60$. La vérification est conforme.
\end{exoresolu}

\begin{attention}[Pièges classiques sur les diagrammes de Venn]
\begin{itemize}
\item Ne pas confondre « $A$ seulement » ($\mathrm{card}(A)-\mathrm{card}(A\cap B)$) et $\mathrm{card}(A)$.
\item Ne pas oublier la zone « ni $A$ ni $B$ » (complémentaire de $A\cup B$) dans le total.
\item Dans un tableau à double entrée, vérifier que la somme des totaux partiels lignes = somme des totaux partiels colonnes = total général.
\end{itemize}
\end{attention}

\courssub{Arbres de choix et p-uplets}

\begin{definition}[p-uplet (ou p-liste) et principe multiplicatif]
Un \cle{p-uplet} (ou \cle{p-liste}) est une liste \emph{ordonnée} de $p$ éléments. Le \cle{principe multiplicatif} affirme que si un choix se décompose en $k$ étapes successives indépendantes, avec $n_1$ choix à l'étape 1, $n_2$ à l'étape 2, ..., $n_k$ à l'étape $k$, alors le nombre total de choix est $n_1\times n_2\times\cdots\times n_k$.
\end{definition}

\begin{propriete}[Nombre de p-uplets]
Soit un ensemble $E$ à $n$ éléments. Le nombre de p-uplets de longueur $p$ ($p\in\mathbb{N}^*$) formés avec des éléments de $E$ est :
\begin{itemize}
\item \textbf{Avec répétition autorisée} : $n^p$ (chaque case offre $n$ choix).
\item \textbf{Sans répétition} : $n(n-1)\cdots(n-p+1)$, noté $A_n^p$ (arrangement de $p$ parmi $n$).
\end{itemize}
\end{propriete}

\begin{experience}[Arbre de choix : visualiser le principe multiplicatif]
Construisons l'arbre des codes à 2 chiffres formés avec $\{0,1\}$ (répétition autorisée).
\begin{popfigure}
\begin{tikzpicture}[level distance=1.8cm, sibling distance=2.5cm, edge from parent/.style={draw, popblue, thick}]
  \node[circle, draw, popblue, fill=popblueL, minimum width=0.7cm] {Départ}
    child {node[circle, draw, poporange, fill=poporangeL, minimum width=0.7cm] {0}
      child {node[rectangle, draw, popgreen, fill=popgreenL, rounded corners] {00}}
      child {node[rectangle, draw, popgreen, fill=popgreenL, rounded corners] {01}}
    }
    child {node[circle, draw, poporange, fill=poporangeL, minimum width=0.7cm] {1}
      child {node[rectangle, draw, popgreen, fill=popgreenL, rounded corners] {10}}
      child {node[rectangle, draw, popgreen, fill=popgreenL, rounded corners] {11}}
    };
\end{tikzpicture}
\legende{Arbre de choix pour un code à 2 chiffres (0 ou 1) : $2^2=4$ codes possibles}
\end{popfigure}
Chaque chemin de la racine à une feuille correspond à un code. On compte $2\times 2 = 2^2 = 4$ feuilles.
\end{experience}

\begin{methode}[Compter les p-uplets avec le principe multiplicatif]
Pour compter des listes \emph{ordonnées} de $p$ éléments (codes, numéros, classements, mots) :
\begin{enumerate}
\item \cle{Se poser deux questions} : l'ordre compte-t-il ? (sinon ce n'est pas un p-uplet) ; un élément peut-il être répété ?
\item \cle{Construire l'arbre ou les cases} : on dessine $p$ cases successives et on écrit dans chacune le nombre de choix possibles.
\item \cle{Appliquer le principe multiplicatif} : on multiplie les nombres de choix.
\begin{itemize}
\item Avec répétition autorisée : $n^p$ listes de $p$ éléments pris parmi $n$.
\item Sans répétition : $n(n-1)\cdots(n-p+1)=A_n^p$ (arrangement).
\end{itemize}
\item \cle{Conclure} par une phrase qui rappelle ce qui a été compté.
\end{enumerate}
\textbf{Réflexe} : code PIN, plaque d'immatriculation, numéro de téléphone $\Rightarrow$ répétition possible ; podium, bureau élu, tirage sans remise $\Rightarrow$ pas de répétition.
\end{methode}

\begin{exoresolu}[Codes et classements]
\textbf{Partie A.} Un code de coffre est formé de 4 chiffres choisis parmi les chiffres de 0 à 9 (un chiffre peut être répété). Combien de codes différents peut-on former ?

\textbf{Partie B.} À l'issue d'une course cycliste entre Douala et Yaoundé, 8 coureurs sont en lice pour les trois premières places (médailles d'or, d'argent et de bronze distinctes). Combien de podiums différents sont possibles ?
\tcblower
\textbf{Partie A.} Un code est une liste \emph{ordonnée} de $p=4$ chiffres pris parmi $n=10$ chiffres, \emph{avec répétition} possible (le code 1111 est valide). On remplit 4 cases, chacune offrant 10 choix :
\[ \underbrace{10}_{1^{\text{er}}\text{ chiffre}}\times\underbrace{10}_{2^{\text{e}}}\times\underbrace{10}_{3^{\text{e}}}\times\underbrace{10}_{4^{\text{e}}}=10^4=\num{10000}. \]
On peut donc former \textbf{\num{10000} codes différents} (de 0000 à 9999).

\textbf{Partie B.} Un podium est une liste \emph{ordonnée} (l'ordre des médailles compte) de $p=3$ coureurs pris parmi $n=8$, \emph{sans répétition} (un coureur ne peut occuper deux places). C'est le nombre d'arrangements :
\[ A_8^3=8\times 7\times 6=336. \]
Il y a donc \textbf{336 podiums différents} possibles.
\end{exoresolu}

\courssub{Factorielle, permutations et arrangements}

\begin{definition}[Factorielle]
Pour $n\in\mathbb{N}^*$, la \cle{factorielle} de $n$, notée $n!$, est le produit des $n$ premiers entiers naturels non nuls :
\[ n!=n\times(n-1)\times\cdots\times 2\times 1. \]
Par convention, $0!=1$.
\end{definition}

\begin{definition}[Permutation]
Une \cle{permutation} de $n$ éléments distincts est une liste ordonnée de ces $n$ éléments (chacun utilisé exactement une fois). C'est un arrangement de $n$ parmi $n$.
\end{definition}

\begin{propriete}[Nombre de permutations]
Le nombre de permutations de $n$ éléments distincts est $n!$.
\emph{Valeurs usuelles à connaître :} $3!=6$, $4!=24$, $5!=120$, $6!=720$, $7!=5040$.
\end{propriete}

\begin{methode}[Compter les permutations]
Pour ranger \emph{tous} les éléments d'un ensemble de $n$ objets \emph{distincts} dans un ordre donné :
\begin{enumerate}
\item \cle{Vérifier} que tous les objets sont utilisés, chacun une seule fois, et que l'ordre compte.
\item \cle{Appliquer} : le nombre de permutations de $n$ éléments est $n!$.
\item \cle{Calculer} proprement la factorielle.
\end{enumerate}
\textbf{Réflexe} : « de combien de façons peut-on ranger / aligner / classer $n$ personnes ? » $\Rightarrow n!$.
\end{methode}

\begin{exoresolu}[Photo de classe]
Les 6 membres du bureau d'une coopérative agricole de Bafoussam se placent en file indienne pour une photo officielle.
\begin{enumerate}
\item De combien de façons différentes peuvent-ils se placer ?
\item De combien de façons peuvent-ils se placer si le président et le secrétaire veulent être aux deux extrémités de la file ?
\end{enumerate}
\tcblower
\textbf{1.} Placer 6 personnes distinctes en file, c'est choisir un ordre sur les 6 personnes : c'est une permutation de 6 éléments.
\[ 6!=6\times 5\times 4\times 3\times 2\times 1=720. \]
Il y a \textbf{720 façons} différentes de se placer.

\textbf{2.} On décompose le placement en étapes successives :
\begin{itemize}
\item \emph{Étape 1 : choisir qui occupe les extrémités.} Le président et le secrétaire occupent les deux bouts : 2 façons (président à gauche ou à droite).
\item \emph{Étape 2 : placer les 4 autres membres} sur les 4 places centrales : c'est une permutation de 4 éléments, soit $4!=24$ façons.
\end{itemize}
Par le principe multiplicatif :
\[ 2\times 4!=2\times 24=48. \]
Il y a donc \textbf{48 façons} de se placer avec le président et le secrétaire aux extrémités.
\end{exoresolu}

\begin{definition}[Arrangement]
Un \cle{arrangement} de $p$ éléments parmi $n$ ($0\le p\le n$) est une liste \emph{ordonnée sans répétition} de $p$ éléments distincts choisis parmi $n$. On note $A_n^p$ ce nombre.
\end{definition}

\begin{propriete}[Formule des arrangements]
\[ A_n^p=n(n-1)\cdots(n-p+1)=\frac{n!}{(n-p)!} \qquad\text{(produit de $p$ facteurs).} \]
Cas particuliers : $A_n^0=1$, $A_n^1=n$, $A_n^n=n!$.
\end{propriete}

\begin{methode}[Calculer $A_n^p$]
\begin{enumerate}
\item \cle{Vérifier les conditions} : $0\le p\le n$, ordre important, pas de répétition.
\item \cle{Appliquer la formule} : $A_n^p=\dfrac{n!}{(n-p)!}$ (produit de $p$ facteurs décroissants).
\item \cle{Contrôler le nombre de facteurs} : le produit doit contenir exactement $p$ facteurs.
\end{enumerate}
\textbf{Réflexe} : bureau (président, trésorier, secrétaire) élu parmi $n$ candidats, tiercé dans l'ordre, mots de $p$ lettres distinctes $\Rightarrow A_n^p$.
\end{methode}

\begin{exoresolu}[Élection du bureau d'un club]
Le club de mathématiques d'un lycée de Bafoussam compte 15 membres. On veut élire un bureau composé d'un président, d'un vice-président et d'un trésorier, un membre ne pouvant cumuler deux postes.
\begin{enumerate}
\item Combien de bureaux différents peut-on former ?
\item Calculer $A_{15}^3$ à l'aide des factorielles et retrouver le résultat.
\end{enumerate}
\tcblower
\textbf{1.} Les trois postes sont \emph{distincts} (l'ordre compte : être élu président n'est pas être élu trésorier) et un membre ne peut pas occuper deux postes (pas de répétition). Il s'agit donc du nombre d'arrangements de 3 éléments parmi 15 :
\[ A_{15}^3=15\times 14\times 13. \]
Calcul : $15\times 14=210$, puis $210\times 13=\num{2730}$. On peut former \textbf{\num{2730} bureaux différents}.

\textbf{2.} Avec la formule factorielle :
\[ A_{15}^3=\frac{15!}{(15-3)!}=\frac{15!}{12!}=\frac{15\times 14\times 13\times \cancel{12!}}{\cancel{12!}}=15\times 14\times 13=\num{2730}. \]
On retrouve bien le même résultat : \textbf{\num{2730} bureaux}. La simplification de $12!$ évite tout calcul de grandes factorielles.
\end{exoresolu}

\courssub{Combinaisons et parties d'un ensemble}

\begin{definition}[Combinaison]
Une \cle{combinaison} de $p$ éléments parmi $n$ ($0\le p\le n$) est un \cle{sous-ensemble} (ou \cle{partie}) à $p$ éléments : l'ordre ne compte pas. On note $C_n^p$ ou $\binom{n}{p}$ ce nombre.
\end{definition}

\begin{propriete}[Formule des combinaisons et propriétés]
\[ C_n^p=\frac{A_n^p}{p!}=\frac{n!}{p!\,(n-p)!}. \]
Propriétés fondamentales (exigibles) :
\begin{itemize}
\item Symétrie : $C_n^p=C_n^{\,n-p}$.
\item Valeurs limites : $C_n^0=C_n^n=1$, $C_n^1=n$.
\item Relation de Pascal : $C_n^p=C_{n-1}^{p-1}+C_{n-1}^{p}$ (pour $1\le p\le n-1$).
\item Relation arrangements/combinaisons : $A_n^p=p!\,C_n^p$.
\end{itemize}
\end{propriete}

\begin{propriete}[Nombre de parties d'un ensemble fini]
Un ensemble à $n$ éléments possède exactement $2^n$ parties (y compris $\varnothing$ et l'ensemble lui-même) :
\[ \sum_{p=0}^{n}C_n^p=2^n. \]
\end{propriete}

\begin{experience}[Arbre des parties : pourquoi $2^n$ ?]
Pour un ensemble $E=\{a,b,c\}$ ($n=3$), chaque élément a deux destins : « dans la partie » ou « pas dans la partie ». L'arbre des choix donne $2\times2\times2=2^3=8$ parties :
$\varnothing,\ \{a\},\ \{b\},\ \{c\},\ \{a,b\},\ \{a,c\},\ \{b,c\},\ \{a,b,c\}$.
\end{experience}

\begin{methode}[Calculer $C_n^p$ et le nombre de parties]
\begin{enumerate}
\item \cle{Vérifier} que l'ordre n'intervient pas (délégation, tirage simultané, équipe sans rôles).
\item \cle{Appliquer la formule} : $C_n^p=\dfrac{n!}{p!\,(n-p)!}$.
\item \cle{Utiliser les propriétés} pour simplifier : symétrie $C_n^p=C_n^{n-p}$, relation de Pascal.
\item \cle{Nombre total de parties} : $2^n$.
\end{enumerate}
\textbf{Réflexe} : « on choisit simultanément », « combien de groupes / délégations / mains de cartes » $\Rightarrow C_n^p$.
\end{methode}

\begin{exoresolu}[Délégation pour la FENASSCO]
Pour représenter son lycée aux jeux FENASSCO, un professeur d'EPS doit choisir 4 athlètes parmi 10 élèves volontaires.
\begin{enumerate}
\item Combien de délégations différentes peut-il former ?
\item Combien de délégations contiennent exactement les deux meilleurs athlètes, Amina et Bassong ?
\item Combien de parties (y compris la partie vide) l'ensemble des 10 volontaires possède-t-il ?
\end{enumerate}
\tcblower
\textbf{1.} Les 4 athlètes sont choisis \emph{simultanément}, sans ordre ni rôle particulier : c'est une combinaison de 4 éléments parmi 10 :
\[ C_{10}^4=\frac{10!}{4!\,6!}=\frac{10\times 9\times 8\times 7}{4\times 3\times 2\times 1}=\frac{\num{5040}}{24}=210. \]
Le professeur peut former \textbf{210 délégations} différentes.

\textbf{2.} Si Amina et Bassong sont imposés, il reste à choisir $4-2=2$ athlètes parmi les $10-2=8$ autres volontaires, toujours sans ordre :
\[ C_8^2=\frac{8!}{2!\,6!}=\frac{8\times 7}{2}=28. \]
Il y a \textbf{28 délégations} contenant Amina et Bassong.

\textbf{3.} Le nombre total de parties d'un ensemble à 10 éléments est :
\[ 2^{10}=\num{1024}. \]
L'ensemble des volontaires possède donc \textbf{\num{1024} parties} (y compris l'ensemble vide et l'ensemble plein).
\end{exoresolu}

\courssub{Triangle de Pascal et formule du binôme}

\begin{definition}[Triangle de Pascal]
Le \cle{triangle de Pascal} est un triangle de nombres où la ligne $n$ (en partant de $n=0$) contient les coefficients $C_n^0, C_n^1, \dots, C_n^n$. Il se construit par la règle : chaque nombre (sauf les 1 aux bords) est la somme des deux nombres juste au-dessus de lui (relation de Pascal).
\end{definition}

\begin{propriete}[Formule du binôme de Newton]
Pour tout $n\in\mathbb{N}$ et pour tous réels (ou complexes) $a,b$ :
\[ (a+b)^n=\sum_{p=0}^{n}C_n^p\,a^{\,n-p}b^{\,p}
= C_n^0 a^n + C_n^1 a^{n-1}b + \cdots + C_n^{n-1} a b^{n-1} + C_n^n b^n. \]
\emph{Preuve non exigible au Probatoire, mais la relation de Pascal et l'utilisation de la formule le sont.}
\end{propriete}

\begin{methode}[Développer $(a+b)^n$]
\begin{enumerate}
\item \cle{Construire le triangle de Pascal} jusqu'à la ligne $n$ grâce à $C_n^p=C_{n-1}^{p-1}+C_{n-1}^{p}$.
\item \cle{Écrire la formule du binôme} : $(a+b)^n=\sum_{p=0}^{n}C_n^p a^{n-p}b^p$.
\item \cle{Disposer les puissances} : exposants de $a$ décroissent de $n$ à $0$, ceux de $b$ croissent de $0$ à $n$ ; somme des exposants $=n$ à chaque terme.
\item \cle{Attention aux signes} : pour $(a-b)^n$, écrire $(a+(-b))^n$ : les termes où l'exposant de $b$ est impair sont négatifs.
\end{enumerate}
\end{methode}

\begin{exoresolu}[Développements au programme du Probatoire]
\begin{enumerate}
\item Construire les lignes $n=0$ à $n=5$ du triangle de Pascal.
\item Développer et réduire $(a+b)^5$.
\item Développer et réduire $(2x-1)^4$.
\end{enumerate}
\tcblower
\textbf{1.} Triangle de Pascal (chaque nombre est la somme des deux nombres au-dessus) :
\begin{center}
\begin{tabular}{c|cccccc}
$n=0$ & 1 & & & & & \\
$n=1$ & 1 & 1 & & & & \\
$n=2$ & 1 & 2 & 1 & & & \\
$n=3$ & 1 & 3 & 3 & 1 & & \\
$n=4$ & 1 & 4 & 6 & 4 & 1 & \\
$n=5$ & 1 & 5 & 10 & 10 & 5 & 1 \\
\end{tabular}
\end{center}

\textbf{2.} La ligne $n=5$ donne les coefficients $1,\ 5,\ 10,\ 10,\ 5,\ 1$. Par la formule du binôme :
\begin{align*}
(a+b)^5 &= C_5^0 a^5 + C_5^1 a^4 b + C_5^2 a^3 b^2 + C_5^3 a^2 b^3 + C_5^4 a b^4 + C_5^5 b^5 \\
&= a^5+5a^4b+10a^3b^2+10a^2b^3+5ab^4+b^5.
\end{align*}

\textbf{3.} On écrit $(2x-1)^4=(2x+(-1))^4$ et on applique le binôme avec $a=2x$, $b=-1$, et les coefficients de la ligne $n=4$ : $1,\ 4,\ 6,\ 4,\ 1$ :
\begin{align*}
(2x-1)^4 &= 1\cdot(2x)^4+4\cdot(2x)^3(-1)+6\cdot(2x)^2(-1)^2+4\cdot(2x)(-1)^3+1\cdot(-1)^4 \\
&= 16x^4-4\times 8x^3+6\times 4x^2-4\times 2x+1 \\
&= 16x^4-32x^3+24x^2-8x+1.
\end{align*}
Les signes alternent car les puissances impaires de $(-1)$ valent $-1$. Conclusion : $(2x-1)^4=16x^4-32x^3+24x^2-8x+1$.
\end{exoresolu}

\begin{savaistu}[Le triangle de Pascal au Cameroun et dans le monde]
Bien que nommé d'après Blaise Pascal (1623--1662), ce triangle était connu en Chine (Jia Xian, XI\textsuperscript{e} siècle ; Yang Hui, XIII\textsuperscript{e} siècle), en Perse (Al-Karaji, X\textsuperscript{e} siècle ; Omar Khayyam, XI\textsuperscript{e} siècle) et en Inde (Pingala, II\textsuperscript{e} siècle av. J.-C.) bien avant lui. Au Cameroun, les coefficients binomiaux interviennent dans la modélisation des réseaux de téléphonie mobile (répartition des canaux), dans les algorithmes de compression d'images (transformée en cosinus discrète) et dans les calculs de probabilités pour les jeux de hasard (Loterie Nationale) ou les risques agricoles (assurance indicielle).
\end{savaistu}

\courssub{Synthèse : choisir le bon outil de dénombrement}

\begin{aretenir}[Tableau récapitulatif des outils de dénombrement]
\begin{center}
\begin{tabularx}{\linewidth}{|l|X|X|X|}\hline
\textbf{Situation} & \textbf{Ordre} & \textbf{Répétition} & \textbf{Formule} \\ \hline
Permutation ($p=n$) & Oui & Non & $n!$ \\ \hline
Arrangement ($p<n$) & Oui & Non & $A_n^p=\dfrac{n!}{(n-p)!}$ \\ \hline
p-uplet avec répétition & Oui & Oui & $n^p$ \\ \hline
Combinaison (partie) & Non & Non & $C_n^p=\dfrac{n!}{p!(n-p)!}$ \\ \hline
Nombre de parties & -- & -- & $2^n$ \\ \hline
\end{tabularx}
\end{center}
Relation clé (exigible) : $A_n^p = p! \times C_n^p$.
\end{aretenir}

\begin{methode}[Le bon réflexe face à un énoncé]
Avant tout calcul, se poser les deux questions décisives :
\begin{enumerate}
\item \cle{L'ordre compte-t-il ?}
\begin{itemize}
\item \textbf{Oui} $\to$ liste / arrangement. Un élément peut-il se répéter ?
\begin{itemize}
\item Répétition possible $\to$ $n^p$.
\item Pas de répétition $\to$ $A_n^p=\dfrac{n!}{(n-p)!}$ (et $n!$ si $p=n$ : permutation).
\end{itemize}
\item \textbf{Non} $\to$ combinaison : $C_n^p=\dfrac{n!}{p!(n-p)!}$.
\end{itemize}
\item \cle{Traduire les contraintes} (« au moins un », « exactement deux », « ni... ni... ») avec les ensembles : passage au complémentaire, formule du crible.
\item \cle{Rédiger} : annoncer l'outil choisi et \emph{justifier} le choix (ordre, répétition) avant d'écrire la formule.
\end{enumerate}
\end{methode}

\begin{exoresolu}[Tirages de jetons — synthèse complète]
Une urne contient 6 jetons numérotés de 1 à 6. On en tire successivement 3 jetons.
\begin{enumerate}
\item On tire avec remise et on note le numéro à chaque tirage : combien de résultats ?
\item On tire sans remise et on note les numéros dans l'ordre de sortie : combien de résultats ?
\item On tire les 3 jetons simultanément : combien de résultats ?
\end{enumerate}
\tcblower
\textbf{1.} Tirage successif \emph{avec remise} : l'ordre compte et un même jeton peut ressortir plusieurs fois. C'est une 3-liste d'éléments pris parmi 6 :
\[ 6^3=216. \]
Il y a \textbf{216 résultats} possibles.

\textbf{2.} Tirage successif \emph{sans remise} : l'ordre compte mais il n'y a pas de répétition. C'est un arrangement de 3 éléments parmi 6 :
\[ A_6^3=6\times 5\times 4=120. \]
Il y a \textbf{120 résultats} possibles.

\textbf{3.} Tirage \emph{simultané} : l'ordre ne compte pas et il n'y a pas de répétition. C'est une combinaison de 3 éléments parmi 6 :
\[ C_6^3=\frac{6!}{3!\,3!}=\frac{6\times 5\times 4}{3\times 2\times 1}=\frac{120}{6}=20. \]
Il y a \textbf{20 résultats} possibles.

\textbf{Remarque de synthèse} : on vérifie la relation $A_6^3=C_6^3\times 3!$ ($120=20\times 6$) : chaque partie de 3 jetons peut être ordonnée de $3!=6$ façons. Cette relation $A_n^p=p!\,C_n^p$ est exigible au Probatoire.
\end{exoresolu}

\begin{attention}[Les 5 erreurs qui coûtent des points au Probatoire]
\begin{enumerate}
\item \cle{Confondre arrangement et combinaison.} Avant d'écrire $A_n^p$ ou $C_n^p$, justifie en une phrase : « l'ordre compte car les postes sont distincts » ou « l'ordre ne compte pas car le tirage est simultané ». Un résultat juste avec le mauvais outil annoncé est sanctionné.
\item \cle{Oublier la condition $p\le n$ et la non-répétition.} $A_n^p$ et $C_n^p$ n'existent que pour $0\le p\le n$ et supposent des éléments \emph{distincts}. Si l'énoncé autorise la répétition (code, numéro), c'est $n^p$, jamais $A_n^p$.
\item \cle{Se tromper dans le nombre de facteurs de $A_n^p$.} $A_n^p$ est un produit de \emph{exactement $p$ facteurs} : $A_8^3=8\times 7\times 6$ (et non $8\times 7\times 6\times 5$). Compte les facteurs avant de calculer.
\item \cle{Oublier de diviser par $p!$ dans $C_n^p$.} L'erreur classique : écrire $C_{10}^4=\dfrac{10\times 9\times 8\times 7}{1}$. Le dénominateur $p!$ est indispensable ; retiens $C_n^p=\dfrac{A_n^p}{p!}$.
\item \cle{Mal gérer les signes dans $(a-b)^n$.} Écris $(a-b)^n=(a+(-b))^n$ : le terme en $C_n^p a^{n-p}b^p$ porte le signe $(-1)^p$. Vérifie aussi que la somme des exposants de chaque terme vaut $n$ : c'est un contrôle rapide et efficace.
\end{enumerate}
\end{attention}

\begin{exercice}[Entraînement Probatoire - Dénombrement]
\begin{enumerate}
\item Un code d'accès à un distributeur automatique de billets (GAB) de la BICEC est composé de 4 chiffres (0 à 9). Combien de codes possibles ? Si le code ne doit pas commencer par 0, combien ?
\item On tire 5 cartes d'un jeu de 52 cartes. Combien de mains différentes (ordre indifférent) ? Combien de façons de les poser en ligne sur la table (ordre important) ?
\item Développer $(3x-2)^4$ à l'aide du triangle de Pascal.
\item Dans une classe de 40 élèves, 22 font du football, 18 font du basketball, 10 font les deux. Combien ne font aucun des deux sports ? Représenter par un diagramme de Venn.
\item Démontrer que $C_n^p = C_n^{n-p}$ à l'aide de la formule factorielle. Donner une interprétation ensembliste.
\end{enumerate}
\end{exercice}

\begin{corrige}[Exercice 1]
\textbf{1.} Code à 4 chiffres, répétition autorisée : $10^4 = 10\,000$.
Si le 1er chiffre $\neq 0$ : 9 choix pour le 1er, 10 pour les 3 suivants $\Rightarrow 9\times 10^3 = 9\,000$.

\textbf{2.} Main de 5 cartes (ordre indifférent) : $C_{52}^5 = \frac{52!}{5!47!} = 2\,598\,960$.
Poser en ligne (ordre important) : $A_{52}^5 = 52\times 51\times 50\times 49\times 48 = 311\,875\,200$.

\textbf{3.} $(3x-2)^4 = (3x+(-2))^4$. Coeffs ligne 4 : 1, 4, 6, 4, 1.
$= 1(3x)^4 + 4(3x)^3(-2) + 6(3x)^2(-2)^2 + 4(3x)(-2)^3 + 1(-2)^4$
$= 81x^4 - 216x^3 + 216x^2 - 96x + 16$.

\textbf{4.} $\Omega=40$, $F=22$, $B=18$, $F\cap B=10$.
$F\cup B = 22+18-10=30$.
Ni l'un ni l'autre : $40-30=10$ élèves.
Venn : Football seulement $12$, Basketball seulement $8$, Les deux $10$, Aucun $10$. Total $40$.

\textbf{5.} $C_n^{n-p} = \frac{n!}{(n-p)!(n-(n-p))!} = \frac{n!}{(n-p)!p!} = C_n^p$.
Interprétation : Choisir $p$ éléments parmi $n$, c'est équivalent à choisir les $n-p$ éléments qu'on \emph{laisse de côté}.
\end{corrige}

\newpage

\coursec{Exercices}

\courssub{Je vérifie mes connaissances}

\begin{exercice}[QCM : le bon outil]
Pour chaque situation, choisir la bonne réponse en justifiant brièvement.
\begin{enumerate}
\item Le nombre de façons de ranger 5 livres distincts sur une étagère est :
\begin{enumerate}
\item $5^5$ \qquad \item $5!$ \qquad \item $C_5^5$ \qquad \item $A_5^2$
\end{enumerate}
\item Le nombre de délégations de 3 élèves choisies parmi 10 élèves (sans ordre) est :
\begin{enumerate}
\item $A_{10}^3$ \qquad \item $10^3$ \qquad \item $C_{10}^3$ \qquad \item $3!$
\end{enumerate}
\item Le nombre de codes PIN à 4 chiffres (chiffres pouvant se répéter) est :
\begin{enumerate}
\item $A_{10}^4$ \qquad \item $C_{10}^4$ \qquad \item $10^4$ \qquad \item $4!$
\end{enumerate}
\item Le nombre de parties d'un ensemble à 8 éléments est :
\begin{enumerate}
\item $8!$ \qquad \item $2^8$ \qquad \item $C_8^2$ \qquad \item $8^2$
\end{enumerate}
\end{enumerate}
\end{exercice}

\begin{exercice}[Vrai ou faux, en justifiant]
Dire si chaque affirmation est vraie ou fausse, en justifiant la réponse.
\begin{enumerate}
\item Pour tous entiers $n$ et $p$ avec $0 \le p \le n$ : $A_n^p = C_n^p \times p!$.
\item $C_n^p = C_n^{n-p}$ pour tout $0 \le p \le n$.
\item Le coefficient de $a^3 b^2$ dans le développement de $(a+b)^5$ est $C_5^3$.
\item Si $A$ et $B$ sont deux parties d'un ensemble fini $E$, alors $\mathrm{Card}(A \cup B) = \mathrm{Card}(A) + \mathrm{Card}(B)$.
\item $C_n^0 + C_n^1 + \cdots + C_n^n = 2^n$.
\end{enumerate}
\end{exercice}

\begin{exercice}[Définitions et formules]
\begin{enumerate}
\item Donner la définition du \cle{cardinal} d'un ensemble fini.
\item Rappeler la formule donnant $\mathrm{Card}(A \cup B)$ en fonction de $\mathrm{Card}(A)$, $\mathrm{Card}(B)$ et $\mathrm{Card}(A \cap B)$.
\item Définir $n!$ (factorielle $n$) et donner la valeur de $0!$.
\item Donner les expressions de $A_n^p$ et de $C_n^p$ à l'aide de factorielles.
\item Énoncer la \cle{formule du binôme de Newton} pour $(a+b)^n$.
\end{enumerate}
\end{exercice}

\begin{exercice}[Calculs directs]
Calculer sans calculatrice, en détaillant les étapes :
\begin{enumerate}
\item $6!$ ; $\dfrac{9!}{7!}$ ; $\dfrac{12!}{10! \times 2!}$.
\item $A_8^2$ ; $A_7^3$ ; $A_{10}^3$.
\item $C_6^2$ ; $C_9^4$ ; $C_{10}^3$.
\item Vérifier que $A_{10}^3 = C_{10}^3 \times 3!$.
\item Simplifier, pour $n \ge 1$ : $\dfrac{(n+1)!}{(n-1)!}$.
\end{enumerate}
\end{exercice}

\courssub{Je m'entraîne}

\begin{exercice}[Plaques d'immatriculation]
Une plaque d'immatriculation camerounaise est formée de $2$ lettres (alphabet de $26$ lettres) suivies de $3$ chiffres.
\begin{enumerate}
\item Combien de plaques peut-on former au total ?
\item Combien de plaques peut-on former si tous les caractères (lettres et chiffres) sont deux à deux distincts ?
\item Combien de plaques commencent par la lettre C et se terminent par un chiffre pair ?
\end{enumerate}
\end{exercice}

\begin{exercice}[Le mot CAMEROUN]
On utilise les $8$ lettres du mot CAMEROUN (toutes distinctes).
\begin{enumerate}
\item Combien de mots de $5$ lettres distinctes (ayant un sens ou non) peut-on former ?
\item Parmi ces mots, combien commencent par une voyelle ?
\item Combien de mots de $8$ lettres peut-on former en utilisant toutes les lettres ? Combien commencent par C et se terminent par N ?
\end{enumerate}
\end{exercice}

\begin{exercice}[Formation d'un comité]
Un comité de $4$ personnes doit être constitué parmi $7$ hommes et $5$ femmes.
\begin{enumerate}
\item Combien de comités peut-on former au total ?
\item Combien de comités comportent exactement $2$ femmes ?
\item Combien de comités comportent au moins une femme ?
\item Combien de comités comportent au plus un homme ?
\end{enumerate}
\end{exercice}

\begin{exercice}[Triangle de Pascal et binôme]
\begin{enumerate}
\item Construire le triangle de Pascal jusqu'à la ligne $n = 6$.
\item Développer $(a+b)^5$ et $(a-b)^4$.
\item À l'aide du triangle de Pascal, développer $(2x - 3)^4$ et déterminer le coefficient de $x^2$.
\item Montrer que $C_n^0 + C_n^1 + \cdots + C_n^n = 2^n$ (on pourra utiliser la formule du binôme avec des valeurs bien choisies de $a$ et $b$). En déduire le nombre de parties d'un ensemble de $6$ éléments.
\end{enumerate}
\end{exercice}

\courssub{Je me prépare au Probatoire}

\begin{exercice}[Sports au lycée — type Probatoire]
Dans un lycée de $500$ élèves, on a constaté que $280$ élèves pratiquent le football, $150$ pratiquent le handball et $60$ pratiquent les deux sports.
\begin{enumerate}
\item Représenter la situation par un diagramme de Venn en faisant figurer les effectifs de chaque zone.
\item Déterminer le nombre d'élèves pratiquant :
\begin{enumerate}
\item le football seulement ;
\item le handball seulement ;
\item au moins un des deux sports.
\end{enumerate}
\item En déduire le nombre d'élèves ne pratiquant aucun de ces deux sports.
\item On choisit au hasard $3$ élèves du lycée pour une enquête.
\begin{enumerate}
\item Combien de groupes de $3$ élèves peut-on former ?
\item Combien de ces groupes sont formés uniquement d'élèves pratiquant le football ?
\item Combien de ces groupes contiennent au moins un élève ne pratiquant aucun des deux sports ?
\end{enumerate}
\end{enumerate}
\end{exercice}

\begin{exercice}[Finale du 100 m — type Probatoire]
Huit coureurs, dont deux Camerounais nommés Essomba et Mbarga, participent à la finale du $100$ mètres lors des championnats nationaux. On suppose qu'il n'y a pas d'ex æquo.
\begin{enumerate}
\item Déterminer le nombre de classements complets possibles des $8$ coureurs.
\item On appelle \emph{podium} l'ensemble ordonné des trois premiers (or, argent, bronze).
\begin{enumerate}
\item Combien de podiums différents sont possibles ?
\item Combien de podiums comportent Essomba en première position ?
\item Combien de podiums comportent au moins l'un des deux Camerounais ?
\end{enumerate}
\item À l'issue de la course, on choisit $3$ coureurs parmi les $8$ pour un contrôle antidopage (le choix est sans ordre).
\begin{enumerate}
\item Combien de choix possibles y a-t-il ?
\item Combien de choix comportent exactement l'un des deux Camerounais ?
\end{enumerate}
\item Vérifier, à l'aide des résultats précédents, la relation $A_8^3 = C_8^3 \times 3!$.
\end{enumerate}
\end{exercice}

\begin{exercice}[Restaurant et menus — problème de synthèse]
Un restaurant de Yaoundé propose chaque midi $4$ entrées, $5$ plats de résistance et $3$ desserts. Un \emph{menu complet} est constitué d'une entrée, d'un plat et d'un dessert.
\begin{enumerate}
\item Combien de menus complets différents peut-on composer ?
\item Deux amis, Alima et Bassile, commandent chacun un menu complet. Combien de couples de menus peut-on obtenir si :
\begin{enumerate}
\item les deux amis choisissent librement (leurs menus peuvent être identiques) ?
\item les deux amis veulent des menus entièrement différents (entrée, plat et dessert tous différents) ?
\end{enumerate}
\item Le restaurateur souhaite afficher sur une ardoise un échantillon de $2$ entrées, $3$ plats et $2$ desserts choisis parmi sa carte. Combien d'ardoises différentes peut-il composer ?
\item Le restaurateur note que le nombre total de menus complets s'écrit $4 \times 5 \times 3 = 60$. Il affirme : « le nombre de couples ordonnés $(\text{entrée}, \text{plat})$ est $A_9^2$ ». Cette affirmation est-elle correcte ? Justifier et donner la valeur exacte.
\item Pour une fête, le restaurant prépare $n$ plats différents et constate qu'il peut former exactement $120$ menus complets avec ses $4$ entrées et $3$ desserts. Déterminer $n$, puis calculer $C_n^2$ et interpréter ce résultat dans le contexte.
\end{enumerate}
\end{exercice}

\newpage

\coursec{Corrigés des exercices}

\begin{corrige}[Exercice 1 — QCM : le bon outil]
\begin{enumerate}
\item \textbf{Réponse b : $5!$.} Ranger $5$ livres distincts sur une étagère, c'est ordonner les $5$ livres : c'est une \cle{permutation} de $5$ éléments, il y en a $5! = 120$.
\item \textbf{Réponse c : $C_{10}^3$.} Une délégation est un choix \emph{sans ordre} et \emph{sans répétition} de $3$ élèves parmi $10$ : c'est une \cle{combinaison}, $C_{10}^3 = 120$.
\item \textbf{Réponse c : $10^4$.} Un code PIN est une \cle{$p$-liste avec répétition} : l'ordre compte et les chiffres peuvent se répéter. Il y a $10$ choix pour chacune des $4$ positions, soit $10^4 = 10\,000$ codes.
\item \textbf{Réponse b : $2^8$.} Le nombre de parties d'un ensemble à $n$ éléments est $2^n$ (pour chaque élément, deux choix : le prendre ou non). Ici $2^8 = 256$.
\end{enumerate}
\end{corrige}

\begin{corrige}[Exercice 2 — Vrai ou faux]
\begin{enumerate}
\item \textbf{Vrai.} Pour former un arrangement de $p$ éléments parmi $n$, on peut d'abord choisir les $p$ éléments ($C_n^p$ façons), puis les ordonner ($p!$ façons). D'où $A_n^p = C_n^p \times p!$. Vérification par les formules : $C_n^p \times p! = \dfrac{n!}{p!(n-p)!} \times p! = \dfrac{n!}{(n-p)!} = A_n^p$.
\item \textbf{Vrai.} Choisir $p$ éléments parmi $n$ revient à choisir les $n-p$ éléments que l'on \emph{laisse}. Par le calcul : $C_n^{n-p} = \dfrac{n!}{(n-p)!\,p!} = C_n^p$.
\item \textbf{Vrai.} Dans le développement de $(a+b)^5$, le terme en $a^3 b^2$ a pour coefficient $C_5^2 = C_5^3 = 10$ (on choisit les $2$ facteurs qui fournissent $b$, ou de manière équivalente les $3$ qui fournissent $a$).
\item \textbf{Faux.} La formule correcte est $\mathrm{Card}(A \cup B) = \mathrm{Card}(A) + \mathrm{Card}(B) - \mathrm{Card}(A \cap B)$ : les éléments communs sont comptés deux fois dans la somme, il faut les retrancher une fois. L'égalité proposée n'est vraie que si $A \cap B = \varnothing$. Contre-exemple : $A = \{1;2\}$, $B = \{2;3\}$ ; alors $\mathrm{Card}(A \cup B) = 3 \neq 2 + 2 = 4$.
\item \textbf{Vrai.} En appliquant la formule du binôme avec $a = b = 1$ : $2^n = (1+1)^n = \sum_{p=0}^{n} C_n^p \times 1^{n-p} \times 1^p = C_n^0 + C_n^1 + \cdots + C_n^n$. C'est aussi le nombre total de parties d'un ensemble à $n$ éléments.
\end{enumerate}
\end{corrige}

\begin{corrige}[Exercice 3 — Définitions et formules]
\begin{enumerate}
\item Le \cle{cardinal} d'un ensemble fini $E$, noté $\mathrm{Card}(E)$, est le nombre d'éléments de $E$.
\item Pour deux parties $A$ et $B$ d'un ensemble fini :
\[ \mathrm{Card}(A \cup B) = \mathrm{Card}(A) + \mathrm{Card}(B) - \mathrm{Card}(A \cap B). \]
\item Pour $n \ge 1$, $n! = n \times (n-1) \times \cdots \times 2 \times 1$ (produit des entiers de $1$ à $n$). Par convention, $0! = 1$.
\item Pour $0 \le p \le n$ :
\[ A_n^p = \frac{n!}{(n-p)!} \qquad \text{et} \qquad C_n^p = \frac{n!}{p!\,(n-p)!} = \frac{A_n^p}{p!}. \]
\item \textbf{Formule du binôme de Newton} : pour tous nombres $a$ et $b$ et tout entier $n \ge 1$ :
\[ (a+b)^n = \sum_{p=0}^{n} C_n^p \, a^{n-p} b^p = C_n^0 a^n + C_n^1 a^{n-1} b + \cdots + C_n^{n-1} a b^{n-1} + C_n^n b^n. \]
\end{enumerate}
\end{corrige}

\begin{corrige}[Exercice 4 — Calculs directs]
\begin{enumerate}
\item $6! = 6 \times 5 \times 4 \times 3 \times 2 \times 1 = 720$.
\[ \frac{9!}{7!} = \frac{9 \times 8 \times 7!}{7!} = 9 \times 8 = 72. \]
\[ \frac{12!}{10! \times 2!} = \frac{12 \times 11 \times 10!}{10! \times 2} = \frac{132}{2} = 66. \]
\item $A_8^2 = 8 \times 7 = 56$ ; \quad $A_7^3 = 7 \times 6 \times 5 = 210$ ; \quad $A_{10}^3 = 10 \times 9 \times 8 = 720$.
\item $C_6^2 = \dfrac{6 \times 5}{2!} = \dfrac{30}{2} = 15$ ; \quad $C_9^4 = \dfrac{9 \times 8 \times 7 \times 6}{4!} = \dfrac{3024}{24} = 126$ ;
\[ C_{10}^3 = \frac{10 \times 9 \times 8}{3!} = \frac{720}{6} = 120. \]
\item $C_{10}^3 \times 3! = 120 \times 6 = 720 = A_{10}^3$. La relation est vérifiée.
\item Pour $n \ge 1$ : $(n+1)! = (n+1) \times n \times (n-1)!$, donc
\[ \frac{(n+1)!}{(n-1)!} = \frac{(n+1) \times n \times (n-1)!}{(n-1)!} = n(n+1) = n^2 + n. \]
\end{enumerate}
\end{corrige}

\begin{corrige}[Exercice 5 — Plaques d'immatriculation]
\begin{enumerate}
\item Les répétitions sont autorisées : $26$ choix pour chaque lettre, $10$ choix pour chaque chiffre. Par le principe multiplicatif :
\[ 26 \times 26 \times 10 \times 10 \times 10 = 26^2 \times 10^3 = 676 \times 1000 = 676\,000 \text{ plaques}. \]
\item Caractères deux à deux distincts : $26$ choix pour la première lettre, $25$ pour la deuxième ; puis $10$, $9$, $8$ choix pour les trois chiffres :
\[ 26 \times 25 \times 10 \times 9 \times 8 = 650 \times 720 = 468\,000 \text{ plaques}. \]
\item La première lettre est imposée (C : $1$ choix), la deuxième lettre est libre ($26$ choix), les deux premiers chiffres sont libres ($10 \times 10$ choix) et le dernier chiffre est pair ($5$ choix : $0, 2, 4, 6, 8$) :
\[ 1 \times 26 \times 10 \times 10 \times 5 = 13\,000 \text{ plaques}. \]
\end{enumerate}
\end{corrige}

\begin{corrige}[Exercice 6 — Le mot CAMEROUN]
Le mot CAMEROUN possède $8$ lettres distinctes : $4$ voyelles (A, E, O, U) et $4$ consonnes (C, M, R, N).
\begin{enumerate}
\item Un mot de $5$ lettres distinctes est un arrangement de $5$ lettres parmi $8$ :
\[ A_8^5 = 8 \times 7 \times 6 \times 5 \times 4 = 6720 \text{ mots}. \]
\item La première lettre est une voyelle : $4$ choix. Les $4$ lettres suivantes sont choisies et ordonnées parmi les $7$ lettres restantes : $A_7^4 = 7 \times 6 \times 5 \times 4 = 840$ choix. Total :
\[ 4 \times A_7^4 = 4 \times 840 = 3360 \text{ mots}. \]
\item Utiliser les $8$ lettres, c'est les permuter : $8! = 40\,320$ mots. Si le mot commence par C et se termine par N, ces deux places sont fixées et il reste à permuter les $6$ lettres restantes : $6! = 720$ mots.
\end{enumerate}
\end{corrige}

\begin{corrige}[Exercice 7 — Formation d'un comité]
Il y a $7 + 5 = 12$ personnes. Un comité est un choix \emph{sans ordre} : on utilise les combinaisons.
\begin{enumerate}
\item Nombre total de comités : $C_{12}^4 = \dfrac{12 \times 11 \times 10 \times 9}{4!} = \dfrac{11\,880}{24} = 495$.
\item Exactement $2$ femmes : on choisit $2$ femmes parmi $5$ \emph{et} $2$ hommes parmi $7$ :
\[ C_5^2 \times C_7^2 = 10 \times 21 = 210 \text{ comités}. \]
\item « Au moins une femme » est le contraire de « aucune femme » (comité de $4$ hommes) :
\[ C_{12}^4 - C_7^4 = 495 - 35 = 460 \text{ comités}. \]
\item « Au plus un homme » signifie $0$ homme ou $1$ homme :
\begin{itemize}
\item $0$ homme, $4$ femmes : $C_5^4 = 5$ comités ;
\item $1$ homme, $3$ femmes : $C_7^1 \times C_5^3 = 7 \times 10 = 70$ comités.
\end{itemize}
Total : $5 + 70 = 75$ comités.
\end{enumerate}
\end{corrige}

\begin{corrige}[Exercice 8 — Triangle de Pascal et binôme]
\begin{enumerate}
\item Triangle de Pascal jusqu'à $n = 6$ (chaque coefficient est la somme des deux situés au-dessus) :
\begin{center}
\begin{tabular}{|c|c|c|c|c|c|c|c|}
\hline
\rowcolor{popblueL} $n$ & $C_n^0$ & $C_n^1$ & $C_n^2$ & $C_n^3$ & $C_n^4$ & $C_n^5$ & $C_n^6$ \\
\hline
$0$ & $1$ & & & & & & \\
$1$ & $1$ & $1$ & & & & & \\
$2$ & $1$ & $2$ & $1$ & & & & \\
$3$ & $1$ & $3$ & $3$ & $1$ & & & \\
$4$ & $1$ & $4$ & $6$ & $4$ & $1$ & & \\
$5$ & $1$ & $5$ & $10$ & $10$ & $5$ & $1$ & \\
$6$ & $1$ & $6$ & $15$ & $20$ & $15$ & $6$ & $1$ \\
\hline
\end{tabular}
\end{center}
\item D'après la ligne $n = 5$ :
\[ (a+b)^5 = a^5 + 5a^4 b + 10 a^3 b^2 + 10 a^2 b^3 + 5 a b^4 + b^5. \]
D'après la ligne $n = 4$, en remplaçant $b$ par $-b$ (signes alternés) :
\[ (a-b)^4 = a^4 - 4 a^3 b + 6 a^2 b^2 - 4 a b^3 + b^4. \]
\item On applique la formule du binôme avec $a = 2x$ et $b = -3$, ligne $n = 4$ :
\begin{align*}
(2x-3)^4 &= (2x)^4 + 4(2x)^3(-3) + 6(2x)^2(-3)^2 + 4(2x)(-3)^3 + (-3)^4 \\
&= 16x^4 - 96x^3 + 216x^2 - 216x + 81.
\end{align*}
Le coefficient de $x^2$ est $C_4^2 \times 2^2 \times (-3)^2 = 6 \times 4 \times 9 = 216$.
\item Appliquons la formule du binôme avec $a = 1$ et $b = 1$ :
\[ (1+1)^n = \sum_{p=0}^{n} C_n^p \, 1^{n-p} \, 1^p \quad \Longrightarrow \quad C_n^0 + C_n^1 + \cdots + C_n^n = 2^n. \]
Or $C_n^p$ est le nombre de parties à $p$ éléments ; la somme compte donc toutes les parties d'un ensemble à $n$ éléments, qui est $2^n$. Pour $n = 6$ : $2^6 = 64$ parties (vérification : $1+6+15+20+15+6+1 = 64$).
\end{enumerate}
\end{corrige}

\begin{corrige}[Exercice 9 — Sports au lycée (type Probatoire)]
Notons $F$ l'ensemble des footballeurs ($280$ élèves) et $H$ celui des handballeurs ($150$ élèves), avec $\mathrm{Card}(F \cap H) = 60$.
\begin{enumerate}
\item Diagramme de Venn : la zone commune contient $60$ élèves ; la zone « football seulement » contient $280 - 60 = 220$ ; la zone « handball seulement » contient $150 - 60 = 90$.
\begin{popfigure}
\begin{tikzpicture}[scale=0.9]
\draw[thick, popline] (0,0) rectangle (8.6,4.6);
\node[popmuted] at (8.1,4.2) {$E$ : $500$};
\draw[thick, popblue, fill=popblueL] (3.1,2.3) circle (1.7);
\draw[thick, poporange, fill=poporangeL, fill opacity=0.6] (5.3,2.3) circle (1.7);
\node[popblue] at (1.9,3.7) {$F$};
\node[poporange] at (6.6,3.7) {$H$};
\node at (2.4,2.3) {$220$};
\node at (4.2,2.3) {$60$};
\node at (6.0,2.3) {$90$};
\node[popdark] at (7.7,0.6) {$130$};
\end{tikzpicture}
\legende{Diagramme de Venn de la répartition des élèves.}
\end{popfigure}
\item 
\begin{enumerate}
\item Football seulement : $280 - 60 = 220$ élèves.
\item Handball seulement : $150 - 60 = 90$ élèves.
\item Au moins un des deux sports :
\[ \mathrm{Card}(F \cup H) = 280 + 150 - 60 = 370 \text{ élèves}. \]
\end{enumerate}
\item Aucun des deux sports : $500 - 370 = 130$ élèves.
\item 
\begin{enumerate}
\item Groupes de $3$ élèves parmi $500$ (sans ordre) :
\[ C_{500}^3 = \frac{500 \times 499 \times 498}{6} = \frac{124\,251\,000}{6} = 20\,708\,500. \]
\item Groupes formés uniquement de footballeurs ($280$ élèves pratiquent le football) :
\[ C_{280}^3 = \frac{280 \times 279 \times 278}{6} = \frac{21\,717\,360}{6} = 3\,619\,560. \]
\item « Au moins un élève ne pratiquant aucun sport » est le contraire de « les $3$ élèves pratiquent tous au moins un sport » ($370$ élèves) :
\[ C_{500}^3 - C_{370}^3 = 20\,708\,500 - \frac{370 \times 369 \times 368}{6} = 20\,708\,500 - 8\,373\,840 = 12\,334\,660. \]
\end{enumerate}
\end{enumerate}
\textbf{Barème indicatif (sur 5 pts)} : diagramme correct avec les trois zones et l'extérieur : $1$ pt ; question 2 : $1$ pt ; question 3 : $0,5$ pt ; question 4.a : $0,5$ pt ; 4.b : $1$ pt ; 4.c (passage au contraire) : $1$ pt.\par
\textbf{Points de vigilance} : citer explicitement la formule $\mathrm{Card}(F \cup H) = \mathrm{Card}(F) + \mathrm{Card}(H) - \mathrm{Card}(F \cap H)$ ; ne pas confondre « football » ($280$) et « football seulement » ($220$) ; justifier l'emploi des combinaisons (choix sans ordre) et annoncer le passage à l'événement contraire en 4.c.
\end{corrige}

\begin{corrige}[Exercice 10 — Finale du 100 m (type Probatoire)]
\begin{enumerate}
\item Un classement complet est une permutation des $8$ coureurs : $8! = 40\,320$ classements.
\item 
\begin{enumerate}
\item Un podium est une liste ordonnée sans répétition de $3$ coureurs parmi $8$ : $A_8^3 = 8 \times 7 \times 6 = 336$ podiums.
\item Essomba est premier ($1$ choix pour la première place) ; il reste à ordonner $2$ coureurs parmi les $7$ restants : $A_7^2 = 7 \times 6 = 42$ podiums.
\item « Au moins l'un des deux Camerounais » est le contraire de « aucun des deux ». Les podiums sans Essomba ni Mbarga sont des arrangements de $3$ coureurs parmi les $6$ autres : $A_6^3 = 6 \times 5 \times 4 = 120$. Donc :
\[ A_8^3 - A_6^3 = 336 - 120 = 216 \text{ podiums}. \]
\end{enumerate}
\item 
\begin{enumerate}
\item Choix sans ordre de $3$ coureurs parmi $8$ : $C_8^3 = \dfrac{8 \times 7 \times 6}{6} = 56$ choix.
\item Exactement l'un des deux Camerounais : on choisit lequel des deux ($2$ choix), puis $2$ coureurs parmi les $6$ autres :
\[ 2 \times C_6^2 = 2 \times 15 = 30 \text{ choix}. \]
\end{enumerate}
\item Vérification : $C_8^3 \times 3! = 56 \times 6 = 336 = A_8^3$. C'est cohérent : tout podium s'obtient en choisissant $3$ coureurs ($C_8^3$ façons) puis en les ordonnant sur les trois marches ($3!$ façons).
\end{enumerate}
\textbf{Barème indicatif (sur 5 pts)} : question 1 : $0,5$ pt ; 2.a : $1$ pt ; 2.b : $0,5$ pt ; 2.c : $1$ pt ; 3.a : $0,5$ pt ; 3.b : $1$ pt ; question 4 : $0,5$ pt.\par
\textbf{Points de vigilance} : distinguer clairement la situation ordonnée (podium : arrangements) de la situation non ordonnée (contrôle : combinaisons) et le justifier dans la rédaction ; en 2.c, expliciter le raisonnement par le contraire ; en 3.b, ne pas oublier le facteur $2$ (choix du Camerounais retenu).
\end{corrige}

\begin{corrige}[Exercice 11 — Restaurant et menus (problème de synthèse)]
\begin{enumerate}
\item Un menu complet est un triplet (entrée, plat, dessert) : par le principe multiplicatif,
\[ 4 \times 5 \times 3 = 60 \text{ menus}. \]
\item 
\begin{enumerate}
\item Chaque ami choisit librement parmi $60$ menus, les choix sont indépendants et peuvent être identiques :
\[ 60 \times 60 = 3600 \text{ couples de menus}. \]
\item Menus entièrement différents : pour chaque composante, les choix d'Alima et de Bassile doivent différer. Alima choisit librement ($4 \times 5 \times 3$), puis Bassile a $3$ entrées, $4$ plats et $2$ desserts possibles :
\[ (4 \times 5 \times 3) \times (3 \times 4 \times 2) = 60 \times 24 = 1440 \text{ couples}. \]
\end{enumerate}
\item Une ardoise est un choix \emph{sans ordre} de $2$ entrées parmi $4$, $3$ plats parmi $5$ et $2$ desserts parmi $3$ :
\[ C_4^2 \times C_5^3 \times C_3^2 = 6 \times 10 \times 3 = 180 \text{ ardoises}. \]
\item L'affirmation est \textbf{fausse}. Une entrée et un plat ne sont pas deux éléments d'un même ensemble de $9$ plats : ils proviennent de deux ensembles distincts ($4$ entrées, $5$ plats). Le nombre de couples ordonnés (entrée, plat) est donné par le principe multiplicatif :
\[ 4 \times 5 = 20 \quad (\text{et non } A_9^2 = 72). \]
\item Le nombre de menus est $4 \times n \times 3 = 12n$. L'équation $12n = 120$ donne $n = 10$ plats. Alors :
\[ C_{10}^2 = \frac{10 \times 9}{2} = 45. \]
\textbf{Interprétation} : $45$ est le nombre de façons de choisir $2$ plats différents parmi les $10$ plats du restaurant (par exemple pour composer une formule « duo de plats »).
\end{enumerate}
\textbf{Barème indicatif (sur 5 pts)} : question 1 : $0,5$ pt ; 2.a : $0,5$ pt ; 2.b : $1$ pt ; question 3 : $1$ pt ; question 4 (justification incluse) : $1$ pt ; question 5 (équation, résolution, interprétation) : $1$ pt.\par
\textbf{Points de vigilance} : le principe multiplicatif s'applique à des choix successifs dans des ensembles \emph{distincts} ; ne pas utiliser $A_n^p$ lorsque les éléments proviennent de catégories différentes ; en 2.b, bien comprendre que « entièrement différents » impose une différence sur \emph{chacune} des trois composantes ; toute question de modélisation gagne à être conclue par une phrase d'interprétation dans le contexte.
\end{corrige}

\newpage

\coursec{Fiche bilan}

\begin{popfigure}
\begin{tikzpicture}[mindmap, grow cyclic, every node/.style={concept, font=\small},
concept color=popblue,
level 1/.append style={level distance=3.6cm, sibling angle=60},
level 2/.append style={level distance=2.2cm, sibling angle=45, font=\scriptsize}]
\node[concept color=popblue, text=white, font=\bfseries]{Dénombrement}
child[concept color=poporange]{ node {Cardinal et outils}
  child { node {Diagrammes de Venn} }
  child { node {Tableaux à double entrée} }
  child { node {$\mathrm{Card}(A\cup B)$} }
}
child[concept color=popgreen]{ node {p-uplets}
  child { node {Avec répétition : $n^p$} }
  child { node {Sans répétition} }
  child { node {Arbre de choix} }
}
child[concept color=poppurple]{ node {Factorielle et permutations}
  child { node {$n!$} }
  child { node {$n!$ permutations} }
}
child[concept color=poppink]{ node {Arrangements}
  child { node {$A_n^p=\frac{n!}{(n-p)!}$} }
  child { node {Ordre important} }
}
child[concept color=popgold]{ node {Combinaisons}
  child { node {$C_n^p=\frac{n!}{p!(n-p)!}$} }
  child { node {Ordre sans importance} }
}
child[concept color=popdark]{ node {Binôme de Newton}
  child { node {Triangle de Pascal} }
  child { node {$(a+b)^n$} }
};
\end{tikzpicture}
\legende{Carte mentale de la leçon : les six grandes familles d'outils du dénombrement.}
\end{popfigure}

\begin{aretenir}[L'essentiel de la leçon]
\textbf{Définitions clés.}
\begin{itemize}
\item Le \cle{cardinal} d'un ensemble fini $E$, noté $\mathrm{Card}(E)$, est le nombre d'éléments de $E$.
\item Un \cle{p-uplet} (ou p-liste) est une liste ordonnée de $p$ éléments, avec ou sans répétition.
\item Une \cle{permutation} d'un ensemble à $n$ éléments est un rangement ordonné de ses $n$ éléments.
\item Un \cle{arrangement} de $p$ éléments parmi $n$ est une liste ordonnée de $p$ éléments distincts (avec $0\le p\le n$).
\item Une \cle{combinaison} de $p$ éléments parmi $n$ est une partie (sous-ensemble) à $p$ éléments : l'ordre ne compte pas.
\end{itemize}

\textbf{Formules à connaître par cœur.}
\begin{center}
\begin{tabularx}{\linewidth}{|l|X|}
\hline
\rowcolor{popblueL} \textbf{Situation} & \textbf{Formule} \\
\hline
Réunion de deux ensembles & $\mathrm{Card}(A\cup B)=\mathrm{Card}(A)+\mathrm{Card}(B)-\mathrm{Card}(A\cap B)$ \\
\hline
Produit cartésien & $\mathrm{Card}(E\times F)=\mathrm{Card}(E)\times\mathrm{Card}(F)$ \\
\hline
p-uplets avec répétition & $n^p$ \\
\hline
Factorielle & $n!=n\times(n-1)\times\cdots\times 2\times 1$, avec $0!=1$ \\
\hline
Permutations de $n$ éléments & $n!$ \\
\hline
Arrangements & $A_n^p=\dfrac{n!}{(n-p)!}=n(n-1)\cdots(n-p+1)$ \\
\hline
Combinaisons & $C_n^p=\dfrac{n!}{p!\,(n-p)!}=\dfrac{A_n^p}{p!}$ \\
\hline
Nombre de parties de $E$ ($\mathrm{Card}(E)=n$) & $2^n$ \\
\hline
Relation de Pascal & $C_n^p=C_{n-1}^{p-1}+C_{n-1}^{p}$ \quad ($1\le p\le n-1$) \\
\hline
Formule du binôme & $(a+b)^n=\displaystyle\sum_{p=0}^{n} C_n^p\, a^{n-p}b^{p}$ \\
\hline
\end{tabularx}
\end{center}

\textbf{Méthodes express.}
\begin{itemize}
\item \textbf{Ordre + répétition possible} (code PIN, plaque d'immatriculation) : $n^p$.
\item \textbf{Ordre + sans répétition} (podium, tirage successif sans remise) : $A_n^p$.
\item \textbf{Pas d'ordre} (équipe, main de cartes, tirage simultané) : $C_n^p$.
\item \textbf{Tous les éléments rangés} : $n!$.
\item Pour un développement de $(a+b)^n$ : construire le triangle de Pascal jusqu'à la ligne $n$, puis faire décroître les puissances de $a$ et croître celles de $b$.
\end{itemize}
\end{aretenir}

\begin{attention}[Vérifications et pièges classiques]
\begin{itemize}
\item Un résultat de dénombrement est toujours un \cle{entier positif ou nul} : si tu obtiens une fraction ou un nombre négatif, il y a une erreur.
\item Ne confonds pas $A_n^p$ et $C_n^p$ : demande-toi toujours si l'\cle{ordre} des éléments choisis modifie le résultat. Par exemple, $A_5^2=20$ alors que $C_5^2=10$.
\item Dans la formule d'inclusion-exclusion, n'oublie pas de \emph{retrancher} $\mathrm{Card}(A\cap B)$, sinon les éléments communs sont comptés deux fois.
\item La relation de Pascal n'est valable que pour $1\le p\le n-1$ ; aux bords on a $C_n^0=C_n^n=1$.
\item $0!=1$ est une convention indispensable : elle garantit que $C_n^n=\dfrac{n!}{n!\,0!}=1$.
\end{itemize}
\end{attention}

\begin{exemplebox}[Deux applications rapides pour vérifier ta maîtrise]
\textbf{1.} Dans une classe de Première C à Yaoundé, $30$ élèves étudient l'allemand ou l'espagnol : $18$ étudient l'allemand, $15$ l'espagnol. Le nombre d'élèves étudiant les deux langues est :
\[ \mathrm{Card}(A\cap B)=18+15-30=3. \]
\textbf{2.} Le bureau d'un club compte un président, un secrétaire et un trésorier, choisis parmi $10$ membres, sans cumul de postes. L'ordre (le poste) compte et il n'y a pas répétition : le nombre de bureaux possibles est
\[ A_{10}^{3}=10\times 9\times 8=720. \]
Si l'on choisit simplement une délégation de $3$ membres sans distinction de rôle, le nombre de choix est
\[ C_{10}^{3}=\dfrac{10\times 9\times 8}{3!}=\dfrac{720}{6}=120. \]
\end{exemplebox}

\begin{aretenir}[Checklist avant le Probatoire]
\begin{itemize}
\item[$\square$] Je sais calculer $\mathrm{Card}(A\cup B)$ avec la formule d'inclusion-exclusion.
\item[$\square$] Je sais construire et exploiter un diagramme de Venn et un tableau à double entrée.
\item[$\square$] Je sais dessiner un arbre de choix pour dénombrer des issues.
\item[$\square$] Je sais distinguer une situation « avec répétition » d'une situation « sans répétition ».
\item[$\square$] Je sais calculer $n!$ et je me souviens que $0!=1$.
\item[$\square$] Je sais calculer $A_n^p$ et $C_n^p$ avec une calculatrice ou à la main.
\item[$\square$] Je sais choisir entre $n^p$, $A_n^p$, $C_n^p$ et $n!$ selon l'ordre et la répétition.
\item[$\square$] Je sais que le nombre de parties d'un ensemble à $n$ éléments est $2^n$.
\item[$\square$] Je sais construire le triangle de Pascal grâce à $C_n^p=C_{n-1}^{p-1}+C_{n-1}^p$.
\item[$\square$] Je sais développer $(a+b)^n$ avec les coefficients du binôme.
\item[$\square$] Je vérifie toujours la cohérence du résultat (un nombre de cas est un entier positif).
\end{itemize}
\end{aretenir}

\findecours

\end{document}
